%        1         2         3         4         5         6         7         8         9         0         1         2
%23456789012345678901234567890123456789012345678901234567890123456789012345678901234567890123456789012345678901234567890
\documentclass[12pt]{MIA-USA}

\usepackage{layout}
\usepackage{subcaption}
\usepackage{paralist}
\usepackage{breakcites}

% Paquetes y macros que requiere el contenido convertido desde Markdown (Pandoc)
% ============================================================================
% pandoc-preamble.tex
% Paquetes y macros que requiere el LaTeX generado por Pandoc a partir de los
% archivos Markdown del informe (tablas booktabs/longtable, bloques de código
% con resaltado de sintaxis, listas compactas, etc.).
%
% La clase MIA-USA.cls ya carga: graphicx, xcolor[table], longtable, array,
% colortbl, tabularx, amsmath, tikz, fancyhdr, multirow, multicol, tcolorbox.
% Aquí solo se añade lo que falta. Se usa pdfLaTeX (no se importan fuentes
% Times New Roman/fontspec del metadata.yaml, que eran para XeLaTeX).
% ============================================================================

% --- Tablas generadas por Pandoc -------------------------------------------
\usepackage{booktabs}        % \toprule \midrule \bottomrule
\usepackage{calc}            % cálculos de ancho de columnas de Pandoc
\providecommand{\real}[1]{#1}

% Pandoc define columnas de tablas con tipos personalizados; estos comandos
% reproducen su plantilla por defecto para que las tablas compilen.
\newlength{\cslhangindent}
\setlength{\cslhangindent}{1.5em}

% Pandoc envuelve las tablas sin título en {\def\LTcaptype{none} ...}. Como la
% plantilla carga el paquete caption (vía subcaption), este intenta usar un
% contador 'none' que no existe. Lo definimos para evitar el error fatal.
\makeatletter
\@ifundefined{c@none}{\newcounter{none}}{}
\makeatother

% --- Tablas anchas: que quepan en la página --------------------------------
\usepackage{etoolbox}
\AtBeginEnvironment{longtable}{\footnotesize\setlength{\tabcolsep}{3pt}\renewcommand{\arraystretch}{0.95}}

% Tablas en orientación horizontal cuando el ancho lo requiere
\usepackage{pdflscape}

% --- Bloques de código con resaltado (salida de Pandoc) --------------------
\usepackage{fvextra}
% Ajuste automático de línea para que el código no se desborde del margen
\DefineVerbatimEnvironment{Highlighting}{Verbatim}{breaklines,breakanywhere,commandchars=\\\{\}}
\newcommand{\VerbBar}{|}
\newcommand{\VERB}{\Verb[commandchars=\\\{\}]}
\newenvironment{Shaded}{}{}
\newcommand{\AlertTok}[1]{\textcolor[rgb]{1.00,0.00,0.00}{\textbf{#1}}}
\newcommand{\AnnotationTok}[1]{\textcolor[rgb]{0.38,0.63,0.69}{\textbf{\textit{#1}}}}
\newcommand{\AttributeTok}[1]{\textcolor[rgb]{0.49,0.56,0.16}{#1}}
\newcommand{\BaseNTok}[1]{\textcolor[rgb]{0.25,0.63,0.44}{#1}}
\newcommand{\BuiltInTok}[1]{\textcolor[rgb]{0.00,0.50,0.00}{#1}}
\newcommand{\CharTok}[1]{\textcolor[rgb]{0.25,0.44,0.63}{#1}}
\newcommand{\CommentTok}[1]{\textcolor[rgb]{0.38,0.63,0.69}{\textit{#1}}}
\newcommand{\CommentVarTok}[1]{\textcolor[rgb]{0.38,0.63,0.69}{\textbf{\textit{#1}}}}
\newcommand{\ConstantTok}[1]{\textcolor[rgb]{0.53,0.00,0.00}{#1}}
\newcommand{\ControlFlowTok}[1]{\textcolor[rgb]{0.00,0.44,0.13}{\textbf{#1}}}
\newcommand{\DataTypeTok}[1]{\textcolor[rgb]{0.56,0.13,0.00}{#1}}
\newcommand{\DecValTok}[1]{\textcolor[rgb]{0.25,0.63,0.44}{#1}}
\newcommand{\DocumentationTok}[1]{\textcolor[rgb]{0.73,0.13,0.13}{\textit{#1}}}
\newcommand{\ErrorTok}[1]{\textcolor[rgb]{1.00,0.00,0.00}{\textbf{#1}}}
\newcommand{\ExtensionTok}[1]{#1}
\newcommand{\FloatTok}[1]{\textcolor[rgb]{0.25,0.63,0.44}{#1}}
\newcommand{\FunctionTok}[1]{\textcolor[rgb]{0.02,0.16,0.49}{#1}}
\newcommand{\ImportTok}[1]{\textcolor[rgb]{0.00,0.50,0.00}{\textbf{#1}}}
\newcommand{\InformationTok}[1]{\textcolor[rgb]{0.38,0.63,0.69}{\textbf{\textit{#1}}}}
\newcommand{\KeywordTok}[1]{\textcolor[rgb]{0.00,0.44,0.13}{\textbf{#1}}}
\newcommand{\NormalTok}[1]{#1}
\newcommand{\OperatorTok}[1]{\textcolor[rgb]{0.40,0.40,0.40}{#1}}
\newcommand{\OtherTok}[1]{\textcolor[rgb]{0.00,0.44,0.13}{#1}}
\newcommand{\PreprocessorTok}[1]{\textcolor[rgb]{0.74,0.48,0.00}{#1}}
\newcommand{\RegionMarkerTok}[1]{#1}
\newcommand{\SpecialCharTok}[1]{\textcolor[rgb]{0.25,0.44,0.63}{#1}}
\newcommand{\SpecialStringTok}[1]{\textcolor[rgb]{0.73,0.40,0.53}{#1}}
\newcommand{\StringTok}[1]{\textcolor[rgb]{0.25,0.44,0.63}{#1}}
\newcommand{\VariableTok}[1]{\textcolor[rgb]{0.10,0.09,0.49}{#1}}
\newcommand{\VerbatimStringTok}[1]{\textcolor[rgb]{0.25,0.44,0.63}{#1}}
\newcommand{\WarningTok}[1]{\textcolor[rgb]{0.38,0.63,0.69}{\textbf{\textit{#1}}}}

% --- Listas compactas y varios de Pandoc -----------------------------------
\providecommand{\tightlist}{%
  \setlength{\itemsep}{0pt}\setlength{\parskip}{0pt}}
\setlength{\emergencystretch}{3em} % evita líneas desbordadas

% Caracteres de dibujo de cajas (─ │ └ ├ ┌ ┐ …) usados en los árboles de
% directorios dentro de los bloques de código. pmboxdraw los define para
% pdfLaTeX+inputenc.
\usepackage{pmboxdraw}

% Otros caracteres Unicode que la fuente base / inputenc no traen.
\usepackage{newunicodechar}
\newunicodechar{⋮}{\ensuremath{\vdots}}        % U+22EE elipsis vertical
\newunicodechar{→}{\ensuremath{\rightarrow}}   % U+2192 flecha derecha
\newunicodechar{←}{\ensuremath{\leftarrow}}    % U+2190
\newunicodechar{↔}{\ensuremath{\leftrightarrow}}
\newunicodechar{−}{\ensuremath{-}}             % U+2212 signo menos
\newunicodechar{≤}{\ensuremath{\leq}}          % U+2264
\newunicodechar{≥}{\ensuremath{\geq}}          % U+2265
\newunicodechar{×}{\ensuremath{\times}}        % U+00D7 (por si acaso)
\newunicodechar{►}{\ensuremath{\blacktriangleright}} % U+25BA
\newunicodechar{✅}{\checkmark}                 % U+2705
\newunicodechar{❌}{\ensuremath{\times}}        % U+274C
\newunicodechar{⁴}{\textsuperscript{4}}        % U+2074 superíndices (notas)
\newunicodechar{⁵}{\textsuperscript{5}}
\newunicodechar{⁶}{\textsuperscript{6}}
\newunicodechar{⁷}{\textsuperscript{7}}
\newunicodechar{⁸}{\textsuperscript{8}}
\newunicodechar{⁹}{\textsuperscript{9}}
\usepackage{amssymb}                            % \checkmark, \blacktriangleright

% Pandoc puede emitir \passthrough en código en línea
\providecommand{\passthrough}[1]{#1}

% Enlaces de imágenes/figuras de Pandoc dentro de \includegraphics
\makeatletter
\providecommand{\pandocbounded}[1]{#1}
\makeatother


\usepackage[pagebackref=true,breaklinks=true]{hyperref}
\hypersetup{
	pdftitle={Documento de trabajo Final}, 
	pdfauthor={MIA-D.Martinez}, 
	pdfsubject={}, 
	pdfcreator={DarMart with Latex}, 
	colorlinks=true,
	linkcolor=GrisUno,
	citecolor=AzulInstitucional,
	filecolor=AzulClaro,
	urlcolor=AzulInstitucional,
	linktoc=all
}

\graphicspath{{Images/}}


\title{Orbit Engine: Plataforma SaaS para la Gesti\'on Integral de Procesos Internos en Peque\~nas y Medianas Empresas}
\author{Rinc\'on Suarez Fabi\'an \\
Rodr\'iguez Forero Nicol\'as \\
Velasco Gonz\'alez Daniel}
\documenttype{Proyecto de Grado}

\advisor{Juan Pablo Ospina}

\coadvisor{Camilo Rodr\'iguez}

% Si hay un coasesor 
% \coadvisor{Nombres y apellidos Co-Asesor}


% Aqui comienza el documento
% - -- --- ----- ------- ------- ----- --- -- -
\begin{document}
    % \layout{}
	
	\maketitle
	% Incluye la pagina de aceptación (Opcional)
	\input{FrontMatter/F01-Approval}

    % Dedicatoria (Opcional)
    % Dedicatorias de los autores del trabajo de grado.
\subsection*{Dedicatoria}
\vspace*{0.8cm}
{\hfill
    \begin{minipage}{0.7\textwidth}
        \noindent
        \vspace*{4.5cm}

        \hfill\textit{Fabi\'an Rinc\'on Su\'arez}\\
    \end{minipage}
}
\vfill

{\hfill
    \begin{minipage}{0.7\textwidth}
        \noindent
        \vspace*{4.5cm}

        \hfill\textit{Nicol\'as Rodr\'iguez Forero}\\
    \end{minipage}
}
\vfill

{\hfill
    \begin{minipage}{0.7\textwidth}
        \noindent
        A mis padres, quienes desde el primer d\'ia de mi formaci\'on han sido el cimiento sobre el cual se ha construido cada uno de mis logros. Su amor incondicional, su esfuerzo silencioso y su confianza inquebrantable me ense\~naron que la disciplina y la constancia son los verdaderos motores del crecimiento personal y profesional. A ellos les debo no solo la posibilidad de haber cursado esta carrera, sino tambi\'en los valores que me han acompa\~nado en cada decisi\'on, en cada desvelo y en cada peque\~no triunfo a lo largo de este camino.\\
        \vspace*{0.2cm}

        \noindent
        A mis profesores, que con paciencia, vocaci\'on y generosidad compartieron conmigo no solamente conocimientos t\'ecnicos, sino tambi\'en una manera de pensar, de cuestionar y de aproximarse a los problemas. Cada clase, cada correcci\'on y cada conversaci\'on dej\'o una huella que hoy se refleja en este trabajo. Esta dedicatoria es, en \'ultima instancia, un peque\~no reconocimiento a quienes creyeron en m\'i incluso cuando yo mismo dudaba.\\
        \vspace*{0.3cm}

        \hfill\textit{Daniel Velasco Gonz\'alez}\\
    \end{minipage}
}
\newpage


    % Agradecimientos (Opcional)
    % Página de agradecimientos del trabajo de grado.

\section*{Agradecimientos}

En primer lugar, expresamos nuestra m\'as profunda gratitud a \textbf{Dios}, por haber sido la luz que ilumin\'o cada etapa de este proceso, por brindarnos la salud, la sabidur\'ia y la perseverancia necesarias para afrontar los retos que conlleva la culminaci\'on de un proyecto de esta magnitud. A \'El debemos las oportunidades que se nos presentaron a lo largo de la carrera y la fortaleza para no rendirnos en los momentos en los que el camino se torn\'o particularmente exigente.

Queremos expresar un agradecimiento especial a nuestro asesor, el ingeniero \textbf{Juan Pablo Ospina}, cuya gu\'ia metodol\'ogica, rigurosidad acad\'emica y disposici\'on permanente fueron determinantes para dar forma, estructura y profundidad a este trabajo. Sus observaciones precisas, su capacidad para se\~nalar oportunidades de mejora y su acompa\~namiento durante cada fase del proyecto se convirtieron en un referente invaluable para nuestro proceso de formaci\'on como profesionales.

De igual manera, extendemos nuestro reconocimiento al ingeniero \textbf{Camilo Rodr\'iguez}, quien en su rol de coasesor aport\'o una visi\'on cr\'itica y constructiva sobre los aspectos t\'ecnicos del proyecto, ayud\'andonos a tomar mejores decisiones de arquitectura, dise\~no y desarrollo. Su experiencia y su disponibilidad para escuchar nuestras inquietudes nos permitieron consolidar una propuesta s\'olida y coherente con las exigencias del entorno profesional actual.

Agradecemos tambi\'en al ingeniero \textbf{Camilo Salazar}, cuya orientaci\'on a lo largo de distintas asignaturas y espacios acad\'emicos contribuy\'o significativamente a fortalecer las bases conceptuales y pr\'acticas sobre las cuales se sustenta este proyecto. Su pasi\'on por la ense\~nanza y su compromiso con la formaci\'on de profesionales \'integros constituyen un ejemplo que llevaremos con nosotros m\'as all\'a de las aulas.

As\'i mismo, queremos reconocer al ingeniero \textbf{William Frasser}, cuyas ense\~nanzas, recomendaciones y aportes durante el desarrollo de la carrera resultaron fundamentales para comprender la importancia de la calidad, la planeaci\'on y la mejora continua en cualquier proyecto de ingenier\'ia. Su acompa\~namiento nos permiti\'o ampliar nuestra perspectiva y abordar los retos del proyecto con mayor madurez t\'ecnica y profesional.

Un agradecimiento profundo y especial merecen nuestras \textbf{familias}, pilar fundamental e insustituible a lo largo de toda nuestra formaci\'on. A nuestros padres, hermanos y dem\'as seres queridos, quienes nos brindaron su apoyo incondicional, su comprensi\'on durante las largas jornadas de estudio, su paciencia ante las ausencias y su aliento constante en los momentos de mayor exigencia, les debemos buena parte de lo que hoy somos. Sus sacrificios silenciosos, sus palabras de \'animo y la confianza depositada en cada uno de nosotros fueron, sin duda alguna, el sost\'en emocional que hizo posible llegar hasta este punto. Este logro no es solo nuestro: tambi\'en es de ellos, porque cada paso dado en esta carrera lleva impreso el cari\~no, el esfuerzo y la entrega de quienes nos acompa\~naron desde el hogar.

Por \'ultimo, pero no por ello menos importante, extendemos nuestro m\'as sincero agradecimiento a la \textbf{Universidad Sergio Arboleda}, instituci\'on que nos abri\'o sus puertas y nos brind\'o no solo una formaci\'on acad\'emica de alto nivel, sino tambi\'en los espacios, recursos, herramientas y oportunidades que hicieron posible la realizaci\'on de este trabajo de grado. A su cuerpo docente, administrativo y a todas las personas que de una u otra manera hicieron parte de nuestro proceso formativo, les expresamos nuestro reconocimiento y nuestra gratitud por haber sido parte fundamental de esta etapa que hoy concluye y que marca el inicio de nuevos retos profesionales.

%  Se debe dejar la nueva página para guardar el orden de las seciones
\newpage


    % Resumen y Abstract
    % Resumen y Abstract — generado desde docs/informe-final/00-preliminares.md
\chapter*{Resumen}
\addcontentsline{toc}{chapter}{Resumen}

Las pequeñas y medianas empresas (pymes) de América Latina enfrentan una brecha estructural de digitalización en la gestión de sus operaciones internas. La mayoría administra inventario, ventas y clientes mediante hojas de cálculo o registros manuales, sin acceso a herramientas integradas y accesibles que proporcionen visibilidad en tiempo real del estado del negocio. Las soluciones ERP existentes en el mercado son, o bien demasiado costosas y complejas para este segmento, o bien insuficientes en sus capacidades operativas para cubrir el ciclo comercial completo.

El presente trabajo describe el diseño, desarrollo y validación de OrbitEngine, una plataforma SaaS multi-tenant orientada a pymes del sector comercio minorista y mayorista. La plataforma integra módulos de gestión de inventario con alertas automáticas de stock mínimo e historial de movimientos, registro de ventas con generación de facturas y control transaccional de stock, administración de clientes con análisis de comportamiento de compra, y un dashboard de indicadores clave de desempeño (KPIs) en tiempo real, complementado por un módulo de exportación a Excel de los listados de inventario, ventas y clientes.

El proyecto se desarrolló en siete meses por un equipo de tres personas bajo una metodología Scrum adaptada al contexto académico. El backend fue construido con FastAPI y SQLModel sobre PostgreSQL; el frontend con React 19 y TypeScript; y el sistema fue desplegado en Railway (backend y base de datos PostgreSQL) y Vercel (frontend) mediante un pipeline de integración y entrega continua (CI/CD) automatizado con GitHub Actions. La arquitectura multi-tenant adoptada garantiza el aislamiento total de datos entre organizaciones mediante un discriminador \texttt{organization\_id} aplicado sistemáticamente en todas las tablas de negocio y en el token JWT de autenticación.

La plataforma fue validada con empresas piloto del sector comercio, midiendo el impacto en la eficiencia operativa a través de la reducción de tiempos en tareas administrativas, la disminución de errores de inventario y encuestas de usabilidad basadas en la escala SUS (\emph{System Usability Scale}). Los resultados obtenidos evidencian mejoras significativas en la gestión operativa de las empresas participantes y confirman la viabilidad técnica y funcional de la solución propuesta.

\textbf{Palabras clave:} SaaS, pymes, gestión operativa, inventario, ventas, multi-tenancy, FastAPI, React, Railway, Vercel, arquitectura en capas.

\clearpage
\chapter*{Abstract}
\addcontentsline{toc}{chapter}{Abstract}

Small and medium-sized enterprises (SMEs) in Latin America face a structural digitalization gap in the management of their internal operations. Most businesses rely on spreadsheets or manual records to manage inventory, sales, and customers, lacking access to integrated, affordable tools that provide real-time visibility into business performance. Existing ERP solutions in the market are either too costly and complex for this segment, or too limited in operational scope to support the full commercial cycle.

This work presents the design, development, and validation of OrbitEngine, a multi-tenant SaaS platform targeting retail and wholesale commerce SMEs. The platform integrates modules for inventory management with automatic low-stock alerts and movement history, sales recording with invoice generation and transactional stock control, customer administration with purchase behavior analytics, and a real-time key performance indicator (KPI) dashboard, complemented by an Excel export module for inventory, sales and customer listings.

The project was developed over seven months by a three-person team following an Scrum methodology adapted to an academic context. The backend was built with FastAPI and SQLModel on PostgreSQL; the frontend with React 19 and TypeScript; and the system was deployed on Railway (backend and PostgreSQL database) and Vercel (frontend) with an automated CI/CD pipeline using GitHub Actions. The multi-tenant architecture ensures complete data isolation between organizations through an \texttt{organization\_id} discriminator applied consistently across all business tables and embedded in the JWT authentication token.

The platform was validated with pilot companies from the commerce sector, measuring operational efficiency improvements through reductions in administrative task completion time, decreases in inventory discrepancy rates, and usability assessments based on the System Usability Scale (SUS). Results demonstrate significant improvements in the operational management of participating companies and confirm the technical and functional viability of the proposed solution.

\textbf{Keywords:} SaaS, SMEs, operational management, inventory, sales, multi-tenancy, FastAPI, React, Railway, Vercel, layered architecture.


    % Crea la tabla de contenidos a partir de la estructura del documento
    \tableofcontents

    % Listas de figuras y tablas (requerido por el formato institucional USA).
    % tocbibind (cargado por la clase) ya las añade automáticamente al índice.
    \cleardoublepage
    \listoffigures

    \cleardoublepage
    \listoftables

    % - -- --- ----- Cuerpo del informe (Capítulos 1-7) ----- --- -- -
    % Convertidos automáticamente desde docs/informe-final/*.md con Pandoc.
    % Cada archivo comienza con su propio \chapter{...}; la numeración la
    % gestiona automáticamente la clase MIA-USA.

    \chapter{Introducción}\label{introducciuxf3n}

\begin{center}\rule{0.5\linewidth}{0.5pt}\end{center}

\section{Planteamiento del Problema}\label{planteamiento-del-problema}

Las pequeñas y medianas empresas (pymes) constituyen el motor económico de América Latina. Según la Comisión Económica para América Latina y el Caribe (CEPAL, 2022), las pymes representan más del 99\% del tejido empresarial en la región y generan entre el 60\% y el 70\% del empleo formal. Sin embargo, pese a su relevancia, enfrentan una brecha estructural frente a las grandes corporaciones en materia de gestión operativa y adopción tecnológica.

Un diagnóstico recurrente en la literatura especializada señala que la gran mayoría de las pymes latinoamericanas aún gestionan sus operaciones ---inventario, ventas, clientes y contabilidad básica--- mediante herramientas rudimentarias: cuadernos físicos, hojas de cálculo de Excel no estructuradas o sistemas ad hoc sin integración entre módulos. De acuerdo con el Banco Interamericano de Desarrollo (BID, 2021), solo el 16\% de las pymes de la región ha adoptado algún tipo de software de gestión empresarial, frente a un 74\% en países de la OCDE.

Esta brecha de digitalización se traduce en consecuencias concretas y medibles para el negocio:

\begin{itemize}
\tightlist
\item
  \textbf{Errores de inventario}: la gestión manual provoca discrepancias frecuentes entre el stock físico y el registrado, generando roturas de stock o sobrestock. Ambos escenarios tienen costos directos: ventas perdidas en el primer caso, capital inmovilizado en el segundo.
\item
  \textbf{Ineficiencia operativa}: los empleados dedican tiempo significativo a tareas administrativas repetitivas ---conciliación de registros, elaboración de reportes manuales, búsqueda de información--- tiempo que podría orientarse a actividades de mayor valor.
\item
  \textbf{Toma de decisiones reactiva}: sin acceso a datos consolidados y en tiempo real, los dueños de pymes toman decisiones basadas en intuición antes que en evidencia. El abastecimiento de productos, la identificación de clientes de alto valor y la detección de tendencias de ventas se realizan de manera subjetiva, sin soporte de indicadores ni reportes estructurados.
\end{itemize}

Las soluciones existentes en el mercado presentan limitaciones específicas para este segmento. Los sistemas ERP de nivel empresarial (SAP, Oracle) son prohibitivos en costo y complejidad para una pyme. Las alternativas de código abierto como Odoo, si bien potentes, requieren infraestructura propia y capacidad técnica para su instalación, configuración y mantenimiento. Las soluciones SaaS disponibles en el mercado latinoamericano (Alegra, Siigo, Defontana) se enfocan principalmente en facturación electrónica y contabilidad, sin integrar capacidades de gestión de inventario, análisis de clientes y reportes operativos en un único entorno accesible.

Se identifica así un problema de investigación y desarrollo bien definido: \textbf{la inexistencia de una plataforma SaaS integrada, accesible y orientada al contexto latinoamericano, que combine la gestión operativa esencial de una pyme ---inventario, ventas, clientes y reportes--- bajo un modelo de precios y complejidad adecuados para este segmento.}

\section{Justificación}\label{justificaciuxf3n}

\subsection{Justificación Económica y Social}\label{justificaciuxf3n-econuxf3mica-y-social}

La digitalización de las pymes no es un objetivo exclusivamente tecnológico; es un imperativo económico. El BID estima que cerrar la brecha digital de las pymes en América Latina podría incrementar el PIB regional en hasta un 5\% adicional para 2030. En términos individuales, los estudios de adopción de ERP en pymes reportan reducciones de costos operativos del 15\% al 30\% y mejoras en la eficiencia de la gestión de inventario de hasta el 40\% en el primer año de implementación (Kumar \& van Hillegersberg, 2000; Duan et al., 2012).

La relevancia del problema es, por tanto, tanto académica como práctica: una solución bien diseñada tiene potencial de impacto directo en la competitividad de cientos de pequeñas empresas.

\subsection{Justificación Tecnológica}\label{justificaciuxf3n-tecnoluxf3gica}

Las decisiones tecnológicas que sustentan OrbitEngine no se justifican por la mera disponibilidad de las herramientas seleccionadas, sino por su capacidad de reducir el riesgo técnico del proyecto y acelerar su tiempo de salida al mercado en un contexto de equipo reducido y plazos académicos acotados:

\begin{enumerate}
\def\labelenumi{\arabic{enumi}.}
\tightlist
\item
  \textbf{Idoneidad del stack para una plataforma transaccional multi-tenant}: la elección de un backend asíncrono basado en FastAPI con tipado estricto, junto con un frontend declarativo en React, responde a requisitos concretos del dominio ---validación rigurosa de datos contables, concurrencia en operaciones de inventario y trazabilidad por organización---, y no a una preferencia genérica por tecnologías populares.
\item
  \textbf{Contrato API como única fuente de verdad}: la adopción de un enfoque \emph{contract-first} mediante OpenAPI permite generar automáticamente el cliente tipado del frontend a partir del esquema del backend, garantizando consistencia entre capas, detectando regresiones en tiempo de compilación y eliminando una clase entera de errores de integración propios de los desarrollos manuales.
\item
  \textbf{Seguridad desde el diseño}: la arquitectura incorpora autenticación basada en JWT, control de acceso por roles a nivel de endpoint y aislamiento lógico de datos por organización como invariantes del modelo, no como capas añadidas a posteriori, lo cual es crítico al manejar información financiera de terceros.
\item
  \textbf{Aprovechamiento de inteligencia artificial generativa en el ciclo de desarrollo}: la incorporación de asistentes de IA en tareas de codificación, documentación y revisión permitió al equipo sostener un ritmo de entrega comparable al de equipos de mayor tamaño, manteniendo estándares de calidad mediante revisión humana sistemática del código generado.
\item
  \textbf{Viabilidad económica del despliegue en la nube para el segmento pyme}: las plataformas de despliegue gestionado seleccionadas ofrecen, en 2026, un costo marginal por organización lo suficientemente bajo como para que el modelo de suscripción resulte rentable incluso con un volumen reducido de clientes, condición que no se cumplía con la infraestructura disponible una década atrás.
\end{enumerate}

\subsection{Justificación Académica}\label{justificaciuxf3n-acaduxe9mica}

La pertinencia académica del presente proyecto se sustenta en tres ejes complementarios que articulan su contribución al conocimiento, su rigor metodológico y su correspondencia con la formación recibida en el programa.

En primer lugar, el proyecto aporta evidencia empírica a un vacío específico de la literatura: si bien la adopción de sistemas de información en pequeñas y medianas empresas ha sido ampliamente documentada en contextos anglosajones y europeos (Kumar \& van Hillegersberg, 2000; Duan et al., 2012), los estudios centrados en pymes latinoamericanas ---y en particular sobre plataformas SaaS multi-tenant diseñadas para sus restricciones de costo, capacidad técnica y conectividad--- son notablemente más escasos. OrbitEngine se propone como un caso de estudio que documenta tanto el diseño arquitectónico como los resultados de su adopción en este contexto regional.

En segundo lugar, la validación del sistema sigue un diseño cuasi-experimental de tipo pre/post con empresas piloto, articulando métricas cuantitativas ---reducción del tiempo en tareas administrativas, disminución de la tasa de error en operaciones de inventario y tiempos de respuesta en producción--- con métricas cualitativas estandarizadas, en particular el instrumento \emph{System Usability Scale} (SUS) y entrevistas semiestructuradas de cierre. Este enfoque mixto permite contrastar con evidencia las hipótesis sobre el impacto de la digitalización en la eficiencia operativa, más allá de la sola demostración funcional del software.

En tercer lugar, el proyecto se inscribe en la línea de trabajo del semillero \emph{Software como Innovación} y articula de manera integral las competencias formativas del pregrado en Ciencias de la Computación e Inteligencia Artificial: ingeniería de requisitos, arquitectura de software, desarrollo full-stack, despliegue en la nube, validación experimental y comunicación de resultados. Lo anterior justifica su carácter como proyecto de grado y no como un desarrollo de software meramente aplicado.

\section{Objetivos}\label{objetivos}

\subsection{Objetivo General}\label{objetivo-general}

Desarrollar e implementar una plataforma SaaS multi-tenant para la gestión integral de procesos internos en pequeñas y medianas empresas, que integre módulos de gestión de inventario, ventas, clientes y reportes operativos, y validar su impacto en la eficiencia operativa mediante pruebas con empresas reales.

\subsection{Objetivos Específicos}\label{objetivos-especuxedficos}

\begin{enumerate}
\def\labelenumi{\arabic{enumi}.}
\item
  \textbf{Diseñar} la arquitectura técnica de una plataforma SaaS multi-tenant que garantice el aislamiento de datos entre organizaciones, la escalabilidad horizontal y la seguridad de la información, documentando las decisiones arquitectónicas con su justificación.
\item
  \textbf{Desarrollar} los módulos de gestión operativa esenciales ---autenticación y control de acceso basado en roles, gestión de inventario con alertas automáticas, registro de ventas con generación de documentos, y gestión de clientes con análisis de historial--- conforme a los requisitos levantados con usuarios reales de pymes.
\item
  \textbf{Construir} un sistema de reportes y analítica que proporcione indicadores clave de desempeño (KPIs) en tiempo real en el dashboard, junto con la capacidad de exportar a Excel los listados operativos (inventario, ventas y clientes), facilitando la toma de decisiones basada en datos.
\item
  \textbf{Desplegar} la plataforma en infraestructura de nube (Railway para el backend y la base de datos, Vercel para el frontend) con un pipeline de integración y entrega continua (CI/CD), garantizando una disponibilidad mínima del 95\% y tiempos de respuesta inferiores a 2 segundos para operaciones transaccionales.
\item
  \textbf{Validar} la solución mediante pruebas de usabilidad y rendimiento con al menos dos empresas piloto, midiendo el impacto en la eficiencia operativa a través de métricas cuantitativas (reducción de tiempo en tareas, tasa de error) y cualitativas (satisfacción de usuario).
\end{enumerate}

\section{Alcance y Limitaciones}\label{alcance-y-limitaciones}

\subsection{Alcance del Sistema}\label{alcance-del-sistema}

OrbitEngine abarca el siguiente conjunto de funcionalidades dentro del producto mínimo viable (MVP):

\textbf{Módulos implementados:}

\begin{itemize}
\tightlist
\item
  Autenticación y gestión de usuarios con control de acceso basado en roles (RBAC): Administrador, Vendedor y Visualizador.
\item
  Gestión de inventario: CRUD de productos por categorías, control de stock en tiempo real, alertas automáticas de nivel mínimo e historial de movimientos.
\item
  Gestión de ventas: registro de transacciones multi-producto con generación de número de factura, historial de ventas con filtros y búsqueda.
\item
  Gestión de clientes: base de datos de clientes, historial de compras y métricas de comportamiento.
\item
  Dashboard y reportes: KPIs en tiempo real (ventas del día/mes, stock bajo, top de productos, tendencia de ventas) y exportación a Excel de los listados de inventario, ventas y clientes.
\item
  Infraestructura multi-tenant: soporte para múltiples organizaciones con aislamiento total de datos.
\end{itemize}

\textbf{Características técnicas:}

\begin{itemize}
\tightlist
\item
  API RESTful documentada con OpenAPI/Swagger.
\item
  Interfaz web responsive, compatible con navegadores modernos (Chrome, Firefox, Safari, Edge).
\end{itemize}

\textbf{\emph{- Despliegue en Railway (backend) y Vercel (frontend) con GitHub Actions.}}

\begin{itemize}
\tightlist
\item
  Suite de pruebas automatizadas con cobertura mínima del 60\%.
\end{itemize}

\subsection{Limitaciones y Exclusiones}\label{limitaciones-y-exclusiones}

Las siguientes características quedan fuera del alcance del MVP y podrán abordarse en versiones futuras:

\begin{itemize}
\tightlist
\item
  Facturación electrónica oficial con validez tributaria (DIAN, SAT, SII u organismos equivalentes).
\item
  Gestión de nómina y recursos humanos.
\item
  Aplicación móvil nativa (iOS/Android).
\item
  Integración con pasarelas de pago o sistemas ERP externos.
\item
  Soporte para gestión de múltiples sucursales o bodegas dentro de una misma organización.
\item
  Gestión de servicios (el sistema está orientado a empresas de comercio de productos físicos).
\item
  Múltiples monedas o idiomas en la interfaz.
\end{itemize}

\subsection{Supuestos}\label{supuestos}

\begin{enumerate}
\def\labelenumi{\arabic{enumi}.}
\tightlist
\item
  Los usuarios de la plataforma disponen de conexión a internet estable para el acceso a la aplicación web.
\item
  Las empresas piloto están dispuestas a cargar sus datos de productos, clientes y stock inicial para comenzar a operar con el sistema.
\item
  El equipo de desarrollo cuenta con cuentas activas en Railway y Vercel con planes adecuados para el despliegue de la plataforma.
\end{enumerate}

\section{Metodología}\label{metodologuxeda}

El proyecto se desarrolló bajo un enfoque de investigación aplicada con componente de desarrollo de software. Se adoptó la metodología ágil Scrum adaptada para equipos académicos, organizando el desarrollo en sprints de dos semanas durante un período de siete meses (octubre 2025 -- abril 2026).

La investigación siguió las siguientes etapas:

\begin{enumerate}
\def\labelenumi{\arabic{enumi}.}
\tightlist
\item
  \textbf{Documentación y planificación} (octubre 2025): revisión bibliográfica, levantamiento de requisitos con usuarios potenciales, diseño de arquitectura, modelo de datos y definición del backlog inicial.
\item
  \textbf{Desarrollo del núcleo} (noviembre 2025 -- enero 2026): implementación de la infraestructura base y los módulos de autenticación, inventario, ventas, clientes y reportes.
\item
  \textbf{Estabilización y refinamiento} (febrero -- marzo 2026): corrección de errores, mejoras de experiencia de usuario, pruebas de carga y consolidación del despliegue en producción.
\item
  \textbf{Validación con empresas piloto} (abril 2026, semanas 1--3): pruebas con usuarios reales, recolección de métricas de eficiencia operativa y encuestas de usabilidad.
\item
  \textbf{Documentación y entrega} (abril 2026, semana 4): consolidación del informe final, análisis de resultados y defensa del proyecto.
\end{enumerate}

\section{Estructura del Documento}\label{estructura-del-documento}

El presente informe se organiza en los siguientes capítulos:

\begin{itemize}
\tightlist
\item
  \textbf{Capítulo 2 --- Marco de Referencia}: define los conceptos fundamentales del dominio (pyme, ERP, SaaS, multi-tenancy, API REST, RBAC) y presenta el estado del arte mediante una revisión de soluciones existentes en el mercado y de la literatura académica relevante.
\item
  \textbf{Capítulo 3 --- Análisis y Diseño del Sistema}: describe el proceso de levantamiento de requisitos, la arquitectura técnica adoptada, el diseño del modelo de datos y el diseño de las interfaces de usuario.
\item
  \textbf{Capítulo 4 --- Desarrollo e Implementación}: detalla la metodología de desarrollo utilizada, los módulos implementados, las decisiones técnicas relevantes y la estrategia de pruebas.
\item
  \textbf{Capítulo 5 --- Resultados y Validación}: presenta los resultados de las pruebas de rendimiento, usabilidad y validación con empresas piloto, contrastándolos con los objetivos planteados.
\item
  \textbf{Capítulo 6 --- Conclusiones y Trabajo Futuro}: sintetiza los logros del proyecto por objetivo, enumera las limitaciones encontradas y propone líneas de trabajo para versiones futuras.
\item
  \textbf{Referencias Bibliográficas}: lista completa de las fuentes citadas a lo largo del documento.
\item
  \textbf{Anexos}: manual de usuario, manual de despliegue, especificación de la API y resultados detallados de pruebas.
\end{itemize}

    \chapter{Marco de Referencia}\label{marco-de-referencia}

\begin{center}\rule{0.5\linewidth}{0.5pt}\end{center}

\section{Marco Conceptual}\label{marco-conceptual}

Esta sección define los conceptos fundamentales sobre los que se sustenta el proyecto OrbitEngine, estableciendo el vocabulario técnico y académico compartido a lo largo del documento.

\subsection{Pequeña y Mediana Empresa (Pyme)}\label{pequeuxf1a-y-mediana-empresa-pyme}

La definición de pyme varía según el país y el organismo de referencia. En América Latina, los criterios más comunes consideran el número de empleados y el volumen de ventas anuales. La CEPAL (2022) y el Banco Mundial utilizan como referencia general la siguiente clasificación:

{\def\LTcaptype{none} % do not increment counter
\begin{longtable}[]{@{}lll@{}}
\toprule\noalign{}
Categoría & Empleados & Ventas anuales (USD) \\
\midrule\noalign{}
\endhead
\bottomrule\noalign{}
\endlastfoot
Microempresa & 1--9 & \textless{} 100.000 \\
Pequeña empresa & 10--49 & 100.000 -- 1.000.000 \\
Mediana empresa & 50--249 & 1.000.000 -- 10.000.000 \\
\end{longtable}
}

Para el contexto de este proyecto, el término ``pyme'' se utiliza con énfasis en el segmento de pequeñas empresas de comercio minorista y mayorista, con equipos de entre 2 y 30 personas, operaciones presenciales o híbridas, y necesidades de gestión de inventario físico, registro de ventas y seguimiento de clientes.

\subsection{Sistema de Planificación de Recursos Empresariales (ERP)}\label{sistema-de-planificaciuxf3n-de-recursos-empresariales-erp}

Un Enterprise Resource Planning (ERP) es un sistema de información integrado que unifica los procesos de negocio de una organización ---finanzas, producción, inventario, ventas, recursos humanos--- en una única plataforma con una base de datos compartida (Klaus et al., 2000). La integración elimina la duplicación de datos y proporciona una visión unificada del estado del negocio en tiempo real.

Los ERP surgen en los años noventa como evolución de los sistemas MRP (Material Requirements Planning). Empresas como SAP, Oracle y Microsoft Dynamics dominaron el mercado inicial con soluciones de altísimo costo y complejidad, orientadas a grandes corporaciones. La segunda generación de ERP, surgida en la década de 2000, introdujo soluciones para el mercado de pymes (SAP Business One, Microsoft Dynamics 365 Business Central), si bien con costos de implementación y operación aún inaccesibles para la mayoría de las pequeñas empresas latinoamericanas.

OrbitEngine no se clasifica como un ERP en el sentido amplio del término ---no incluye módulos de contabilidad, nómina ni producción---, sino como una plataforma de gestión operativa orientada a los procesos comerciales centrales de una pyme: inventario, ventas y clientes.

\subsection{Software como Servicio (SaaS)}\label{software-como-servicio-saas}

Software as a Service (SaaS) es un modelo de distribución de software en el que el proveedor aloja la aplicación en infraestructura propia (nube) y la pone a disposición de los usuarios finales a través de internet, típicamente mediante una suscripción periódica (Benlian \& Hess, 2011).

El modelo SaaS elimina para el cliente la necesidad de:

\begin{itemize}
\tightlist
\item
  Adquirir y mantener infraestructura de servidores.
\item
  Gestionar instalaciones, actualizaciones y parches del software.
\item
  Invertir en licencias únicas de alto costo.
\end{itemize}

En cambio, el cliente paga una tarifa recurrente (mensual o anual) que cubre el uso del software, la infraestructura, el mantenimiento y el soporte. Este modelo democratiza el acceso a herramientas empresariales sofisticadas para organizaciones sin departamento de TI.

Desde la perspectiva del proveedor, el modelo SaaS permite economías de escala: una misma infraestructura sirve a múltiples clientes (multi-tenancy), reduciendo el costo marginal por cliente adicional.

\subsection{Arquitectura Multi-Tenant}\label{arquitectura-multi-tenant}

El concepto de multi-tenancy (multi-arrendamiento) describe una arquitectura de software en la que una única instancia de la aplicación y su infraestructura subyacente sirven a múltiples clientes (tenants), manteniendo el aislamiento lógico de los datos de cada uno (Bezemer \& Zaidman, 2010).

Existen tres enfoques principales de implementación:

{\def\LTcaptype{none} % do not increment counter
\begin{longtable}[]{@{}
  >{\raggedright\arraybackslash}p{(\linewidth - 6\tabcolsep) * \real{0.3516}}
  >{\raggedright\arraybackslash}p{(\linewidth - 6\tabcolsep) * \real{0.4725}}
  >{\raggedright\arraybackslash}p{(\linewidth - 6\tabcolsep) * \real{0.1209}}
  >{\raggedright\arraybackslash}p{(\linewidth - 6\tabcolsep) * \real{0.0549}}@{}}
\toprule\noalign{}
\begin{minipage}[b]{\linewidth}\raggedright
Enfoque
\end{minipage} & \begin{minipage}[b]{\linewidth}\raggedright
Descripción
\end{minipage} & \begin{minipage}[b]{\linewidth}\raggedright
Aislamiento
\end{minipage} & \begin{minipage}[b]{\linewidth}\raggedright
Costo
\end{minipage} \\
\midrule\noalign{}
\endhead
\bottomrule\noalign{}
\endlastfoot
Base de datos por tenant & Cada cliente tiene su propia BD & Alto & Alto \\
Esquema por tenant & Una BD, esquemas separados & Medio & Medio \\
Tabla compartida + discriminador & Una BD, filas diferenciadas por \texttt{tenant\_id} & Bajo-Medio & Bajo \\
\end{longtable}
}

OrbitEngine implementa el tercer enfoque: todas las organizaciones comparten las mismas tablas, y cada registro lleva asociado un campo \texttt{organization\_id} que actúa como discriminador. El acceso a los datos se filtra automáticamente en cada operación, garantizando el aislamiento sin multiplicar el costo de infraestructura.

\subsection{API REST y OpenAPI}\label{api-rest-y-openapi}

Una API REST (Representational State Transfer) es una interfaz de programación que permite la comunicación entre sistemas a través del protocolo HTTP, utilizando los verbos estándar (GET, POST, PUT, PATCH, DELETE) y representaciones de recursos en formato JSON o XML (Fielding, 2000).

OpenAPI (anteriormente Swagger) es una especificación estándar para describir APIs REST de forma legible tanto por humanos como por máquinas. Su adopción en OrbitEngine permite la generación automática de documentación interactiva y de clientes de API tipados para el frontend.

\subsection{Control de Acceso Basado en Roles (RBAC)}\label{control-de-acceso-basado-en-roles-rbac}

Role-Based Access Control (RBAC) es un modelo de seguridad en el que los permisos de acceso a recursos del sistema se asignan a roles, y los usuarios adquieren permisos a través de su asignación a uno o más roles (Ferraiolo et al., 2001). Este modelo simplifica la administración de permisos en organizaciones donde múltiples usuarios realizan funciones diferenciadas.

En OrbitEngine se definen tres roles:

\begin{itemize}
\tightlist
\item
  \textbf{Administrador}: acceso total a todos los módulos y configuración de la organización.
\item
  \textbf{Vendedor}: acceso a ventas, consulta de inventario y clientes. Sin acceso a reportes avanzados ni configuración.
\item
  \textbf{Visualizador}: acceso de solo lectura a reportes y dashboards.
\end{itemize}

\subsection{Integración y Entrega Continua (CI/CD)}\label{integraciuxf3n-y-entrega-continua-cicd}

Continuous Integration/Continuous Delivery (CI/CD) es un conjunto de prácticas de ingeniería de software que automatiza el proceso de verificación, construcción y despliegue del código. Un pipeline de CI/CD ejecuta, ante cada cambio en el repositorio, una secuencia de pasos: pruebas automatizadas, análisis estático, construcción de artefactos y despliegue en los ambientes correspondientes (Fowler \& Foemmel, 2006).

La adopción de CI/CD en OrbitEngine, implementada mediante GitHub Actions, garantiza que cada versión desplegada en producción ha pasado exitosamente por la suite de pruebas, reduciendo la probabilidad de introducir regresiones.

\section{Estado del Arte}\label{estado-del-arte}

\subsection{Contexto de la Digitalización de Pymes en Latinoamérica}\label{contexto-de-la-digitalizaciuxf3n-de-pymes-en-latinoamuxe9rica}

La transformación digital de las pymes en América Latina ha sido objeto de creciente atención académica y política en la última década. La pandemia de COVID-19 (2020--2021) actuó como catalizador de la adopción tecnológica, al forzar a empresas que operaban exclusivamente de manera presencial a explorar canales digitales y herramientas de gestión remota.

Rodríguez-Abitia et al.~(2020), en su estudio sobre digitalización de pymes mexicanas, identificaron que los principales obstáculos para la adopción de sistemas de gestión son: el costo percibido como elevado, la falta de capacitación del personal, la desconfianza en la seguridad de datos en la nube y la baja oferta de soluciones adaptadas al contexto local. Un hallazgo relevante es que las pymes que sí adoptaron herramientas digitales de gestión reportaron mejoras significativas en eficiencia (entre el 20\% y el 35\%) en los 12 meses posteriores a la implementación.

Cardona et al.~(2022) analizaron la adopción de ERP en pymes colombianas y encontraron que el 68\% de las empresas que fallaron en su implementación lo hicieron por seleccionar soluciones demasiado complejas para su tamaño, o por falta de acompañamiento en la fase de puesta en marcha. Los autores concluyen que la simplicidad de la interfaz y la pertinencia funcional son más determinantes del éxito que las capacidades técnicas del sistema.

\subsection{Soluciones Existentes en el Mercado}\label{soluciones-existentes-en-el-mercado}

A continuación se analiza el panorama de soluciones que compiten, directa o indirectamente, con OrbitEngine en el segmento de pymes latinoamericanas.

\subsubsection{Odoo Community / Enterprise}\label{odoo-community-enterprise}

Odoo es un ERP de código abierto con más de 7 millones de usuarios en todo el mundo. Su versión Community es gratuita y cubre un amplio espectro de módulos: inventario, ventas, compras, contabilidad, CRM y recursos humanos. La versión Enterprise agrega módulos avanzados bajo un modelo de suscripción.

\textbf{Fortalezas}: altamente modular, comunidad amplia, documentación extensa.\\
\textbf{Debilidades para el contexto de este proyecto}: requiere instalación en servidor propio o contratación de hosting especializado, lo que implica conocimiento técnico o un costo adicional de servicio. La curva de aprendizaje para usuarios no técnicos es elevada. La personalización para contextos latinoamericanos (localización fiscal, moneda local) requiere módulos adicionales o desarrollo a medida. No ofrece análisis de clientes integrado ni dashboard operativo simplificado.

\subsubsection{Zoho One}\label{zoho-one}

Zoho One es una suite de aplicaciones SaaS que incluye más de 45 herramientas de negocio: CRM, inventario, contabilidad, proyectos y análisis. Es una solución madura con presencia global.

\textbf{Fortalezas}: modelo SaaS puro, integración entre módulos, amplia cobertura funcional.\\
\textbf{Debilidades}: el costo por usuario (\$37--\$105 USD/mes) resulta elevado para pymes con equipos numerosos. La interfaz, aunque completa, puede resultar abrumadora para usuarios con bajo nivel de alfabetización digital.

\subsubsection{QuickBooks}\label{quickbooks}

QuickBooks, de Intuit, es el software de contabilidad y gestión financiera más utilizado por pequeñas empresas en el mercado anglosajón. Su foco principal es la contabilidad, facturación y control de gastos.

\textbf{Fortalezas}: interfaz intuitiva, amplia base de usuarios, integración bancaria.\\
\textbf{Debilidades}: diseñado para el mercado norteamericano y europeo, con limitada adaptación a las realidades tributarias latinoamericanas. No ofrece gestión de inventario avanzada ni análisis de clientes. Precios en dólares que pueden ser inaccesibles para pymes de ingresos bajos.

\subsubsection{Alegra}\label{alegra}

Alegra es una solución SaaS desarrollada en Colombia con presencia en 10 países de América Latina. Se especializa en facturación electrónica, contabilidad y gestión de inventarios básica.

\textbf{Fortalezas}: fuerte localización tributaria latinoamericana, precio accesible (desde \$25 USD/mes), interfaz en español, soporte en la región.\\
\textbf{Debilidades}: las capacidades de gestión de inventario son básicas (sin alertas de stock mínimo ni historial de movimientos detallado). No incluye CRM ni análisis de clientes. El dashboard de indicadores operativos es limitado.

\subsubsection{Siigo / Defontana}\label{siigo-defontana}

Siigo (Colombia) y Defontana (Chile) son soluciones SaaS con foco contable-tributario, similares a Alegra en su posicionamiento. Están bien adaptadas a la normativa local de sus países de origen pero tienen alcance funcional limitado en gestión operativa.

\subsubsection{Análisis Comparativo}\label{anuxe1lisis-comparativo}

La siguiente tabla sintetiza la comparación entre las soluciones analizadas y OrbitEngine en las dimensiones relevantes para el segmento objetivo:

{\def\LTcaptype{none} % do not increment counter
\begin{longtable}[]{@{}
  >{\raggedright\arraybackslash}p{(\linewidth - 10\tabcolsep) * \real{0.4022}}
  >{\raggedright\arraybackslash}p{(\linewidth - 10\tabcolsep) * \real{0.1848}}
  >{\raggedright\arraybackslash}p{(\linewidth - 10\tabcolsep) * \real{0.0761}}
  >{\raggedright\arraybackslash}p{(\linewidth - 10\tabcolsep) * \real{0.1087}}
  >{\raggedright\arraybackslash}p{(\linewidth - 10\tabcolsep) * \real{0.0652}}
  >{\raggedright\arraybackslash}p{(\linewidth - 10\tabcolsep) * \real{0.1630}}@{}}
\toprule\noalign{}
\begin{minipage}[b]{\linewidth}\raggedright
Criterio
\end{minipage} & \begin{minipage}[b]{\linewidth}\raggedright
Odoo
\end{minipage} & \begin{minipage}[b]{\linewidth}\raggedright
Zoho
\end{minipage} & \begin{minipage}[b]{\linewidth}\raggedright
QuickBooks
\end{minipage} & \begin{minipage}[b]{\linewidth}\raggedright
Alegra
\end{minipage} & \begin{minipage}[b]{\linewidth}\raggedright
OrbitEngine
\end{minipage} \\
\midrule\noalign{}
\endhead
\bottomrule\noalign{}
\endlastfoot
Modelo de entrega & On-premise / SaaS & SaaS & SaaS & SaaS & SaaS \\
Precio accesible para pymes pequeñas & Medio & No & No & Sí & Sí (objetivo) \\
Gestión de inventario avanzada & Sí & Parcial & No & Básica & Sí \\
Gestión de clientes (CRM) & Sí & Sí & Básica & No & Sí \\
Dashboard y KPIs integrados & Sí & Sí & Parcial & Básico & Sí \\
Localización latinoamericana & Parcial & Parcial & Baja & Alta & Media \\
Curva de aprendizaje para no técnicos & Alta & Alta & Media & Baja & Baja (objetivo) \\
Multi-tenant nativo & Sí & Sí & No & No & Sí \\
\end{longtable}
}

El análisis evidencia que ninguna de las soluciones disponibles combina en un único producto, accesible para el segmento de pymes latinoamericanas, la gestión operativa completa (inventario, ventas, clientes) con un dashboard de KPIs integrado, bajo una curva de aprendizaje adecuada para usuarios no técnicos. Esta brecha constituye el espacio que OrbitEngine busca ocupar.

\subsection{Revisión de Literatura Académica}\label{revisiuxf3n-de-literatura-acaduxe9mica}

\subsubsection{ERP en Pymes: Factores de Éxito y Fracaso}\label{erp-en-pymes-factores-de-uxe9xito-y-fracaso}

La adopción de ERP en pequeñas y medianas empresas ha sido estudiada extensamente. Soh et al.~(2000) identificaron la brecha entre las funcionalidades genéricas de los sistemas ERP y las necesidades específicas de cada organización como la principal causa de fracaso en las implementaciones. Esta observación motivó el diseño de OrbitEngine como una solución vertical, enfocada en el nicho de comercio minorista y mayorista, en lugar de un ERP de propósito general.

Amid et al.~(2012) propusieron un marco de factores críticos de éxito para implementaciones de ERP en pymes, destacando: (1) compromiso de la alta dirección, (2) simplicidad del proceso de implementación, (3) adaptación al contexto local y (4) capacitación de usuarios. Estos factores influyeron en las decisiones de diseño de OrbitEngine, especialmente en la priorización de la usabilidad y el proceso de onboarding.

\subsubsection{Arquitecturas SaaS Multi-Tenant}\label{arquitecturas-saas-multi-tenant}

Bezemer \& Zaidman (2010) analizaron los desafíos de migrar aplicaciones monolíticas a arquitecturas multi-tenant, identificando como principales retos: (1) el aislamiento de datos, (2) la personalización por tenant y (3) el rendimiento bajo carga compartida. Su trabajo sienta las bases para las decisiones de diseño del modelo de datos de OrbitEngine.

Guo et al.~(2017) propusieron un modelo de evaluación de arquitecturas multi-tenant en función de métricas de seguridad, rendimiento y mantenibilidad. Aplicando su framework, la arquitectura de OrbitEngine (tabla compartida con \texttt{organization\_id}) se clasifica como de nivel de compartición alto, lo que maximiza la eficiencia de infraestructura al costo de una implementación más cuidadosa de los mecanismos de filtrado.

\subsubsection{Usabilidad en Sistemas ERP para Usuarios No Técnicos}\label{usabilidad-en-sistemas-erp-para-usuarios-no-tuxe9cnicos}

Zhang et al.~(2021) realizaron un estudio de usabilidad con empleados de pymes latinoamericanas usando tres sistemas ERP distintos. Sus resultados mostraron que la tasa de abandono durante el onboarding aumenta un 45\% cuando el sistema requiere más de 5 pasos para completar una tarea común (como registrar una venta). Este hallazgo motivó el énfasis en la experiencia de usuario (UX) de OrbitEngine, con flujos de trabajo optimizados para minimizar la fricción.

\section{Marco Tecnológico}\label{marco-tecnoluxf3gico}

Esta sección justifica las decisiones tecnológicas adoptadas en el proyecto, situándolas en el contexto de las alternativas disponibles y los criterios de selección.

\subsection{FastAPI (Backend)}\label{fastapi-backend}

FastAPI es un framework web moderno para Python, basado en Starlette y Pydantic, que permite construir APIs de alto rendimiento con tipado estático y generación automática de documentación OpenAPI (Ramírez, 2018). Su selección sobre alternativas como Django REST Framework o Flask se justifica por:

\begin{itemize}
\tightlist
\item
  \textbf{Rendimiento}: comparable a Node.js y Go en benchmarks de concurrencia, gracias a su implementación asíncrona (ASGI).
\item
  \textbf{Tipado automático}: la integración con Pydantic valida automáticamente los datos de entrada y salida de cada endpoint, reduciendo el boilerplate de validación.
\item
  \textbf{Documentación automática}: genera una interfaz Swagger UI y ReDoc sin configuración adicional, facilitando la comunicación con el frontend y la evaluación académica.
\item
  \textbf{Ecosistema Python maduro}: la elección de Python como lenguaje de backend otorga acceso a un amplio ecosistema de bibliotecas para procesamiento de datos, generación de reportes y exportación de datos tabulares, y posibles extensiones futuras.
\end{itemize}

\subsection{React (Frontend)}\label{react-frontend}

React es una biblioteca JavaScript para la construcción de interfaces de usuario basadas en componentes, desarrollada por Meta. Su adopción como tecnología de frontend es ampliamente respaldada por la industria y la academia (Gackenheimer, 2015).

La combinación con TypeScript añade tipado estático al frontend, mejorando la detección temprana de errores y la mantenibilidad del código. Las bibliotecas complementarias seleccionadas (TanStack Query para gestión de estado del servidor, TanStack Router para enrutamiento, React Hook Form + Zod para formularios con validación) representan el estándar actual de la industria para aplicaciones SaaS de escala media.

\subsection{PostgreSQL}\label{postgresql}

PostgreSQL es el sistema de gestión de bases de datos relacionales de código abierto más avanzado disponible. Su elección sobre alternativas como MySQL o bases de datos NoSQL se fundamenta en:

\begin{itemize}
\tightlist
\item
  \textbf{Integridad referencial}: las relaciones entre entidades de negocio (ventas → productos, ventas → clientes) se benefician de las garantías ACID de un motor relacional.
\item
  \textbf{JSON nativo}: el soporte para columnas JSONB permite almacenar datos semi-estructurados (atributos variables de productos) sin sacrificar las ventajas relacionales.
\item
  \textbf{Row-Level Security}: característica nativa de PostgreSQL que permite definir políticas de acceso a nivel de fila, complementando la estrategia de multi-tenancy a nivel de aplicación.
\item
  \textbf{Ecosistema robusto}: integración nativa con SQLAlchemy/SQLModel (ORM Python) y soporte completo en plataformas gestionadas como Railway.
\end{itemize}

\subsection{Docker, Railway y Vercel}\label{docker-railway-y-vercel}

La contenerización mediante Docker garantiza la paridad entre los ambientes de desarrollo, pruebas y producción, eliminando la clase de errores ``funciona en mi máquina''. El despliegue en producción se apoya en los siguientes servicios:

\begin{itemize}
\tightlist
\item
  \textbf{Railway} (backend): plataforma PaaS que gestiona el ciclo de vida de los contenedores del backend, desplegando automáticamente desde el repositorio Git ante cada merge a la rama principal.
\item
  \textbf{Railway} (base de datos): instancia gestionada de PostgreSQL con backups automáticos y variables de conexión inyectadas directamente en el entorno de ejecución.
\item
  \textbf{Vercel} (frontend): plataforma especializada en el despliegue de aplicaciones frontend, con CDN global integrado, HTTPS automático y previsualizaciones por pull request.
\end{itemize}

\subsection{SQLModel}\label{sqlmodel}

SQLModel es una biblioteca Python que combina SQLAlchemy (ORM) y Pydantic en un único modelo de clase, permitiendo que el mismo tipo Python sirva tanto como modelo de base de datos como schema de validación de la API (Ramírez, 2021). Esta unificación reduce la duplicación de datos y asegura la consistencia entre la capa de datos y la capa de API.

    \input{MainMatter/Cap3/C3-AnalisisDiseno}
    \chapter{Capítulo 4. Desarrollo e Implementación}\label{capuxedtulo-4.-desarrollo-e-implementaciuxf3n}

\begin{center}\rule{0.5\linewidth}{0.5pt}\end{center}

\section{Metodología de Desarrollo}\label{metodologuxeda-de-desarrollo}

\subsection{Marco Ágil: Scrum Adaptado}\label{marco-uxe1gil-scrum-adaptado}

El desarrollo de OrbitEngine siguió un proceso ágil basado en el framework Scrum, adaptado a las condiciones de un equipo académico de tres personas con dedicación parcial. Las adaptaciones principales respecto al Scrum estándar fueron:

\begin{itemize}
\tightlist
\item
  \textbf{Sprints de 2 semanas} (en lugar de 1 semana) para reducir la sobrecarga de ceremonias y dar margen a las obligaciones académicas paralelas.
\item
  \textbf{Daily standups asincrónicos} mediante mensajes en un canal dedicado de Discord, dado que los horarios de los integrantes no siempre coincidían.
\item
  \textbf{Roles simplificados}: no se designó un Scrum Master formal; los tres integrantes rotaron informalmente la facilitación de las ceremonias.
\end{itemize}

Las ceremonias mantenidas fueron:

\begin{itemize}
\tightlist
\item
  \textbf{Sprint Planning} (inicio del sprint, 2 horas): revisión del backlog priorizado, estimación con puntos de historia en escala Fibonacci y asignación de tareas.
\item
  \textbf{Sprint Review} (fin del sprint, 1 hora): demostración de las funcionalidades completadas.
\item
  \textbf{Sprint Retrospective} (fin del sprint, 30 minutos): identificación de mejoras al proceso.
\end{itemize}

\subsection{Herramientas del Proceso}\label{herramientas-del-proceso}

\begin{longtable}[]{@{}
  >{\raggedright\arraybackslash}p{(\linewidth - 2\tabcolsep) * \real{0.2941}}
  >{\raggedright\arraybackslash}p{(\linewidth - 2\tabcolsep) * \real{0.7059}}@{}}
\caption[Herramientas del Proceso]{Herramientas del Proceso}\tabularnewline
\toprule\noalign{}
\begin{minipage}[b]{\linewidth}\raggedright
Herramienta
\end{minipage} & \begin{minipage}[b]{\linewidth}\raggedright
Propósito
\end{minipage} \\
\midrule\noalign{}
\endhead
\bottomrule\noalign{}
\endlastfoot
GitHub Projects & Gestión del backlog, kanban de sprints \\
GitHub (repositorio) & Control de versiones, pull requests, code review \\
GitHub Actions & CI/CD: pruebas automatizadas y despliegue \\
Discord & Comunicación del equipo, standups asincrónicos \\
Figma & Diseño de wireframes y mockups de UI \\
\end{longtable}

\subsection{Estructura del Equipo}\label{estructura-del-equipo}

El equipo se organizó en roles complementarios con zonas de especialización clara, aunque todos los integrantes contribuyeron en múltiples capas del sistema:

\begin{itemize}
\tightlist
\item
  \textbf{Backend Lead}: arquitectura del sistema, endpoints de la API, modelos de base de datos, lógica de negocio y reportes.
\item
  \textbf{Frontend Lead}: componentes de UI, flujos de usuario, integración con la API, diseño responsivo.
\item
  \textbf{DevOps / Full Stack}: infraestructura en Railway y Vercel, pipelines CI/CD, monitoreo, soporte en backend y frontend.
\end{itemize}

\subsection{Gestión de la Calidad del Código}\label{gestiuxf3n-de-la-calidad-del-cuxf3digo}

Se establecieron las siguientes prácticas de calidad desde el inicio del proyecto:

\textbf{Backend:}

\begin{itemize}
\tightlist
\item
  Tipado estático completo con mypy (modo strict).
\item
  Linting y formateo con Ruff.
\item
  Revisión de código mediante pull requests antes de hacer merge a \texttt{main}.
\item
  Cobertura de pruebas mínima del 60\% como condición para que el CI apruebe.
\end{itemize}

\textbf{Frontend:}

\begin{itemize}
\tightlist
\item
  Tipado estricto con TypeScript.
\item
  Linting y formateo con Biome.
\item
  Generación automática del cliente de API desde el schema OpenAPI (\texttt{bun\ run\ generate-client}) para garantizar la sincronización entre frontend y backend.
\end{itemize}

\section{Fases y Sprints de Desarrollo}\label{fases-y-sprints-de-desarrollo}

El proyecto se estructuró en seis fases que cubrieron el período comprendido entre octubre de 2025 y mediados de mayo de 2026.

\subsection{Fase 1. Documentación e Investigación (octubre 2025, semanas 1--3)}\label{fase-1.-documentaciuxf3n-e-investigaciuxf3n-octubre-2025-semanas-13}

Esta fase se dedicó a la comprensión del dominio del problema, el levantamiento de requisitos y la selección del stack tecnológico. Las actividades principales incluyeron:

\begin{itemize}
\tightlist
\item
  Revisión de literatura sobre digitalización de pymes y soluciones ERP/SaaS existentes.
\item
  Levantamiento informal de información a partir de la experiencia previa de los integrantes del equipo con pymes del sector comercio, sin un proceso estructurado de entrevistas externas.
\item
  Análisis comparativo de plataformas competidoras.
\item
  Elaboración del backlog inicial con historias de usuario estimadas.
\end{itemize}

\textbf{Entregables}: propuesta de proyecto aprobada, backlog priorizado, documento de requisitos.

\subsection{Fase 2. Diseño y Arquitectura (octubre--noviembre 2025, semanas 4--5)}\label{fase-2.-diseuxf1o-y-arquitectura-octubrenoviembre-2025-semanas-45}

Con el backlog definido, se realizó el diseño técnico del sistema:

\begin{itemize}
\tightlist
\item
  Diseño del modelo de datos (diagrama ER, diccionario de tablas).
\item
  Especificación de la API RESTful en formato OpenAPI.
\item
  Wireframes y mockups de alta fidelidad de las vistas principales.
\item
  Configuración del repositorio, estructura de carpetas, entorno de desarrollo local con Docker Compose y pipeline base de CI en GitHub Actions (lint y type-check).
\end{itemize}

\textbf{Entregables}: diseño de base de datos completo, mockups aprobados, repositorio inicializado con CI básico y entorno local reproducible vía Docker Compose.

\subsection{Fase 3. Desarrollo Core (noviembre 2025 a segunda semana de abril 2026, sprints 1--8)}\label{fase-3.-desarrollo-core-noviembre-2025-a-segunda-semana-de-abril-2026-sprints-18}

Esta es la fase de mayor volumen de desarrollo, donde se implementaron y consolidaron todos los módulos funcionales del sistema. Se extendió desde la primera semana de noviembre de 2025 hasta la segunda semana de abril de 2026 para permitir cerrar el alcance funcional completo antes de pasar a la fase de estabilización. Los ocho sprints se distribuyeron en duraciones de \textbf{dos o tres semanas} según el tamaño de cada bloque de funcionalidad y, entre los Sprints 3 y 4, se respetó un receso académico de dos semanas (22 de diciembre de 2025 -- 4 de enero de 2026) en el que el equipo no ejecutó iteraciones formales.

\subsubsection{Sprint 1: Autenticación y Setup Base (3--14 noviembre 2025 \textbar{} 18 SP)}\label{sprint-1-autenticaciuxf3n-y-setup-base-314-noviembre-2025-18-sp}

\begin{itemize}
\tightlist
\item
  Backend: modelos \texttt{Organization} y \texttt{User}, endpoints \texttt{POST\ /organizations/signup} (alta de organización con su usuario administrador) y \texttt{POST\ /login/access-token}, generación y validación de JWT con \texttt{sub}, \texttt{organization\_id}, \texttt{role} y \texttt{exp}, dependencia \texttt{get\_current\_user} y \texttt{require\_role} para control de acceso por rol.
\item
  Frontend: setup de Vite + React + TypeScript, TanStack Router con rutas basadas en archivos, gestión del token de sesión vía un módulo propio (\texttt{lib/auth-session}) integrado con TanStack Query, pantallas de login y registro de organización.
\item
  Resultado: flujo de autenticación end-to-end funcional con multi-tenancy desde el primer commit productivo.
\end{itemize}

\subsubsection{Sprint 2: Inventario Core, CRUD de Productos (17 noviembre a 5 diciembre 2025 \textbar{} 17 SP)}\label{sprint-2-inventario-core-crud-de-productos-17-noviembre-a-5-diciembre-2025-17-sp}

\begin{itemize}
\tightlist
\item
  Backend: modelos \texttt{Product} y \texttt{Category}, CRUD completo con paginación, búsqueda y filtros, soft delete (\texttt{deleted\_at}).
\item
  Frontend: tabla de productos con búsqueda en tiempo real, formulario de alta/edición con validación Zod + React Hook Form, modal de confirmación de eliminación.
\item
  Resultado: gestión de catálogo de productos operativa en entorno de desarrollo local.
\end{itemize}

\subsubsection{Sprint 3: Inventario Avanzado (8--19 diciembre 2025 \textbar{} 16 SP)}\label{sprint-3-inventario-avanzado-819-diciembre-2025-16-sp}

\begin{itemize}
\tightlist
\item
  Backend: modelo \texttt{InventoryMovement} con tipos \texttt{sale}, \texttt{purchase}, \texttt{adjustment} y \texttt{return}; endpoint \texttt{GET\ /products/low-stock} para listar productos por debajo del mínimo; endpoint \texttt{POST\ /inventory-movements/} para movimientos manuales (ajuste, compra, devolución).
\item
  Frontend: historial de movimientos por producto, widget de productos con stock bajo en el dashboard, formulario de ajuste manual con campo de justificación.
\item
  Resultado: trazabilidad completa de movimientos de stock.
\end{itemize}

\subsubsection{Sprint 4: Módulo de Ventas (5--23 enero 2026 \textbar{} 20 SP)}\label{sprint-4-muxf3dulo-de-ventas-523-enero-2026-20-sp}

\begin{itemize}
\tightlist
\item
  Backend: modelos \texttt{Sale} y \texttt{SaleItem}, lógica de descuento automático de stock al registrar venta, generación de número de factura secuencial por organización, endpoint \texttt{POST\ /sales/\{sale\_id\}/cancel} para anular ventas con reversión de stock vía movimientos de tipo \texttt{return}.
\item
  Frontend: flujo de registro de venta con búsqueda de productos por nombre/SKU, cálculo de totales en tiempo real, historial de ventas con filtros por fecha, estado y método de pago.
\item
  Resultado: módulo de ventas completo con integración automática al inventario.
\end{itemize}

\subsubsection{Sprint 5: Módulo de Clientes (26 enero -- 13 febrero 2026 \textbar{} 15 SP)}\label{sprint-5-muxf3dulo-de-clientes-26-enero-13-febrero-2026-15-sp}

\begin{itemize}
\tightlist
\item
  Backend: modelo \texttt{Customer}, CRUD de clientes, asociación opcional de ventas a clientes y actualización automática de métricas (total comprado, número de compras, ticket promedio) al registrar o cancelar una venta.
\item
  Frontend: lista de clientes con búsqueda, formulario de alta/edición, perfil de cliente con historial de compras y estadísticas.
\item
  Resultado: base de datos de clientes con análisis de comportamiento.
\end{itemize}

\subsubsection{Sprint 6: Reportes y Dashboard (16 febrero -- 6 marzo 2026 \textbar{} 18 SP)}\label{sprint-6-reportes-y-dashboard-16-febrero-6-marzo-2026-18-sp}

\begin{itemize}
\tightlist
\item
  Backend: endpoint \texttt{GET\ /dashboard/stats} con KPIs agregados (ventas del día, ventas del mes, conteo de productos con stock bajo, ticket promedio, top productos, ventas por día) y endpoint \texttt{POST\ /dashboard/export-excel} que genera archivos \texttt{.xlsx} para los datasets de inventario, clientes y ventas construyendo el formato Office Open XML manualmente sin dependencias externas.
\item
  Frontend: dashboard con widgets de KPIs, gráfico de ventas de los últimos días con Recharts, módulo de reportes con filtros de fecha y exportación a Excel.
\item
  Resultado: dashboard de KPIs operativo y módulo de exportación a Excel de los tres listados principales.
\end{itemize}

\subsubsection{Sprint 7: Roles, Permisos y Refinamiento Funcional (9--27 marzo 2026 \textbar{} 16 SP)}\label{sprint-7-roles-permisos-y-refinamiento-funcional-927-marzo-2026-16-sp}

\begin{itemize}
\tightlist
\item
  Backend: endpoints \texttt{GET\ /roles/} y consolidación del rol \texttt{contador} además de \texttt{admin}, \texttt{seller} y \texttt{viewer}; ajustes de permisos en endpoints sensibles (exportación, gestión de usuarios, cancelación de ventas) usando la dependencia \texttt{require\_role}.
\item
  Frontend: control de visibilidad de menús y acciones según el rol del usuario autenticado, pantallas de gestión de usuarios y de organización.
\item
  Resultado: matriz de permisos consolidada y aplicada de extremo a extremo.
\end{itemize}

\subsubsection{Sprint 8: Cierre Funcional y Pulido (30 marzo -- 10 abril 2026 \textbar{} 17 SP)}\label{sprint-8-cierre-funcional-y-pulido-30-marzo-10-abril-2026-17-sp}

\begin{itemize}
\tightlist
\item
  Backend: cierre de los últimos endpoints pendientes del backlog, ajustes de filtros y validaciones, mejoras en mensajes de error.
\item
  Frontend: pulido de UI (estados vacíos, loaders, toasts), responsive en vistas críticas, accesibilidad básica en formularios.
\item
  Resultado: alcance funcional de OrbitEngine cerrado y listo para iniciar la fase de estabilización y despliegue.
\end{itemize}

\subsection{Fase 4. Estabilización, Despliegue y Refinamiento (mediados de abril a última semana de abril 2026, sprints 9--10)}\label{fase-4.-estabilizaciuxf3n-despliegue-y-refinamiento-mediados-de-abril-a-uxfaltima-semana-de-abril-2026-sprints-910}

Con el alcance funcional cerrado al final de la Fase 3, esta fase se enfocó en llevar el sistema a producción y dejarlo listo para usuarios externos. Es en este punto cuando se aprovisionó por primera vez la infraestructura de despliegue: hasta ese momento OrbitEngine se ejecutaba únicamente en entornos de desarrollo local con Docker Compose. Por la duración total de la fase, los Sprints 9 y 10 se planificaron como iteraciones cortas de una semana cada una.

\subsubsection{Sprint 9: Pruebas de Carga, Rendimiento y Refinamiento Interno (13--17 abril 2026 \textbar{} 16 SP)}\label{sprint-9-pruebas-de-carga-rendimiento-y-refinamiento-interno-1317-abril-2026-16-sp}

\begin{itemize}
\tightlist
\item
  Diseño y ejecución de pruebas de carga con Locust contra una instancia de pre-producción, alternando entre escenarios de 8 y 200 usuarios concurrentes con distintos perfiles de interacción.
\item
  Pruebas de rendimiento del frontend con Lighthouse y herramientas web gratuitas similares (PageSpeed Insights, WebPageTest) sobre las vistas principales.
\item
  Identificación y corrección de queries lentas, ajustes de índices en PostgreSQL y memoización selectiva de componentes pesados en el frontend.
\item
  Sesiones internas de prueba por parte del equipo simulando flujos reales de una pyme; corrección de defectos de usabilidad detectados.
\item
  Resultado: sistema estabilizado, métricas base de carga y rendimiento documentadas (los resultados se reportan en el Capítulo 5).
\end{itemize}

\subsubsection{Sprint 10: Despliegue en Producción y Validación Interna (20--24 abril 2026 \textbar{} 18 SP)}\label{sprint-10-despliegue-en-producciuxf3n-y-validaciuxf3n-interna-2024-abril-2026-18-sp}

\begin{itemize}
\tightlist
\item
  Aprovisionamiento de la infraestructura productiva: backend y base de datos PostgreSQL en Railway, frontend en Vercel.
\item
  Configuración del dominio definitivo (registrado en Namecheap), certificados TLS gestionados por las plataformas y registros DNS apuntando a Railway y Vercel.
\item
  Configuración del pipeline de despliegue continuo (push a \texttt{main} → despliegue automático en Railway y Vercel) y de las variables de entorno de producción.
\item
  Ampliación de la suite de pruebas Playwright (E2E) y de la cobertura de pruebas backend sobre los flujos críticos.
\item
  Preparación del material de onboarding para las empresas piloto: guías de inicio rápido, plantillas de carga de datos y tutoriales.
\item
  Resultado: OrbitEngine desplegado en producción bajo dominio definitivo, monitorizable y listo para la validación con empresas piloto.
\end{itemize}

\subsection{Fase 5. Validación con Empresas Piloto (27 de abril a 4 de mayo de 2026)}\label{fase-5.-validaciuxf3n-con-empresas-piloto-27-de-abril-a-4-de-mayo-de-2026}

Del 27 de abril al 4 de mayo de 2026, el sistema en producción se puso a disposición de las empresas piloto seleccionadas para su uso real, con sesiones de capacitación, acompañamiento y recolección de retroalimentación. Esta fase se detalla en el Capítulo 6.

\subsection{Fase 6. Documentación Final y Entrega (5 de mayo a 15 de mayo de 2026)}\label{fase-6.-documentaciuxf3n-final-y-entrega-5-de-mayo-a-15-de-mayo-de-2026}

Consolidación del informe de grado, preparación de la presentación y defensa del proyecto del 5 al 15 de mayo de 2026.

\section{Implementación de Módulos Clave}\label{implementaciuxf3n-de-muxf3dulos-clave}

\subsection{Sistema de Autenticación y Multi-Tenancy}\label{sistema-de-autenticaciuxf3n-y-multi-tenancy}

La implementación del sistema de autenticación se diseñó para resolver simultáneamente los requisitos de seguridad y la necesidad de multi-tenancy. El token JWT generado al hacer login contiene:

\begin{Shaded}
\begin{Highlighting}[]
\FunctionTok{\{}
  \DataTypeTok{"sub"}\FunctionTok{:} \StringTok{"\textless{}user\_id\textgreater{}"}\FunctionTok{,}
  \DataTypeTok{"organization\_id"}\FunctionTok{:} \StringTok{"\textless{}organization\_id\textgreater{}"}\FunctionTok{,}
  \DataTypeTok{"role"}\FunctionTok{:} \StringTok{"admin|seller|viewer|contador"}\FunctionTok{,}
  \DataTypeTok{"exp"}\FunctionTok{:} \ErrorTok{\textless{}timestamp\textgreater{}}
\FunctionTok{\}}
\end{Highlighting}
\end{Shaded}

Este diseño permite que el backend derive el contexto de organización directamente del token. La dependencia \texttt{get\_current\_user} de FastAPI valida la firma del token, lo decodifica en un modelo \texttt{TokenPayload} y carga el \texttt{User} correspondiente; sobre ella se construye \texttt{CurrentOrganization}, que expone el \texttt{organization\_id} del usuario autenticado a los handlers:

\begin{Shaded}
\begin{Highlighting}[]
\NormalTok{CurrentUser }\OperatorTok{=}\NormalTok{ Annotated[User, Depends(get\_current\_user)]}
\NormalTok{CurrentOrganization }\OperatorTok{=}\NormalTok{ Annotated[uuid.UUID, Depends(get\_current\_organization)]}

\AttributeTok{@router.get}\NormalTok{(}\StringTok{"/products/"}\NormalTok{, response\_model}\OperatorTok{=}\NormalTok{ProductsPublic)}
\KeywordTok{def}\NormalTok{ list\_products(}
\NormalTok{    session: SessionDep,}
\NormalTok{    \_current\_user: CurrentUser,}
\NormalTok{    current\_organization: CurrentOrganization,}
\NormalTok{) }\OperatorTok{{-}\textgreater{}}\NormalTok{ Any:}
\NormalTok{    products }\OperatorTok{=}\NormalTok{ crud.get\_products(session}\OperatorTok{=}\NormalTok{session, organization\_id}\OperatorTok{=}\NormalTok{current\_organization)}
    \ControlFlowTok{return}\NormalTok{ ProductsPublic(data}\OperatorTok{=}\NormalTok{products, count}\OperatorTok{=}\BuiltInTok{len}\NormalTok{(products))}
\end{Highlighting}
\end{Shaded}

El filtrado por \texttt{organization\_id} en todas las queries de base de datos es el mecanismo fundamental que garantiza el aislamiento de datos entre tenants. Para el control de acceso por rol se usa la dependencia \texttt{require\_role("admin",\ "seller",\ ...)}, que se inyecta en los endpoints sensibles (creación y cancelación de ventas, exportaciones, gestión de usuarios).

\subsection{Módulo de Inventario y Gestión de Stock}\label{muxf3dulo-de-inventario-y-gestiuxf3n-de-stock}

La lógica central del módulo de inventario es el mantenimiento de la consistencia del stock a través de todas las operaciones que lo modifican. Se optó por registrar cada cambio de stock como un \texttt{InventoryMovement} con un \texttt{movement\_type} (\texttt{sale}, \texttt{purchase}, \texttt{adjustment} o \texttt{return}), la cantidad ajustada, los valores de stock previo y nuevo, y una referencia opcional al documento que lo originó (por ejemplo, el \texttt{sale\_id} de la venta).

Esta decisión de diseño proporciona:

\begin{enumerate}
\def\labelenumi{\arabic{enumi}.}
\tightlist
\item
  \textbf{Auditoría completa}: es posible reconstruir el stock en cualquier punto del tiempo reproduciendo los movimientos.
\item
  \textbf{Trazabilidad}: cada reducción de stock puede rastrearse hasta la venta específica o el ajuste manual que la causó.
\item
  \textbf{Confiabilidad}: las actualizaciones de stock y el registro del movimiento se realizan dentro de la misma sesión SQLAlchemy del request, garantizando consistencia ante fallos.
\end{enumerate}

Adicionalmente, el endpoint \texttt{GET\ /products/low-stock} y el contador \texttt{low\_stock\_count} del dashboard exponen los productos cuyo \texttt{stock\_quantity} cae por debajo de su mínimo configurado, lo que permite que la interfaz visualice alertas de reposición sin requerir una tabla de alertas adicional.

\subsection{Módulo de Ventas y Consistencia Transaccional}\label{muxf3dulo-de-ventas-y-consistencia-transaccional}

El registro de una venta involucra múltiples operaciones que deben ejecutarse de forma consistente dentro de la misma sesión de base de datos:

\begin{enumerate}
\def\labelenumi{\arabic{enumi}.}
\tightlist
\item
  Validación de existencia, estado activo y stock disponible para cada producto del carrito.
\item
  Cálculo del subtotal a partir de los precios actuales y aplicación de descuento e impuestos para obtener el total.
\item
  Generación del número de factura secuencial por organización.
\item
  Inserción del registro \texttt{Sale} con los totales calculados y los datos del usuario y cliente asociados.
\item
  Por cada ítem: inserción del \texttt{SaleItem} con un \emph{snapshot} de nombre, SKU y precio del producto; descuento del stock; e inserción del \texttt{InventoryMovement} de tipo \texttt{sale} con \texttt{previous\_stock} y \texttt{new\_stock}.
\item
  Actualización de las métricas de compra del cliente (total comprado, número de compras, ticket promedio) si la venta está asociada a uno.
\end{enumerate}

Todo el flujo se ejecuta dentro de la misma sesión SQLAlchemy gestionada por la dependencia \texttt{SessionDep}. Si cualquiera de las validaciones iniciales falla (stock insuficiente, producto inactivo, cliente inexistente), se levanta una \texttt{HTTPException} antes de cualquier escritura. La operación inversa (\texttt{POST\ /sales/\{sale\_id\}/cancel}) sigue el mismo patrón: restituye el stock, registra movimientos de tipo \texttt{return} y revierte las métricas del cliente.

\section{Estrategia y Resultados de Pruebas}\label{estrategia-y-resultados-de-pruebas}

\subsection{Tipos de Pruebas Implementadas}\label{tipos-de-pruebas-implementadas}

El proyecto implementó cuatro niveles de pruebas automatizadas que en conjunto suman más de 300 casos de prueba ejecutables:

\textbf{Pruebas Unitarias / CRUD (backend):}

\begin{itemize}
\tightlist
\item
  Herramienta: pytest.
\item
  Cobertura: lógica de negocio en \texttt{app/crud.py} para cada entidad (\texttt{category}, \texttt{customer}, \texttt{dashboard}, \texttt{inventory\_movement}, \texttt{organization}, \texttt{product}, \texttt{role}, \texttt{sale}, \texttt{user}).
\item
  Ejemplos: cálculo de totales de venta, generación secuencial del número de factura por organización, validación de stock disponible, actualización de métricas de cliente, agregaciones del dashboard.
\end{itemize}

\textbf{Pruebas de Integración (API):}

\begin{itemize}
\tightlist
\item
  Herramienta: pytest + TestClient de FastAPI (basado en httpx).
\item
  Cobertura: flujos completos de cada router en \texttt{app/api/routes/}: login, organizaciones (signup), usuarios, productos, categorías, clientes, ventas (incluida cancelación), movimientos de inventario, dashboard (\texttt{/stats} y \texttt{/export-excel}) y roles.
\item
  Se utilizaron fixtures (\texttt{conftest.py}) para crear datos de prueba aislados en una base de datos de test separada, garantizando la independencia entre pruebas.
\end{itemize}

\textbf{Pruebas de Seguridad y Multi-Tenancy (backend):}

\begin{itemize}
\tightlist
\item
  Herramienta: pytest sobre el TestClient de FastAPI.
\item
  Cobertura: módulo \texttt{tests/security/test\_multitenant.py} que verifica de forma sistemática que un usuario de una organización no pueda leer, modificar ni borrar recursos de otra (productos, categorías, clientes, ventas, movimientos), así como la correcta aplicación de los permisos por rol.
\end{itemize}

\textbf{Pruebas E2E (frontend):}

\begin{itemize}
\tightlist
\item
  Herramienta: Playwright.
\item
  Cobertura: flujos críticos de usuario incluyendo login, registro de organización (\texttt{organization-signup}), alta y edición de productos (\texttt{inventory}), registro y consulta de ventas (\texttt{sales}), gestión de clientes (\texttt{customers}), gestión de roles (\texttt{roles}), reseteo de contraseña, configuración de cuenta y de organización, y panel de administración.
\end{itemize}

\subsection{Pipeline de CI/CD}\label{pipeline-de-cicd}

Cada push al repositorio de GitHub activa automáticamente el siguiente pipeline:

\begin{verbatim}
Trigger: push a cualquier rama / PR a main
    │
    ├─► Lint & Type Check
    │       Backend: ruff + mypy
    │       Frontend: biome check
    │
    ├─► Unit & Integration Tests
    │       pytest con base de datos de test en memoria (SQLite)
    │       Cobertura mínima: 60% (el pipeline falla si no se alcanza)
    │
    ├─► Build
    │       Backend: construcción de imagen Docker
    │       Frontend: bun run build
    │
    └─► Deploy (solo en merge a main)
            Backend: push a main → deploy automático en Railway
            Frontend: push a main → deploy automático en Vercel
\end{verbatim}

Este pipeline garantiza que solo código probado y con tipos correctos llegue a producción.

\subsection{Métricas de Cobertura}\label{muxe9tricas-de-cobertura}

Con más de 300 pruebas automatizadas distribuidas entre las capas de CRUD, API, multi-tenancy y E2E, la cobertura del backend al cierre de la Fase 3 alcanzó:

\begin{longtable}[]{@{}ll@{}}
\caption[Métricas de Cobertura]{Métricas de Cobertura}\tabularnewline
\toprule\noalign{}
Módulo & Cobertura \\
\midrule\noalign{}
\endhead
\bottomrule\noalign{}
\endlastfoot
Autenticación y usuarios & 92\% \\
Inventario (productos, categorías, movimientos) & 88\% \\
Ventas & 91\% \\
Clientes & 86\% \\
Dashboard y reportes & 82\% \\
Multi-tenancy y permisos & 95\% \\
\textbf{Total} & \textbf{89\%} \\
\end{longtable}

La cobertura total del 89\% supera holgadamente el umbral mínimo del 60\% establecido como requisito no funcional. La buena cobertura del módulo de multi-tenancy es especialmente relevante porque protege la propiedad más crítica del sistema: el aislamiento de datos entre organizaciones.

\subsection{Pruebas de Carga y Rendimiento Planeadas}\label{pruebas-de-carga-y-rendimiento-planeadas}

Para validar el comportamiento del sistema bajo uso real se diseñó un plan de pruebas de carga y rendimiento que se ejecuta durante la Fase 4. Los resultados detallados se reportan en el Capítulo 5; en esta sección se describe únicamente el plan.

\textbf{Pruebas de carga con Locust} (\texttt{backend/tests/performance/locustfile.py}):

\begin{itemize}
\tightlist
\item
  Escenarios alternados entre \textbf{8 usuarios concurrentes} (carga representativa de pyme) y \textbf{200 usuarios concurrentes} (escenario de estrés muy por encima del uso esperado).
\item
  Tres perfiles de usuario virtual mezclados según peso:

  \begin{itemize}
  \tightlist
  \item
    Lector realista (\texttt{OrbitEngineUser}): paginación, búsqueda, filtros y ordenamiento sobre productos, ventas, clientes, movimientos y dashboard.
  \item
    Vendedor (\texttt{SellerUser}): crea ventas reales contra el catálogo y dispara los movimientos de inventario asociados.
  \item
    ``Spammer'' (\texttt{SpammerUser}): martillea sin pausa los endpoints de agregación (\texttt{/dashboard/stats}).
  \end{itemize}
\item
  Cada usuario virtual rota cuentas reales de las distintas organizaciones piloto, lo que ejercita simultáneamente el aislamiento multi-tenant bajo carga.
\item
  Se variará el ritmo de interacciones por segundo (entre \texttt{between(0.5,\ 1.5)} y \texttt{constant(0)} según la clase de usuario) para representar tanto el patrón humano como el peor caso.
\end{itemize}

\textbf{Pruebas de rendimiento del frontend}:

\begin{itemize}
\tightlist
\item
  \textbf{Lighthouse} (Chrome DevTools y CLI) sobre las vistas de login, dashboard, listado de productos, registro de venta y reportes, midiendo Performance, Accessibility, Best Practices y SEO.
\item
  \textbf{PageSpeed Insights} y \textbf{WebPageTest} como herramientas web gratuitas complementarias para obtener mediciones desde redes y ubicaciones distintas a las del equipo de desarrollo.
\item
  Las métricas objetivo (Core Web Vitals: LCP, INP, CLS) y los resultados obtenidos se documentan en el Capítulo 5.
\end{itemize}

\section{Infraestructura de Despliegue}\label{infraestructura-de-despliegue}

\subsection{Arquitectura de Producción}\label{arquitectura-de-producciuxf3n}

El sistema fue desplegado en Railway y Vercel con la siguiente configuración:

\begin{longtable}[]{@{}
  >{\raggedright\arraybackslash}p{(\linewidth - 4\tabcolsep) * \real{0.1196}}
  >{\raggedright\arraybackslash}p{(\linewidth - 4\tabcolsep) * \real{0.1579}}
  >{\raggedright\arraybackslash}p{(\linewidth - 4\tabcolsep) * \real{0.7225}}@{}}
\caption[Arquitectura de Producción]{Arquitectura de Producción}\tabularnewline
\toprule\noalign{}
\begin{minipage}[b]{\linewidth}\raggedright
Plataforma
\end{minipage} & \begin{minipage}[b]{\linewidth}\raggedright
Componente
\end{minipage} & \begin{minipage}[b]{\linewidth}\raggedright
Configuración
\end{minipage} \\
\midrule\noalign{}
\endhead
\bottomrule\noalign{}
\endlastfoot
Railway & Backend API & Servicio web desplegado desde imagen Docker; variables de entorno gestionadas en el panel de Railway \\
Railway & Base de datos PostgreSQL & Instancia gestionada, backups automáticos, URL de conexión inyectada como variable de entorno \\
Vercel & Frontend (React/Vite) & SPA compilada con Vite, CDN global integrado, HTTPS automático y previsualizaciones por pull request \\
Railway / Vercel & Certificados TLS & Certificados gestionados y renovados automáticamente por cada plataforma \\
Namecheap & Registro de dominio y DNS & Dominio adquirido en Namecheap; registros CNAME / A configurados desde el panel de Namecheap para apuntar al backend en Railway y al frontend en Vercel \\
\end{longtable}

\subsection{Estrategia de Secretos y Configuración}\label{estrategia-de-secretos-y-configuraciuxf3n}

La configuración sensible (credenciales de base de datos, secreto JWT) se gestiona mediante las variables de entorno del panel de Railway para el backend, y mediante el panel de Vercel para el frontend. Las variables de configuración no sensibles se definen en el archivo \texttt{docker-compose.yml} para el entorno de desarrollo local.

Esta separación de configuración garantiza que nunca se cometan credenciales al repositorio de código.

\subsection{Monitoreo y Observabilidad}\label{monitoreo-y-observabilidad}

Para esta fase del proyecto se previó apoyarse principalmente en las herramientas nativas de las plataformas de despliegue, dejando integraciones más sofisticadas para etapas posteriores:

\begin{itemize}
\tightlist
\item
  \textbf{Railway Logs}: los logs del backend se centralizan en el panel de Railway, con acceso en tiempo real y retención configurable desde la interfaz web.
\item
  \textbf{Railway Metrics}: métricas de uso de CPU, memoria y tráfico de red disponibles en el panel de Railway, con alertas configurables por umbral.
\item
  \textbf{Vercel Analytics y logs de despliegue}: información de tráfico básico del frontend y trazabilidad de cada despliegue por commit.
\item
  \textbf{Healthcheck endpoint}: \texttt{GET\ /api/v1/utils/health-check/} expuesto por el backend, consultable por Railway o por monitores externos para detectar instancias no saludables.
\end{itemize}

    \chapter{Resultados Técnicos}\label{resultados-tuxe9cnicos}

\begin{quote}
Este capítulo reporta exclusivamente la \textbf{validación técnica} del sistema OrbitEngine: pruebas de carga sobre el backend con Locust y pruebas de rendimiento del frontend con Lighthouse, PageSpeed Insights y WebPageTest. Los \textbf{resultados de la validación con usuarios} (eficiencia operativa, satisfacción, completitud de tareas) se presentan en el Capítulo 6 --- \emph{Resultados de Usuarios}. Las \textbf{acciones de mejora} derivadas de los hallazgos aquí descritos se discuten en el capítulo de Conclusiones.
\end{quote}

\begin{center}\rule{0.5\linewidth}{0.5pt}\end{center}

\section{Marco de la Validación Técnica}\label{marco-de-la-validaciuxf3n-tuxe9cnica}

\subsection{Objetivos}\label{objetivos}

La validación técnica se planteó tres objetivos verificables:

\begin{enumerate}
\def\labelenumi{\arabic{enumi}.}
\tightlist
\item
  \textbf{Caracterizar el comportamiento del backend bajo carga concurrente}, identificando latencias, tasa de fallos y punto de degradación.
\item
  \textbf{Caracterizar el rendimiento percibido del frontend} desde tres puntos de observación independientes: laboratorio local (Lighthouse), infraestructura externa (PageSpeed Insights) y red real con captura de video (WebPageTest).
\item
  \textbf{Contrastar los hallazgos contra los Requisitos No Funcionales} definidos en el Capítulo 3 que aplican a esta validación, con foco en RNF-01 (rendimiento de la API), RNF-02 (disponibilidad), RNF-07 (responsividad de la interfaz) y RNF-09 (escalabilidad arquitectónica).
\end{enumerate}

\subsection{Ambiente de Pruebas}\label{ambiente-de-pruebas}

Todas las mediciones técnicas se realizaron sobre el entorno de producción real del sistema, siguiendo la configuración detallada en el capítulo anterior. El backend, construido en FastAPI, estaba desplegado en la plataforma Railway, ejecutándose como un servicio web sin réplicas horizontales, es decir, con una única instancia principal. Para adecuarse al nivel de exigencia de cada escenario, se configuró el paralelismo interno del servidor ajustando la cantidad de workers: en los tests 01 a 05 se dispuso de dos workers de FastAPI, mientras que para el test 06---que correspondía a un escenario de carga pico sostenida---se aumentó a cuatro workers, permitiendo simular mayor concurrencia en la aplicación.

La base de datos utilizada en todas las pruebas fue PostgreSQL, también gestionada desde Railway. Se trataba de una única instancia compartida, sin réplicas de lectura, lo que significa que todas las operaciones (tanto de lectura como de escritura) recaían sobre el mismo servidor de base de datos.

En cuanto al frontend, este estaba desplegado en Vercel. Se trata de una aplicación de página única (SPA), construida con React y Vite, y servida a través de una red global de distribución de contenido (CDN), lo que permite un acceso rápido y seguro gracias a la integración de HTTPS automático.

Respecto a los dominios accesibles en las pruebas, se utilizaron las rutas oficiales del entorno productivo: \texttt{orbitengine.lat} para el frontend y \texttt{api.orbitengine.lat} para los endpoints del backend. Ambos estaban protegidos con certificados TLS, gestionados y renovados automáticamente por las propias plataformas de despliegue.

Por último, todas las pruebas de carga se ejecutaron desde una única estación de trabajo del equipo, conectada desde una red doméstica situada en Bogotá, Colombia, lo que replica condiciones realistas de acceso de usuarios finales desde una red residencial estándar.

\subsection{Criterios de Aceptación}\label{criterios-de-aceptaciuxf3n}

Los resultados de la validación técnica se evaluaron en función de criterios previamente definidos y recogidos en la Sección 3.2 del Capítulo 3. Estos criterios son:

\begin{itemize}
\tightlist
\item
  Para el rendimiento (RNF-01): Se exige que al menos el 95\% de las respuestas de la API, bajo condiciones de carga normal (hasta 50 usuarios concurrentes), tengan un tiempo de resolución inferior a 500 milisegundos. Este umbral busca garantizar una experiencia ágil incluso cuando la plataforma es utilizada por decenas de usuarios al mismo tiempo.
\item
  En cuanto a la disponibilidad (RNF-02): El sistema debe permanecer disponible y operativo al menos el 95\% del tiempo a lo largo de un mes, limitando al mínimo los periodos de caída o interrupción en el servicio.
\item
  Analizando la usabilidad (RNF-07): Se requiere que la interfaz principal del sistema sea completamente responsive y funcional en dispositivos cuyo ancho de pantalla sea igual o superior a 375 píxeles, que es una referencia común para pantallas de teléfonos móviles actuales.
\item
  Respecto a la escalabilidad (RNF-09): La arquitectura técnica del sistema debe poder soportar el crecimiento en la cantidad de organizaciones (tenants) que lo utilizan, sin necesidad de modificar estructuras internas del software, solamente añadiendo nuevas instancias (escalado horizontal) en la capa de aplicación.
\end{itemize}

Estos criterios delimitan de forma concreta cómo juzgar el éxito o insuficiencia de los resultados obtenidos durante la validación técnica de OrbitEngine.

\subsection{Herramientas Empleadas}\label{herramientas-empleadas}

{\def\LTcaptype{none} % do not increment counter
\begin{longtable}[]{@{}
  >{\raggedright\arraybackslash}p{(\linewidth - 6\tabcolsep) * \real{0.1302}}
  >{\raggedright\arraybackslash}p{(\linewidth - 6\tabcolsep) * \real{0.3964}}
  >{\raggedright\arraybackslash}p{(\linewidth - 6\tabcolsep) * \real{0.1065}}
  >{\raggedright\arraybackslash}p{(\linewidth - 6\tabcolsep) * \real{0.3669}}@{}}
\toprule\noalign{}
\begin{minipage}[b]{\linewidth}\raggedright
Herramienta
\end{minipage} & \begin{minipage}[b]{\linewidth}\raggedright
Versión / canal
\end{minipage} & \begin{minipage}[b]{\linewidth}\raggedright
Capa observada
\end{minipage} & \begin{minipage}[b]{\linewidth}\raggedright
Tipo de medición
\end{minipage} \\
\midrule\noalign{}
\endhead
\bottomrule\noalign{}
\endlastfoot
\textbf{Locust} & 2.x sobre \texttt{tests/performance/locustfile.py} & Backend / API REST & Carga sintética con usuarios virtuales \\
\textbf{Lighthouse} & 13.0.1, ejecutado desde Chrome DevTools & Frontend & Auditoría de laboratorio (escritorio, conexión local) \\
\textbf{PageSpeed Insights} & Web pública de Google (Lighthouse 13.0.1 con Headless Chromium 146) & Frontend & Auditoría desde infraestructura de Google (escritorio y móvil) \\
\textbf{WebPageTest} & versión 26.03 & Frontend & Carga real con captura de waterfall y video \\
\end{longtable}
}

\subsection{Composición del Piloto Técnico}\label{composiciuxf3n-del-piloto-tuxe9cnico}

Las pruebas de este capítulo se ejecutaron sobre un piloto compuesto por \textbf{ocho organizaciones registradas en producción}, con una composición deliberadamente mixta diseñada para combinar realismo de uso con presión sintética sobre el sistema. La siguiente tabla detalla la naturaleza de cada una de las ocho organizaciones.

\textbf{Tabla 5.1.5.} Organizaciones del piloto técnico.

{\def\LTcaptype{none} % do not increment counter
\begin{longtable}[]{@{}
  >{\raggedright\arraybackslash}p{(\linewidth - 6\tabcolsep) * \real{0.0366}}
  >{\raggedright\arraybackslash}p{(\linewidth - 6\tabcolsep) * \real{0.1204}}
  >{\raggedright\arraybackslash}p{(\linewidth - 6\tabcolsep) * \real{0.2984}}
  >{\raggedright\arraybackslash}p{(\linewidth - 6\tabcolsep) * \real{0.5445}}@{}}
\toprule\noalign{}
\begin{minipage}[b]{\linewidth}\raggedright
N°
\end{minipage} & \begin{minipage}[b]{\linewidth}\raggedright
Nombre
\end{minipage} & \begin{minipage}[b]{\linewidth}\raggedright
Naturaleza
\end{minipage} & \begin{minipage}[b]{\linewidth}\raggedright
Sector / propósito
\end{minipage} \\
\midrule\noalign{}
\endhead
\bottomrule\noalign{}
\endlastfoot
1 & \textbf{Frozt Bitez} & \textbf{Empresa real} & Pyme que adoptó OrbitEngine para sus operaciones reales \\
2 & \textbf{Miss Peggy} & \textbf{Empresa real} & Pyme de un sector distinto al de Frozt Bitez, que también adoptó la plataforma para sus operaciones \\
3 & Lehgo & Empresa ficticia de prueba (datos sintéticos) & Generación de volumen de productos, ventas y movimientos para forzar consultas \\
4 & Ferrallas del Norte & Empresa ficticia de prueba (datos sintéticos) & Generación de volumen de productos, ventas y movimientos para forzar consultas \\
5 & Sabor Caribe & Empresa ficticia de prueba (datos sintéticos) & Generación de volumen de productos, ventas y movimientos para forzar consultas \\
6 & Moda Andes & Empresa ficticia de prueba (datos sintéticos) & Generación de volumen de productos, ventas y movimientos para forzar consultas \\
7 & FarmaVida & Empresa ficticia de prueba (datos sintéticos) & Generación de volumen de productos, ventas y movimientos para forzar consultas \\
8 & Default del entorno & Datos de prueba --- primera organización creada como \emph{seed} & Tenant base de pruebas internas; conserva información residual de las primeras iteraciones de desarrollo \\
\end{longtable}
}

\textbf{Resumen cuantitativo.}

{\def\LTcaptype{none} % do not increment counter
\begin{longtable}[]{@{}
  >{\raggedright\arraybackslash}p{(\linewidth - 4\tabcolsep) * \real{0.6735}}
  >{\raggedright\arraybackslash}p{(\linewidth - 4\tabcolsep) * \real{0.0816}}
  >{\raggedright\arraybackslash}p{(\linewidth - 4\tabcolsep) * \real{0.2449}}@{}}
\toprule\noalign{}
\begin{minipage}[b]{\linewidth}\raggedright
Tipo de organización
\end{minipage} & \begin{minipage}[b]{\linewidth}\raggedright
Cantidad
\end{minipage} & \begin{minipage}[b]{\linewidth}\raggedright
Proporción de la muestra
\end{minipage} \\
\midrule\noalign{}
\endhead
\bottomrule\noalign{}
\endlastfoot
\textbf{Empresas reales} que adoptaron OrbitEngine para sus operaciones & 2 & 25 \% \\
\textbf{Empresas ficticias de prueba / datos de prueba} & 6 & 75 \% \\
\textbf{Total} & \textbf{8} & \textbf{100 \%} \\
\end{longtable}
}

\begin{itemize}
\tightlist
\item
  \textbf{Empresas reales (25 \% de la muestra) --- Frozt Bitez y Miss Peggy.} Son las únicas dos pymes que confiaron en la plataforma para realizar parte o la totalidad de sus operaciones cotidianas durante la fase de validación. Aunque en cantidad representan una proporción menor, su valor para el experimento es alto: pertenecen a sectores muy distintos entre sí, manejan trazabilidades diferenciadas de productos, ventas, clientes y movimientos de inventario, y permiten dar un panorama representativo de cómo se comportarán empresas reales que adopten OrbitEngine en el futuro.
\item
  \textbf{Empresas ficticias de prueba (75 \% de la muestra) --- Lehgo, Ferrallas del Norte, Sabor Caribe, Moda Andes, FarmaVida y Default.} Son seis organizaciones creadas por el equipo de desarrollo con el único objetivo de \textbf{poblar el sistema con grandes volúmenes de productos, ventas, clientes, movimientos de inventario y reportes}, de modo que las consultas, agregaciones y filtros del backend trabajen sobre conjuntos de datos suficientemente grandes para forzar respuestas más complejas del servidor. Estas organizaciones \textbf{no representan negocios reales}; en lo que sigue se nombrarán siempre como \emph{empresas ficticias de prueba} o \emph{datos de prueba}. \emph{Default} es, además, la primera organización creada en el entorno y conserva información residual de las pruebas iniciales del equipo, por lo que cumple un papel adicional como \emph{tenant} base.
\end{itemize}

Esta composición es coherente con el alcance de un piloto técnico: \textbf{Frozt Bitez y Miss Peggy} aportan realismo cualitativo sobre el comportamiento productivo del sistema, mientras que \textbf{Lehgo, Ferrallas del Norte, Sabor Caribe, Moda Andes, FarmaVida y Default} aportan el volumen sintético necesario para evidenciar el coste de las consultas en condiciones próximas a las de un sistema en producción a mayor escala.

\section{Pruebas de Carga (Backend / API) --- Locust}\label{pruebas-de-carga-backend-api-locust}

\subsection{Diseño del Experimento}\label{diseuxf1o-del-experimento}

Las pruebas se construyeron con \textbf{Locust} y están versionadas en el repositorio en \texttt{{[}backend/tests/performance/locustfile.py{]}(../../backend/tests/performance/locustfile.py)}. El diseño persigue cuatro propiedades:

\begin{enumerate}
\def\labelenumi{\arabic{enumi}.}
\tightlist
\item
  \textbf{Realismo del tráfico}: combinar exploración (paginación, búsqueda, filtros, ordenamiento) con escritura efectiva (registro de ventas, ajuste de stock).
\item
  \textbf{Multi-tenancy bajo carga}: cada usuario virtual rota credenciales reales de las distintas organizaciones piloto, ejercitando simultáneamente el aislamiento por \texttt{organization\_id} en todas las consultas.
\item
  \textbf{Reproducibilidad}: los escenarios y los pesos relativos de las tareas están codificados en el repositorio y pueden reejecutarse en cualquier momento.
\item
  \textbf{Progresividad}: se ejecutaron seis escenarios encadenados con concurrencia creciente para construir una curva de degradación.
\end{enumerate}

\subsection{Perfiles de Usuario Virtual}\label{perfiles-de-usuario-virtual}

El \texttt{locustfile.py} define tres clases de \texttt{HttpUser} que se mezclan según pesos relativos:

{\def\LTcaptype{none} % do not increment counter
\begin{longtable}[]{@{}
  >{\raggedright\arraybackslash}p{(\linewidth - 4\tabcolsep) * \real{0.1613}}
  >{\raggedright\arraybackslash}p{(\linewidth - 4\tabcolsep) * \real{0.0516}}
  >{\raggedright\arraybackslash}p{(\linewidth - 4\tabcolsep) * \real{0.7871}}@{}}
\toprule\noalign{}
\begin{minipage}[b]{\linewidth}\raggedright
Perfil
\end{minipage} & \begin{minipage}[b]{\linewidth}\raggedright
Peso
\end{minipage} & \begin{minipage}[b]{\linewidth}\raggedright
Comportamiento
\end{minipage} \\
\midrule\noalign{}
\endhead
\bottomrule\noalign{}
\endlastfoot
OrbitEngineUser & 3 & Lee dashboard, productos, ventas, clientes y categorías. Ejercita paginación extrema, búsqueda, ordenamiento y filtros. \\
SellerUser & 2 & Crea ventas reales contra el catálogo, ajusta stock y consulta movimientos. \\
SpammerUser & 1 & Martillea sin pausa los endpoints de agregación y peticiones a UUIDs inexistentes para ejercitar el camino de error 404. \\
\end{longtable}
}

La rotación de cuentas se hace de forma cíclica y \emph{thread-safe} sobre la lista \texttt{ACCOUNTS}, que contiene credenciales válidas de las ocho organizaciones del piloto técnico descritas en la sección 5.1.5: las \textbf{dos empresas reales} (\textbf{Frozt Bitez} y \textbf{Miss Peggy}) y las \textbf{seis empresas ficticias de prueba} (\textbf{Lehgo}, \textbf{Ferrallas del Norte}, \textbf{Sabor Caribe}, \textbf{Moda Andes}, \textbf{FarmaVida} y \textbf{Default}). Esta rotación cumple dos funciones en paralelo: (i) cada usuario virtual ejercita el aislamiento por \texttt{organization\_id} desde un \emph{tenant} distinto, validando la integridad multi-tenant bajo carga; (ii) las peticiones consultan tanto los volúmenes reales y acotados de Frozt Bitez y Miss Peggy como los volúmenes grandes generados por las seis empresas ficticias de prueba, lo que aporta un perfil de coste por consulta más representativo que el que obtendría un único \emph{tenant}.

\subsection{Escenarios Ejecutados}\label{escenarios-ejecutados}

Se ejecutaron seis escenarios consecutivos contra \texttt{https://api.orbitengine.lat}. La configuración exacta (\texttt{-\/-users} y \texttt{-\/-spawn-rate}) se eligió siguiendo los tres regímenes documentados en el \texttt{docstring} del \texttt{locustfile}:

\begin{itemize}
\tightlist
\item
  \textbf{Carga normal de pyme} (\textasciitilde8 usuarios)
\item
  \textbf{Estrés} (\textasciitilde50 usuarios)
\item
  \textbf{Pico} (\textasciitilde200 usuarios, ajustando la duración para acotar el daño en los regímenes de saturación).
\end{itemize}

{\def\LTcaptype{none} % do not increment counter
\begin{longtable}[]{@{}
  >{\raggedright\arraybackslash}p{(\linewidth - 12\tabcolsep) * \real{0.0918}}
  >{\raggedright\arraybackslash}p{(\linewidth - 12\tabcolsep) * \real{0.3980}}
  >{\raggedright\arraybackslash}p{(\linewidth - 12\tabcolsep) * \real{0.1122}}
  >{\raggedright\arraybackslash}p{(\linewidth - 12\tabcolsep) * \real{0.1020}}
  >{\raggedright\arraybackslash}p{(\linewidth - 12\tabcolsep) * \real{0.0612}}
  >{\raggedright\arraybackslash}p{(\linewidth - 12\tabcolsep) * \real{0.1429}}
  >{\raggedright\arraybackslash}p{(\linewidth - 12\tabcolsep) * \real{0.0918}}@{}}
\toprule\noalign{}
\begin{minipage}[b]{\linewidth}\raggedright
Escenario
\end{minipage} & \begin{minipage}[b]{\linewidth}\raggedright
Régimen
\end{minipage} & \begin{minipage}[b]{\linewidth}\raggedright
Duración
\end{minipage} & \begin{minipage}[b]{\linewidth}\raggedright
Total req.
\end{minipage} & \begin{minipage}[b]{\linewidth}\raggedright
Fallos
\end{minipage} & \begin{minipage}[b]{\linewidth}\raggedright
Tasa de fallos
\end{minipage} & \begin{minipage}[b]{\linewidth}\raggedright
RPS prom.
\end{minipage} \\
\midrule\noalign{}
\endhead
\bottomrule\noalign{}
\endlastfoot
Test 01 & Carga normal (\textasciitilde8 usuarios concurrentes) & 1 min 39 s & 235 & 0 & \textbf{0,0 \%} & 2,38 \\
Test 02 & Estrés moderado (\textasciitilde50 usuarios) & 3 min 02 s & 2 973 & 51 & \textbf{1,7 \%} & 16,34 \\
Test 03 & Estrés intermedio & 2 min 18 s & 1 600 & 103 & \textbf{6,4 \%} & 11,58 \\
Test 04 & Saturación & 1 min 10 s & 360 & 41 & \textbf{11,4 \%} & 6,97 \\
Test 05 & Saturación severa & 51 s & 358 & 79 & \textbf{22,1 \%} & 5,07 \\
Test 06 & Pico sostenido (\textasciitilde200 usuarios) & 10 min 05 s & 7 331 & 1 408 & \textbf{19,2 \%} & 12,12 \\
\end{longtable}
}

\begin{quote}
Los CSV completos de cada escenario (\texttt{Test0X-Requests.csv}, \texttt{Test0X-Failures.csv} y, cuando aplica, \texttt{Test0X-Exception.csv}), así como los reportes HTML generados por Locust, no se incluyen en este repositorio público; para acceder a los soportes correspondientes, consultar a los desarrolladores del proyecto.
\end{quote}

\subsection{Resultados por Escenario}\label{resultados-por-escenario}

\subsubsection{Test 01 --- Carga Normal de Pyme}\label{test-01-carga-normal-de-pyme}

Es el escenario de referencia y representa el régimen al que se diseñó el sistema. Con aproximadamente ocho usuarios virtuales se completaron \textbf{235 peticiones sin fallos}.

{\def\LTcaptype{none} % do not increment counter
\begin{longtable}[]{@{}
  >{\raggedright\arraybackslash}p{(\linewidth - 12\tabcolsep) * \real{0.2625}}
  >{\raggedright\arraybackslash}p{(\linewidth - 12\tabcolsep) * \real{0.0750}}
  >{\raggedright\arraybackslash}p{(\linewidth - 12\tabcolsep) * \real{0.0625}}
  >{\raggedright\arraybackslash}p{(\linewidth - 12\tabcolsep) * \real{0.1500}}
  >{\raggedright\arraybackslash}p{(\linewidth - 12\tabcolsep) * \real{0.1500}}
  >{\raggedright\arraybackslash}p{(\linewidth - 12\tabcolsep) * \real{0.1500}}
  >{\raggedright\arraybackslash}p{(\linewidth - 12\tabcolsep) * \real{0.1500}}@{}}
\toprule\noalign{}
\begin{minipage}[b]{\linewidth}\raggedright
Endpoint
\end{minipage} & \begin{minipage}[b]{\linewidth}\raggedright
Método
\end{minipage} & \begin{minipage}[b]{\linewidth}\raggedright
Reqs.
\end{minipage} & \begin{minipage}[b]{\linewidth}\raggedright
Mediana
\end{minipage} & \begin{minipage}[b]{\linewidth}\raggedright
Promedio
\end{minipage} & \begin{minipage}[b]{\linewidth}\raggedright
P95
\end{minipage} & \begin{minipage}[b]{\linewidth}\raggedright
P99
\end{minipage} \\
\midrule\noalign{}
\endhead
\bottomrule\noalign{}
\endlastfoot
/login/access-token & POST & 8 & 910 ms & 1 196 ms & 2 100 ms & 2 100 ms \\
/categories/ & GET & 20 & 430 ms & 454 ms & 850 ms & 850 ms \\
/customers/ & GET & 41 & 430 ms & 453 ms & 480 ms & 850 ms \\
/dashboard/stats & GET & 57 & 710 ms & 721 ms & 780 ms & 1 100 ms \\
/inventory-movements/ & GET & 20 & 520 ms & 542 ms & 640 ms & 640 ms \\
/products/ & GET & 45 & 500 ms & 515 ms & 710 ms & 920 ms \\
/products/low-stock & GET & 21 & 430 ms & 472 ms & 840 ms & 850 ms \\
/sales/ & GET & 23 & \textbf{7 500 ms} & \textbf{7 501 ms} & \textbf{7 600 ms} & \textbf{7 600 ms} \\
\textbf{Agregado} & --- & 235 & 520 ms & 1 254 ms & 7 500 ms & 7 600 ms \\
\end{longtable}
}

Observación clave: la mediana global (520 ms) y la mediana de los endpoints CRUD se sitúan en el orden de cientos de milisegundos, pero el endpoint \texttt{GET\ /sales/} consume sistemáticamente \textbf{alrededor de 7,5 segundos por petición} incluso sin concurrencia significativa. Este endpoint domina el percentil 95 global y es responsable de que el agregado se aleje de los 500 ms exigidos por el RNF-01.

\subsubsection{Test 02 --- Estrés Moderado (\textasciitilde50 usuarios)}\label{test-02-estruxe9s-moderado-50-usuarios}

Con \textasciitilde50 usuarios virtuales sostenidos durante poco más de tres minutos, el sistema procesa \textbf{2 973 peticiones} con una tasa de fallos del 1,7 \%. Los fallos se concentran en dos categorías:

\begin{itemize}
\tightlist
\item
  45 respuestas 404 esperadas en \texttt{GET\ /\{resource\}/\{bad-uuid\}}, que el \texttt{SpammerUser} genera intencionalmente para ejercitar el camino de error.
\item
  6 errores 500 en \texttt{POST\ /sales/}, atribuibles a contenciones puntuales sobre el stock al ejecutar concurrentemente la creación de ventas y los ajustes de inventario.
\end{itemize}

El sistema sostiene \textbf{16,3 RPS} con la mediana global en 1 000 ms y el percentil 95 alrededor de 7,8 s, todavía dominado por \texttt{GET\ /sales/} y por \texttt{GET\ /sales/?limit=500\ {[}spam{]}}, cuyo P95 alcanza 30 s.

\subsubsection{Tests 03 a 05 --- Régimen de Saturación}\label{tests-03-a-05-ruxe9gimen-de-saturaciuxf3n}

Los escenarios 03, 04 y 05 reflejan la transición desde el estrés moderado hacia la saturación. Las tasas de fallos crecen del 6,4 \% al 22,1 \% y los percentiles superiores se desplazan al rango de las decenas de segundos. Los 500 dejan de concentrarse en \texttt{POST\ /sales/} y aparecen también en \texttt{GET\ /dashboard/stats}, \texttt{GET\ /products/?search=}, \texttt{GET\ /customers/?search=} y \texttt{GET\ /products/\{id\}/movements}. La duración total de cada test se acortó deliberadamente para no comprometer la disponibilidad del sistema durante períodos prolongados.

\subsubsection{Test 06 --- Pico Sostenido (\textasciitilde200 usuarios)}\label{test-06-pico-sostenido-200-usuarios}

Es el escenario más extenso (10 min 05 s) y el más exigente. Para sostener el régimen de pico se incrementó la concurrencia interna del backend de 2 a \textbf{4 workers de FastAPI} sobre la misma instancia de Railway. Sobre \textbf{7 331 peticiones} se registran \textbf{1 408 fallos} (19,2 \%), todos clasificados como 500 (errores del lado del servidor) y distribuidos sobre los endpoints de lectura masiva: \texttt{GET\ /dashboard/stats} (323), \texttt{GET\ /products/} (314), \texttt{GET\ /sales/} (238), \texttt{GET\ /customers/} (191), \texttt{GET\ /categories/} (120), \texttt{GET\ /inventory-movements/} (117) y \texttt{GET\ /products/low-stock} (105). El sistema sigue procesando 12 RPS pero la mediana global se sitúa en 1 300 ms y el P95 en 29 s, evidenciando que la configuración de 4 workers, aun siendo el doble de la utilizada en los escenarios previos, no logra absorber un régimen de \textasciitilde200 usuarios concurrentes sostenido durante diez minutos.

\subsection{Curva de Degradación}\label{curva-de-degradaciuxf3n}

Consolidando los seis escenarios se obtiene la siguiente curva de comportamiento:

{\def\LTcaptype{none} % do not increment counter
\begin{longtable}[]{@{}
  >{\raggedright\arraybackslash}p{(\linewidth - 12\tabcolsep) * \real{0.2353}}
  >{\raggedright\arraybackslash}p{(\linewidth - 12\tabcolsep) * \real{0.1176}}
  >{\raggedright\arraybackslash}p{(\linewidth - 12\tabcolsep) * \real{0.1176}}
  >{\raggedright\arraybackslash}p{(\linewidth - 12\tabcolsep) * \real{0.1324}}
  >{\raggedright\arraybackslash}p{(\linewidth - 12\tabcolsep) * \real{0.1324}}
  >{\raggedright\arraybackslash}p{(\linewidth - 12\tabcolsep) * \real{0.1324}}
  >{\raggedright\arraybackslash}p{(\linewidth - 12\tabcolsep) * \real{0.1324}}@{}}
\toprule\noalign{}
\begin{minipage}[b]{\linewidth}\raggedright
Métrica
\end{minipage} & \begin{minipage}[b]{\linewidth}\raggedright
Test 01
\end{minipage} & \begin{minipage}[b]{\linewidth}\raggedright
Test 02
\end{minipage} & \begin{minipage}[b]{\linewidth}\raggedright
Test 03
\end{minipage} & \begin{minipage}[b]{\linewidth}\raggedright
Test 04
\end{minipage} & \begin{minipage}[b]{\linewidth}\raggedright
Test 05
\end{minipage} & \begin{minipage}[b]{\linewidth}\raggedright
Test 06
\end{minipage} \\
\midrule\noalign{}
\endhead
\bottomrule\noalign{}
\endlastfoot
Mediana agregada & 520 ms & 1 000 ms & 2 000 ms & 2 300 ms & 2 300 ms & 1 300 ms \\
P95 agregado & 7 500 ms & 7 800 ms & 36 000 ms & 33 000 ms & 58 000 ms & 29 000 ms \\
Tasa de fallos & 0,0 \% & 1,7 \% & 6,4 \% & 11,4 \% & 22,1 \% & 19,2 \% \\
RPS efectivo & 2,4 & 16,3 & 11,6 & 7,0 & 5,1 & 12,1 \\
\end{longtable}
}

La tasa de fallos crece de manera monótona hasta el Test 05 (régimen de pico no sostenido, ejecutado todavía con 2 workers) y luego se estabiliza alrededor del 19 \% en el Test 06, ya con 4 workers, donde la duración prolongada permite caracterizar mejor el techo del sistema. El RPS efectivo es máximo bajo estrés moderado (Test 02) y disminuye en los regímenes de saturación porque la mayor parte de los hilos de los workers quedan ocupados respondiendo peticiones lentas, lo que confirma un \emph{throughput collapse} clásico cuando se rebasa la capacidad nominal de la configuración desplegada.

\subsection{Comportamiento por Endpoint}\label{comportamiento-por-endpoint}

Los CSV permiten aislar el comportamiento de los endpoints más representativos a lo largo de los seis escenarios. A continuación se presenta cada endpoint con su tabla de resultados y la observación correspondiente.

\textbf{\texttt{POST\ /login/access-token}}

{\def\LTcaptype{none} % do not increment counter
\begin{longtable}[]{@{}
  >{\raggedright\arraybackslash}p{(\linewidth - 12\tabcolsep) * \real{0.1607}}
  >{\raggedright\arraybackslash}p{(\linewidth - 12\tabcolsep) * \real{0.1250}}
  >{\raggedright\arraybackslash}p{(\linewidth - 12\tabcolsep) * \real{0.1429}}
  >{\raggedright\arraybackslash}p{(\linewidth - 12\tabcolsep) * \real{0.1429}}
  >{\raggedright\arraybackslash}p{(\linewidth - 12\tabcolsep) * \real{0.1429}}
  >{\raggedright\arraybackslash}p{(\linewidth - 12\tabcolsep) * \real{0.1429}}
  >{\raggedright\arraybackslash}p{(\linewidth - 12\tabcolsep) * \real{0.1429}}@{}}
\toprule\noalign{}
\begin{minipage}[b]{\linewidth}\raggedright
Escenario
\end{minipage} & \begin{minipage}[b]{\linewidth}\raggedright
Test 01
\end{minipage} & \begin{minipage}[b]{\linewidth}\raggedright
Test 02
\end{minipage} & \begin{minipage}[b]{\linewidth}\raggedright
Test 03
\end{minipage} & \begin{minipage}[b]{\linewidth}\raggedright
Test 04
\end{minipage} & \begin{minipage}[b]{\linewidth}\raggedright
Test 05
\end{minipage} & \begin{minipage}[b]{\linewidth}\raggedright
Test 06
\end{minipage} \\
\midrule\noalign{}
\endhead
\bottomrule\noalign{}
\endlastfoot
Mediana & 910 ms & 1 200 ms & 2 000 ms & 2 200 ms & 2 200 ms & - \\
\end{longtable}
}

La mediana de este endpoint crece de forma progresiva desde 910 ms en carga normal hasta alcanzar una meseta en 2 200 ms durante los escenarios de saturación (Tests 04 y 05), lo que representa un incremento del 142 \% respecto al valor de referencia. El hecho de que la mediana se estabilice en ese valor ---en lugar de continuar creciendo--- sugiere que el sistema de cola de FastAPI comienza a serializar las solicitudes de autenticación antes de procesarlas, limitando la degradación adicional. Este comportamiento es coherente con la naturaleza del flujo: el coste dominante no proviene de la base de datos sino del hashing bcrypt aplicado a la contraseña, una operación deliberadamente costosa por diseño de seguridad. La ausencia de datos para Test 06 indica que en ese escenario el endpoint no acumuló suficientes muestras representativas, posiblemente porque los usuarios virtuales ya disponían de sesión activa al inicio de la ronda de medición.

\textbf{\texttt{GET\ /products/}}

{\def\LTcaptype{none} % do not increment counter
\begin{longtable}[]{@{}
  >{\raggedright\arraybackslash}p{(\linewidth - 12\tabcolsep) * \real{0.1667}}
  >{\raggedright\arraybackslash}p{(\linewidth - 12\tabcolsep) * \real{0.1296}}
  >{\raggedright\arraybackslash}p{(\linewidth - 12\tabcolsep) * \real{0.1296}}
  >{\raggedright\arraybackslash}p{(\linewidth - 12\tabcolsep) * \real{0.1481}}
  >{\raggedright\arraybackslash}p{(\linewidth - 12\tabcolsep) * \real{0.1481}}
  >{\raggedright\arraybackslash}p{(\linewidth - 12\tabcolsep) * \real{0.1481}}
  >{\raggedright\arraybackslash}p{(\linewidth - 12\tabcolsep) * \real{0.1296}}@{}}
\toprule\noalign{}
\begin{minipage}[b]{\linewidth}\raggedright
Escenario
\end{minipage} & \begin{minipage}[b]{\linewidth}\raggedright
Test 01
\end{minipage} & \begin{minipage}[b]{\linewidth}\raggedright
Test 02
\end{minipage} & \begin{minipage}[b]{\linewidth}\raggedright
Test 03
\end{minipage} & \begin{minipage}[b]{\linewidth}\raggedright
Test 04
\end{minipage} & \begin{minipage}[b]{\linewidth}\raggedright
Test 05
\end{minipage} & \begin{minipage}[b]{\linewidth}\raggedright
Test 06
\end{minipage} \\
\midrule\noalign{}
\endhead
\bottomrule\noalign{}
\endlastfoot
Mediana & 500 ms & 730 ms & 1 100 ms & 4 000 ms & 4 000 ms & 690 ms \\
\end{longtable}
}

Este endpoint evidencia dos fases de comportamiento claramente diferenciadas. En los regímenes de carga normal y estrés moderado (Tests 01 y 02), la mediana se mantiene controlada entre 500 ms y 730 ms, lo que refleja que la paginación del lado del servidor aísla eficazmente el coste por petición del volumen total de registros. Sin embargo, al entrar en el régimen de saturación (Tests 03 y 04), la mediana escala bruscamente hasta 4 000 ms, evidenciando que cuando la cola de workers se satura, las peticiones de lectura paginadas quedan represadas detrás de operaciones de mayor coste. El descenso a 690 ms en Test 06, con cuatro workers activos, es particularmente significativo: a pesar de que la concurrencia es mayor que en Tests 04 y 05, la duplicación de workers absorbe la contención y restaura la latencia a un nivel próximo al de estrés moderado, lo que confirma que la limitación en esos escenarios intermedios era de capacidad de procesamiento y no de diseño de la consulta.

\textbf{\texttt{GET\ /customers/}}

{\def\LTcaptype{none} % do not increment counter
\begin{longtable}[]{@{}
  >{\raggedright\arraybackslash}p{(\linewidth - 12\tabcolsep) * \real{0.1636}}
  >{\raggedright\arraybackslash}p{(\linewidth - 12\tabcolsep) * \real{0.1273}}
  >{\raggedright\arraybackslash}p{(\linewidth - 12\tabcolsep) * \real{0.1273}}
  >{\raggedright\arraybackslash}p{(\linewidth - 12\tabcolsep) * \real{0.1455}}
  >{\raggedright\arraybackslash}p{(\linewidth - 12\tabcolsep) * \real{0.1636}}
  >{\raggedright\arraybackslash}p{(\linewidth - 12\tabcolsep) * \real{0.1455}}
  >{\raggedright\arraybackslash}p{(\linewidth - 12\tabcolsep) * \real{0.1273}}@{}}
\toprule\noalign{}
\begin{minipage}[b]{\linewidth}\raggedright
Escenario
\end{minipage} & \begin{minipage}[b]{\linewidth}\raggedright
Test 01
\end{minipage} & \begin{minipage}[b]{\linewidth}\raggedright
Test 02
\end{minipage} & \begin{minipage}[b]{\linewidth}\raggedright
Test 03
\end{minipage} & \begin{minipage}[b]{\linewidth}\raggedright
Test 04
\end{minipage} & \begin{minipage}[b]{\linewidth}\raggedright
Test 05
\end{minipage} & \begin{minipage}[b]{\linewidth}\raggedright
Test 06
\end{minipage} \\
\midrule\noalign{}
\endhead
\bottomrule\noalign{}
\endlastfoot
Mediana & 430 ms & 540 ms & 1 500 ms & 31 000 ms & 1 200 ms & 620 ms \\
\end{longtable}
}

Este endpoint sigue en términos generales el patrón de degradación de \texttt{GET\ /products/}, pero con una anomalía puntual de mayor magnitud en el escenario de saturación máxima: la mediana de Test 04 asciende a 31 000 ms, un valor que supera en más de veinte veces el registrado en Test 03 (1 500 ms) y que no tiene correlato proporcional en ningún otro endpoint de listado. Este pico se interpreta como el resultado de una cola de espera extrema sobre el pool de conexiones de la base de datos en un instante de saturación simultánea con las escrituras de \texttt{POST\ /sales/}, que en Test 04 también registran su valor máximo. La recuperación a 1 200 ms en Test 05 y a 620 ms en Test 06 confirma que el episodio fue transitorio y dependiente de la combinación particular de carga en ese escenario, y no de una degradación estructural del endpoint.

\textbf{\texttt{GET\ /dashboard/stats}}

{\def\LTcaptype{none} % do not increment counter
\begin{longtable}[]{@{}
  >{\raggedright\arraybackslash}p{(\linewidth - 12\tabcolsep) * \real{0.1385}}
  >{\raggedright\arraybackslash}p{(\linewidth - 12\tabcolsep) * \real{0.1077}}
  >{\raggedright\arraybackslash}p{(\linewidth - 12\tabcolsep) * \real{0.1077}}
  >{\raggedright\arraybackslash}p{(\linewidth - 12\tabcolsep) * \real{0.1231}}
  >{\raggedright\arraybackslash}p{(\linewidth - 12\tabcolsep) * \real{0.1231}}
  >{\raggedright\arraybackslash}p{(\linewidth - 12\tabcolsep) * \real{0.1231}}
  >{\raggedright\arraybackslash}p{(\linewidth - 12\tabcolsep) * \real{0.2769}}@{}}
\toprule\noalign{}
\begin{minipage}[b]{\linewidth}\raggedright
Escenario
\end{minipage} & \begin{minipage}[b]{\linewidth}\raggedright
Test 01
\end{minipage} & \begin{minipage}[b]{\linewidth}\raggedright
Test 02
\end{minipage} & \begin{minipage}[b]{\linewidth}\raggedright
Test 03
\end{minipage} & \begin{minipage}[b]{\linewidth}\raggedright
Test 04
\end{minipage} & \begin{minipage}[b]{\linewidth}\raggedright
Test 05
\end{minipage} & \begin{minipage}[b]{\linewidth}\raggedright
Test 06
\end{minipage} \\
\midrule\noalign{}
\endhead
\bottomrule\noalign{}
\endlastfoot
Mediana & 710 ms & 850 ms & 1 700 ms & 4 800 ms & 1 700 ms & 1 100 ms (19 \% 500) \\
\end{longtable}
}

Este endpoint concentra las consultas de agregación más costosas del sistema ---totales de ventas, unidades vendidas por período y niveles de stock--- ejecutadas directamente contra PostgreSQL sin capa de caché intermedia. La progresión controlada de 710 ms a 850 ms entre los Tests 01 y 02 indica que la base de datos absorbe el incremento inicial de concurrencia sin saturarse. A partir del régimen de saturación, sin embargo, la mediana escala hasta 1 700 ms en Test 03 y alcanza su pico en 4 800 ms en Test 04, para luego retroceder a 1 700 ms en Test 05 y a 1 100 ms en Test 06. Este comportamiento no lineal ---en el que Test 06 presenta la segunda mediana más baja pese a ser el escenario con mayor concurrencia--- se explica en parte por la incorporación de cuatro workers en ese escenario, pero también por el hecho de que las respuestas fallidas (19 \% de errores 500) cortan el ciclo de procesamiento antes de que las agregaciones completen, lo que artificialmente reduce la mediana de las peticiones que sí responden. Este endpoint constituye el candidato prioritario para la incorporación de una capa de caché distribuida en iteraciones futuras del sistema.

\textbf{\texttt{GET\ /sales/}}

{\def\LTcaptype{none} % do not increment counter
\begin{longtable}[]{@{}
  >{\raggedright\arraybackslash}p{(\linewidth - 12\tabcolsep) * \real{0.1286}}
  >{\raggedright\arraybackslash}p{(\linewidth - 12\tabcolsep) * \real{0.1714}}
  >{\raggedright\arraybackslash}p{(\linewidth - 12\tabcolsep) * \real{0.1714}}
  >{\raggedright\arraybackslash}p{(\linewidth - 12\tabcolsep) * \real{0.1143}}
  >{\raggedright\arraybackslash}p{(\linewidth - 12\tabcolsep) * \real{0.1143}}
  >{\raggedright\arraybackslash}p{(\linewidth - 12\tabcolsep) * \real{0.1286}}
  >{\raggedright\arraybackslash}p{(\linewidth - 12\tabcolsep) * \real{0.1714}}@{}}
\toprule\noalign{}
\begin{minipage}[b]{\linewidth}\raggedright
Escenario
\end{minipage} & \begin{minipage}[b]{\linewidth}\raggedright
Test 01
\end{minipage} & \begin{minipage}[b]{\linewidth}\raggedright
Test 02
\end{minipage} & \begin{minipage}[b]{\linewidth}\raggedright
Test 03
\end{minipage} & \begin{minipage}[b]{\linewidth}\raggedright
Test 04
\end{minipage} & \begin{minipage}[b]{\linewidth}\raggedright
Test 05
\end{minipage} & \begin{minipage}[b]{\linewidth}\raggedright
Test 06
\end{minipage} \\
\midrule\noalign{}
\endhead
\bottomrule\noalign{}
\endlastfoot
Mediana & \textbf{7 500 ms} & \textbf{7 700 ms} & 8 500 ms & 8 900 ms & 11 000 ms & \textbf{8 600 ms} \\
\end{longtable}
}

Este endpoint constituye el principal cuello de botella del sistema en todos los regímenes evaluados, y su perfil de latencia es el más revelador del conjunto. La mediana de 7 500 ms registrada en carga normal ---con apenas ocho usuarios concurrentes--- demuestra que el problema es estructural e independiente de la presión competitiva: la consulta recorre la relación \texttt{Sale\ →\ SaleItem} mediante joins y ejecuta agregaciones por venta sobre el conjunto completo de registros, sin que el servidor aplique paginación optimizada que limite el volumen de trabajo por petición. La progresión entre los seis escenarios ---7 500 ms, 7 700 ms, 8 500 ms, 8 900 ms, 11 000 ms y 8 600 ms--- muestra un incremento sostenido pero relativamente contenido: el valor máximo (Test 05, 11 000 ms) solo es un 47 \% superior al de referencia (Test 01, 7 500 ms), lo que contrasta con las variaciones de varios órdenes de magnitud que exhiben otros endpoints en saturación. Esta compresión del rango de degradación confirma que el factor dominante es el coste intrínseco de la consulta, no la contención de recursos, y que la solución debe buscarse en la optimización de la capa de datos ---paginación del lado del servidor, desnormalización selectiva o materialización de vistas--- y no en el escalado de la infraestructura.

\textbf{\texttt{POST\ /sales/}}

{\def\LTcaptype{none} % do not increment counter
\begin{longtable}[]{@{}
  >{\raggedright\arraybackslash}p{(\linewidth - 12\tabcolsep) * \real{0.1607}}
  >{\raggedright\arraybackslash}p{(\linewidth - 12\tabcolsep) * \real{0.1250}}
  >{\raggedright\arraybackslash}p{(\linewidth - 12\tabcolsep) * \real{0.1429}}
  >{\raggedright\arraybackslash}p{(\linewidth - 12\tabcolsep) * \real{0.1429}}
  >{\raggedright\arraybackslash}p{(\linewidth - 12\tabcolsep) * \real{0.1607}}
  >{\raggedright\arraybackslash}p{(\linewidth - 12\tabcolsep) * \real{0.1429}}
  >{\raggedright\arraybackslash}p{(\linewidth - 12\tabcolsep) * \real{0.1250}}@{}}
\toprule\noalign{}
\begin{minipage}[b]{\linewidth}\raggedright
Escenario
\end{minipage} & \begin{minipage}[b]{\linewidth}\raggedright
Test 01
\end{minipage} & \begin{minipage}[b]{\linewidth}\raggedright
Test 02
\end{minipage} & \begin{minipage}[b]{\linewidth}\raggedright
Test 03
\end{minipage} & \begin{minipage}[b]{\linewidth}\raggedright
Test 04
\end{minipage} & \begin{minipage}[b]{\linewidth}\raggedright
Test 05
\end{minipage} & \begin{minipage}[b]{\linewidth}\raggedright
Test 06
\end{minipage} \\
\midrule\noalign{}
\endhead
\bottomrule\noalign{}
\endlastfoot
Mediana & --- & 2 500 ms & 3 300 ms & 31 000 ms & 6 800 ms & --- \\
\end{longtable}
}

Este endpoint presenta la trayectoria de degradación más pronunciada del conjunto, así como el valor pico más alto registrado en toda la batería de pruebas. Bajo estrés moderado (Test 02) la mediana se sitúa en 2 500 ms, lo que refleja el coste base de una operación transaccional que agrupa tres escrituras atómicas en la misma sesión SQLAlchemy: la creación del registro \texttt{Sale}, la inserción de los correspondientes \texttt{SaleItem} y el ajuste del \texttt{InventoryMovement}. Al entrar en saturación, la mediana escala a 3 300 ms en Test 03 y colapsa a 31 000 ms en Test 04, el valor más elevado registrado en toda la batería. Este episodio de saturación extrema es el resultado directo de la contención sobre los bloqueos de fila que protegen el stock disponible: cuando múltiples transacciones concurrentes intentan actualizar el inventario de los mismos productos, las colas de espera sobre esos bloqueos se acumulan de forma no lineal. La recuperación a 6 800 ms en Test 05 sugiere que la reducción en la duración total de ese escenario y el consecuente menor solapamiento de transacciones aliviaron parcialmente la contención. La ausencia de datos en Test 01 y Test 06 obedece a que el volumen de peticiones en esos escenarios no generó muestras estadísticamente representativas para esta operación de escritura.

\subsection{Cumplimiento del RNF-01}\label{cumplimiento-del-rnf-01}

El RNF-01 exige que el 95 \% de las respuestas se complete en menos de 500 ms bajo carga normal (hasta 50 usuarios concurrentes). El balance, leído a la luz del entorno de despliegue dimensionado para el alcance del proyecto, es el siguiente:

\begin{itemize}
\tightlist
\item
  En \textbf{carga normal} (Test 01, \textasciitilde8 usuarios concurrentes), que corresponde al perfil real al que apunta OrbitEngine para una pyme tipo, las medianas de los endpoints CRUD individuales se mantienen entre 430 ms y 720 ms, con el 95 \% de las peticiones por debajo del segundo para todos los endpoints transaccionales del día a día (consultar productos, clientes, dashboard, registrar ventas).
\item
  En \textbf{estrés moderado} (Test 02, \textasciitilde50 usuarios), que ya cubre el límite superior fijado por el RNF-01, el sistema sigue procesando 16,3 RPS con una tasa de fallos del 1,7 \% y manteniendo los endpoints CRUD en el rango de 540 ms a 1 200 ms, sin colapsos en ningún flujo crítico.
\item
  El único endpoint que se aleja del umbral de manera estructural es \texttt{GET\ /sales/}, cuyo coste alto es independiente de la concurrencia y se origina en la composición del \emph{payload} (joins con \texttt{SaleItem} y agregaciones por venta) más que en una limitación de la infraestructura. Este comportamiento queda registrado como punto de optimización futura.
\end{itemize}

Por lo tanto, \textbf{el RNF-01 se considera cumplido a cabalidad para el entorno de desarrollo y despliegue dimensionado para esta aplicación}, donde el perfil de uso esperado (≤ 50 usuarios concurrentes por organización) está totalmente cubierto por la combinación de servicio backend y configuración de workers descritos en la sección 5.1.2. Los regímenes por encima de 50 usuarios concurrentes (Tests 03 a 06) se incluyen únicamente como caracterización del techo del sistema y la optimización requerida para sostenerlos queda planteada como evolución futura, no como criterio de aceptación del MVP.

\section{Pruebas de Rendimiento (Frontend / Web Vitals)}\label{pruebas-de-rendimiento-frontend-web-vitals}

Las pruebas de rendimiento del frontend se ejecutaron sobre el despliegue productivo en Vercel (\texttt{https://orbitengine.lat}), combinando tres herramientas con perspectivas complementarias.

\subsection{Lighthouse (Auditoría de Laboratorio)}\label{lighthouse-auditoruxeda-de-laboratorio}

Se ejecutó Lighthouse desde Chrome DevTools sobre la página pública (\emph{landing}) y sobre las vistas internas del dashboard accedidas con sesión iniciada para \textbf{tres de las ocho organizaciones del piloto técnico} descritas en la sección 5.1.5: \textbf{Miss Peggy} (empresa real), \textbf{Lehgo} (empresa ficticia de prueba) y \textbf{Moda Andes} (empresa ficticia de prueba). Esta selección busca contrastar deliberadamente el comportamiento del frontend bajo dos perfiles de datos muy distintos: por un lado, los volúmenes reales y acotados de \textbf{Miss Peggy}, una de las dos pymes que adoptaron la plataforma; por otro lado, los volúmenes sintéticos elevados que aportan \textbf{Lehgo} y \textbf{Moda Andes} en su rol de empresas ficticias de prueba. La otra empresa real del piloto, \textbf{Frozt Bitez}, no se incluyó en este barrido de Lighthouse para evitar exponer información comercial específica de su operación; sus métricas se ejercitan indirectamente a través de las pruebas de carga del backend (sección 5.2). Los reportes individuales en PDF no se incluyen en este repositorio público; para acceder a los soportes correspondientes, consultar a los desarrolladores del proyecto.

\textbf{Tabla 5.3.1.} Puntajes y Web Vitals de Lighthouse por vista y organización. La columna \emph{Tipo} indica si el \emph{tenant} corresponde a una empresa real del piloto o a un \emph{tenant} sintético de prueba.

{\def\LTcaptype{none} % do not increment counter
\begin{longtable}[]{@{}
  >{\raggedright\arraybackslash}p{(\linewidth - 20\tabcolsep) * \real{0.1392}}
  >{\raggedright\arraybackslash}p{(\linewidth - 20\tabcolsep) * \real{0.1772}}
  >{\raggedright\arraybackslash}p{(\linewidth - 20\tabcolsep) * \real{0.1013}}
  >{\raggedright\arraybackslash}p{(\linewidth - 20\tabcolsep) * \real{0.0633}}
  >{\raggedright\arraybackslash}p{(\linewidth - 20\tabcolsep) * \real{0.1013}}
  >{\raggedright\arraybackslash}p{(\linewidth - 20\tabcolsep) * \real{0.1392}}
  >{\raggedright\arraybackslash}p{(\linewidth - 20\tabcolsep) * \real{0.0380}}
  >{\raggedright\arraybackslash}p{(\linewidth - 20\tabcolsep) * \real{0.0633}}
  >{\raggedright\arraybackslash}p{(\linewidth - 20\tabcolsep) * \real{0.0633}}
  >{\raggedright\arraybackslash}p{(\linewidth - 20\tabcolsep) * \real{0.0506}}
  >{\raggedright\arraybackslash}p{(\linewidth - 20\tabcolsep) * \real{0.0633}}@{}}
\toprule\noalign{}
\begin{minipage}[b]{\linewidth}\raggedright
Vista
\end{minipage} & \begin{minipage}[b]{\linewidth}\raggedright
Organización
\end{minipage} & \begin{minipage}[b]{\linewidth}\raggedright
Tipo
\end{minipage} & \begin{minipage}[b]{\linewidth}\raggedright
Perf.
\end{minipage} & \begin{minipage}[b]{\linewidth}\raggedright
Accesib.
\end{minipage} & \begin{minipage}[b]{\linewidth}\raggedright
Best Pract.
\end{minipage} & \begin{minipage}[b]{\linewidth}\raggedright
SEO
\end{minipage} & \begin{minipage}[b]{\linewidth}\raggedright
FCP
\end{minipage} & \begin{minipage}[b]{\linewidth}\raggedright
LCP
\end{minipage} & \begin{minipage}[b]{\linewidth}\raggedright
TBT
\end{minipage} & \begin{minipage}[b]{\linewidth}\raggedright
CLS
\end{minipage} \\
\midrule\noalign{}
\endhead
\bottomrule\noalign{}
\endlastfoot
Landing (/) & --- & Página & 92 & 96 & 96 & 92 & 1,1 s & 1,6 s & 0 ms & 0 \\
Dashboard & Lehgo & Ficticia & 90 & 96 & 100 & 83 & 1,2 s & 1,6 s & 0 ms & 0,017 \\
Dashboard & \textbf{Miss Peggy} & \textbf{Real} & 91 & 96 & 100 & 83 & 1,1 s & 1,6 s & 0 ms & 0,021 \\
Dashboard & Moda Andes & Ficticia & 91 & 96 & 100 & 83 & 1,1 s & 1,6 s & 0 ms & 0,017 \\
Inventario & Lehgo & Ficticia & 92 & 89 & 100 & 83 & 1,0 s & 1,5 s & 0 ms & 0 \\
Inventario & \textbf{Miss Peggy} & \textbf{Real} & 92 & 89 & 100 & 83 & 1,1 s & 1,5 s & 0 ms & 0 \\
Inventario & Moda Andes & Ficticia & 92 & 89 & 100 & 83 & 1,1 s & 1,5 s & 0 ms & 0 \\
Ventas & Lehgo & Ficticia & 92 & 88 & 100 & 83 & 1,1 s & 1,5 s & 0 ms & 0 \\
Ventas & \textbf{Miss Peggy} & \textbf{Real} & 93 & 88 & 100 & 83 & 1,0 s & 1,5 s & 0 ms & 0 \\
Ventas & Moda Andes & Ficticia & 93 & 88 & 100 & 83 & 1,0 s & 1,5 s & 0 ms & 0 \\
\end{longtable}
}

Las puntuaciones de Performance se mantienen entre \textbf{90 y 93} en todas las vistas y organizaciones, y las diferencias entre los tres \emph{tenants} medidos son menores a 1 punto. Esto es particularmente relevante porque la muestra confronta de manera explícita una empresa real (\textbf{Miss Peggy}) con dos empresas ficticias de prueba que cargan al sistema con volúmenes sintéticos mucho mayores (\textbf{Lehgo} y \textbf{Moda Andes}): el rendimiento percibido del frontend resulta \textbf{insensible al volumen y a los datos específicos} de cada organización dentro del rango medido, lo que valida que el coste del cliente no se ve afectado por los grandes volúmenes generados por las empresas ficticias de prueba. Los Core Web Vitals se sitúan dentro de las bandas verdes establecidas por Google (LCP ≤ 2,5 s, CLS ≤ 0,1, TBT bajo) tanto para los datos reales como para los datos sintéticos.

\subsection{PageSpeed Insights (Infraestructura de Google)}\label{pagespeed-insights-infraestructura-de-google}

PageSpeed Insights ejecutó Lighthouse 13.0.1 desde la infraestructura de Google con dos \emph{form factors} (escritorio y móvil emulado). Los reportes en PDF no se incluyen en este repositorio público; para acceder a los soportes correspondientes, consultar a los desarrolladores del proyecto.

\begin{quote}
\emph{Nota sobre el etiquetado de archivos}: los reportes nominados \texttt{Login\_PC.pdf} y \texttt{Login\_Tel.pdf} analizan la URL \texttt{https://orbitengine.lat/dashboard}. Como PageSpeed se ejecuta sin sesión iniciada, dicha URL fue redirigida automáticamente al formulario de login (\texttt{/login?reason=auth-required}). Por tanto, la página efectivamente medida es la \textbf{pantalla de acceso (Login)}, y bajo esa etiqueta se reportan los resultados.
\end{quote}

\textbf{Tabla 5.3.2.} PageSpeed Insights --- escritorio vs.~móvil.

{\def\LTcaptype{none} % do not increment counter
\begin{longtable}[]{@{}
  >{\raggedright\arraybackslash}p{(\linewidth - 20\tabcolsep) * \real{0.1538}}
  >{\raggedright\arraybackslash}p{(\linewidth - 20\tabcolsep) * \real{0.1429}}
  >{\raggedright\arraybackslash}p{(\linewidth - 20\tabcolsep) * \real{0.1209}}
  >{\raggedright\arraybackslash}p{(\linewidth - 20\tabcolsep) * \real{0.1429}}
  >{\raggedright\arraybackslash}p{(\linewidth - 20\tabcolsep) * \real{0.0659}}
  >{\raggedright\arraybackslash}p{(\linewidth - 20\tabcolsep) * \real{0.0330}}
  >{\raggedright\arraybackslash}p{(\linewidth - 20\tabcolsep) * \real{0.0549}}
  >{\raggedright\arraybackslash}p{(\linewidth - 20\tabcolsep) * \real{0.0549}}
  >{\raggedright\arraybackslash}p{(\linewidth - 20\tabcolsep) * \real{0.0549}}
  >{\raggedright\arraybackslash}p{(\linewidth - 20\tabcolsep) * \real{0.0549}}
  >{\raggedright\arraybackslash}p{(\linewidth - 20\tabcolsep) * \real{0.1209}}@{}}
\toprule\noalign{}
\begin{minipage}[b]{\linewidth}\raggedright
Vista
\end{minipage} & \begin{minipage}[b]{\linewidth}\raggedright
\emph{Form factor}
\end{minipage} & \begin{minipage}[b]{\linewidth}\raggedright
Rendimiento
\end{minipage} & \begin{minipage}[b]{\linewidth}\raggedright
Accesibilidad
\end{minipage} & \begin{minipage}[b]{\linewidth}\raggedright
Recom.
\end{minipage} & \begin{minipage}[b]{\linewidth}\raggedright
SEO
\end{minipage} & \begin{minipage}[b]{\linewidth}\raggedright
FCP
\end{minipage} & \begin{minipage}[b]{\linewidth}\raggedright
LCP
\end{minipage} & \begin{minipage}[b]{\linewidth}\raggedright
TBT
\end{minipage} & \begin{minipage}[b]{\linewidth}\raggedright
CLS
\end{minipage} & \begin{minipage}[b]{\linewidth}\raggedright
Speed Index
\end{minipage} \\
\midrule\noalign{}
\endhead
\bottomrule\noalign{}
\endlastfoot
Landing & Escritorio & \textbf{98} & 96 & 100 & 92 & 0,8 s & 0,8 s & 0 ms & 0 & 1,2 s \\
Landing & Móvil & 83 & 96 & 100 & 92 & 3,4 s & 3,4 s & 80 ms & 0 & 3,8 s \\
Acceso (Login) & Escritorio & \textbf{95} & 94 & 100 & 83 & 0,8 s & 1,4 s & 0 ms & 0 & 1,0 s \\
Acceso (Login) & Móvil & 87 & 94 & 100 & 83 & 3,2 s & 3,2 s & 20 ms & 0,001 & 3,2 s \\
\end{longtable}
}

Las dos observaciones más relevantes son:

\begin{itemize}
\tightlist
\item
  En \textbf{escritorio} todas las puntuaciones de Rendimiento se sitúan por encima de 95, con LCP por debajo de 1,5 s y SI ≤ 1,2 s, lo que confirma que el sitio cumple los Core Web Vitals con holgura para usuarios con conexión de banda ancha y dispositivos modernos.
\item
  En \textbf{móvil} la puntuación cae al rango 83--87, principalmente por LCP en torno a 3,2--3,4 s (banda \emph{needs improvement} de Google, entre 2,5 s y 4,0 s). El TBT permanece por debajo de los 100 ms y la CLS prácticamente nula, por lo que la interactividad y la estabilidad visual no se ven comprometidas; la causa de la caída es el coste de descarga sobre redes móviles emuladas.
\end{itemize}

\subsection{WebPageTest (Red Real con Captura de Video)}\label{webpagetest-red-real-con-captura-de-video}

WebPageTest 26.03 ejecutó pasadas reales sobre las páginas pública (\texttt{/}) y de acceso (\texttt{/dashboard} redirigido a \texttt{/login}), grabando video del rendering progresivo y produciendo el waterfall completo de cada recurso. Los archivos JSON, los CSV de requests y las capturas de los frames del video no se incluyen en este repositorio público; para acceder a los soportes correspondientes, consultar a los desarrolladores del proyecto.

\textbf{Tabla 5.3.3.} Métricas clave reportadas por WebPageTest.

{\def\LTcaptype{none} % do not increment counter
\begin{longtable}[]{@{}
  >{\raggedright\arraybackslash}p{(\linewidth - 18\tabcolsep) * \real{0.1505}}
  >{\raggedright\arraybackslash}p{(\linewidth - 18\tabcolsep) * \real{0.0538}}
  >{\raggedright\arraybackslash}p{(\linewidth - 18\tabcolsep) * \real{0.1290}}
  >{\raggedright\arraybackslash}p{(\linewidth - 18\tabcolsep) * \real{0.0645}}
  >{\raggedright\arraybackslash}p{(\linewidth - 18\tabcolsep) * \real{0.0860}}
  >{\raggedright\arraybackslash}p{(\linewidth - 18\tabcolsep) * \real{0.0538}}
  >{\raggedright\arraybackslash}p{(\linewidth - 18\tabcolsep) * \real{0.0538}}
  >{\raggedright\arraybackslash}p{(\linewidth - 18\tabcolsep) * \real{0.1183}}
  >{\raggedright\arraybackslash}p{(\linewidth - 18\tabcolsep) * \real{0.1290}}
  >{\raggedright\arraybackslash}p{(\linewidth - 18\tabcolsep) * \real{0.1613}}@{}}
\toprule\noalign{}
\begin{minipage}[b]{\linewidth}\raggedright
Vista
\end{minipage} & \begin{minipage}[b]{\linewidth}\raggedright
TTFB
\end{minipage} & \begin{minipage}[b]{\linewidth}\raggedright
Start render
\end{minipage} & \begin{minipage}[b]{\linewidth}\raggedright
FCP
\end{minipage} & \begin{minipage}[b]{\linewidth}\raggedright
LCP
\end{minipage} & \begin{minipage}[b]{\linewidth}\raggedright
TBT
\end{minipage} & \begin{minipage}[b]{\linewidth}\raggedright
CLS
\end{minipage} & \begin{minipage}[b]{\linewidth}\raggedright
Speed Index
\end{minipage} & \begin{minipage}[b]{\linewidth}\raggedright
Fully loaded
\end{minipage} & \begin{minipage}[b]{\linewidth}\raggedright
Bytes recibidos
\end{minipage} \\
\midrule\noalign{}
\endhead
\bottomrule\noalign{}
\endlastfoot
Landing (\texttt{/}) & 46 ms & --- & --- & --- & --- & --- & --- & \textasciitilde752 ms & 3,9 MB \\
Acceso (Login) & 88 ms & 600 ms & 911 ms & 1 044 ms & 59 ms & 0,004 & 785 & 752 ms & 3,9 MB \\
\end{longtable}
}

La secuencia visual del rendering progresivo quedó documentada en cinco capturas (frames sucesivos del video) por vista, que permiten verificar que la interfaz alcanza el estado \emph{Visually Complete} alrededor de los 900--1 000 ms desde el inicio de la navegación. Estas capturas están disponibles a través de los desarrolladores del proyecto.

\subsection{Cumplimiento de Core Web Vitals}\label{cumplimiento-de-core-web-vitals}

Consolidando las tres herramientas:

\begin{itemize}
\tightlist
\item
  \textbf{LCP} (objetivo ≤ 2,5 s): cumplido en escritorio (0,8--1,6 s) y por debajo del umbral en WebPageTest (1,0 s). En móvil emulado se ubica en 3,2--3,4 s (banda \emph{needs improvement}).
\item
  \textbf{CLS} (objetivo ≤ 0,1): cumplido en todas las vistas y herramientas (máximo registrado: 0,021 en Lighthouse Dashboard de Miss Peggy).
\item
  \textbf{TBT / INP} (objetivo TBT \textless{} 200 ms): cumplido con holgura. El máximo registrado es 80 ms en PageSpeed móvil de la landing.
\end{itemize}

\section{Síntesis de Cumplimiento de Requisitos No Funcionales}\label{suxedntesis-de-cumplimiento-de-requisitos-no-funcionales}

A continuación se presenta el cumplimiento consolidado de los Requisitos No Funcionales (RNF) que aplican a la validación técnica:

\begin{itemize}
\item
  RNF-01: P95 de la API menor a 500 ms con hasta 50 usuarios concurrentes.
  Resultado: Cumplido para el entorno de despliegue dimensionado para la aplicación. Bajo carga normal, los endpoints CRUD se mantienen dentro del rango esperado y el régimen de 50 usuarios se sostiene sin colapsos. La optimización para escenarios con más de 50 usuarios concurrentes queda planteada como una evolución futura.
  Evidencia: Ver Sección 5.2.4 (Test 01 y 02) y 5.2.7.
\item
  RNF-02: Disponibilidad mensual mayor o igual al 95 \%.
  Resultado: Cumplido en aproximación. Aunque no se realizó una medición formal de uptime mensual durante la validación, las estadísticas de carga obtenidas (0 \% de fallos en carga normal, 1,7 \% en estrés moderado y comportamiento estable del despliegue durante todas las pruebas) permiten proyectar razonablemente que la plataforma se mantendría por encima del 95 \% de disponibilidad en una ventana mensual bajo el perfil de uso esperado. La medición formal continua, apoyada en los registros nativos de Railway y Vercel, queda planteada para futuras referencias.
  Evidencia: Ver Sección 5.1.2.
\item
  RNF-07: Interfaz responsive y funcional desde 375 px.
  Resultado: Cumplido. Auditorías de Lighthouse y PageSpeed en dispositivos móviles arrojan una accesibilidad igual o superior a 88 sin errores bloqueantes.
  Evidencia: Ver Tablas 5.3.1 y 5.3.2.
\item
  RNF-09: Escalabilidad horizontal de la capa de aplicación.
  Resultado: Cumplido a nivel arquitectónico. El sistema soporta multi-tenancy bajo carga sin filtraciones, como se verifica en la sección 5.2.1, y la degradación bajo picos de carga es la esperada para una configuración de 2 a 4 workers sobre una sola instancia, según el escenario. Esto confirma que la arquitectura admite escalado horizontal sin cambios estructurales.
  Evidencia: Ver Sección 5.2.5.
\end{itemize}

\section{Interpretación y Limitaciones}\label{interpretaciuxf3n-y-limitaciones}

\subsection{Interpretación Breve de los Hallazgos}\label{interpretaciuxf3n-breve-de-los-hallazgos}

En el plano del frontend, las tres herramientas convergen en una imagen coherente: la aplicación obtiene puntuaciones de Lighthouse entre 90 y 93 en todas las vistas y organizaciones, los Core Web Vitals se sitúan en banda verde para escritorio y red real, y las diferencias entre tenants son menores a un punto. Esto indica que el rendimiento percibido por el usuario final no depende del volumen de datos de la organización ni de la vista visitada, sino principalmente del \emph{form factor} y de la calidad de la red, con una penalización esperable en móvil emulado.

En el plano del backend, los seis escenarios de Locust dibujan una curva de degradación característica de un servicio con concurrencia interna acotada (2 workers para los Tests 01--05 y 4 workers para el Test 06) sobre una sola instancia sin réplicas: bajo carga normal (8 usuarios) el sistema atiende sin fallos y con medianas en el rango de cientos de milisegundos, bajo estrés moderado (50 usuarios) sostiene 16 RPS con una tasa de fallos del 1,7 \%, y bajo pico sostenido (\textasciitilde200 usuarios, ya con 4 workers) conserva \textasciitilde12 RPS con una tasa de fallos cercana al 19 \%. La consistencia de las medianas de los endpoints CRUD a lo largo de los escenarios (variaciones de menos del doble entre 8 y 200 usuarios) sugiere que la base de datos no está saturada; el factor limitante observado es la capacidad de procesamiento de los workers configurados sobre la única instancia desplegada.

Confrontados con los criterios del Capítulo 3 (tabla 5.4), los resultados muestran cumplimiento integral en la dimensión de frontend y cumplimiento del backend dentro del rango operativo previsto para la aplicación: hasta 50 usuarios concurrentes por organización el sistema atiende sin colapsos y con latencias acotadas, mientras que los regímenes por encima de ese límite se reportan únicamente como caracterización del techo y como insumo para la planificación de evoluciones futuras. La arquitectura multi-tenant se mantiene íntegra durante toda la batería de pruebas: ninguna petición devuelve datos cruzados entre organizaciones, lo que valida en escenarios de carga real el mecanismo de aislamiento por \texttt{organization\_id} descrito en el Capítulo 4.

\subsection{Limitaciones de Escalabilidad Inherentes al Concepto del Proyecto}\label{limitaciones-de-escalabilidad-inherentes-al-concepto-del-proyecto}

Los límites observados en las pruebas deben leerse a la luz del concepto y del alcance del proyecto, no como deficiencias del producto. OrbitEngine es un trabajo de grado desarrollado por un equipo de tres personas con dedicación parcial durante siete meses, cuyo objetivo es entregar un MVP funcional con presupuesto académico, y las decisiones de infraestructura y de arquitectura son consecuentes con ese marco:

\begin{itemize}
\tightlist
\item
  \textbf{Infraestructura en planes básicos}. El backend se ejecuta sobre un único servicio en Railway, sin réplicas horizontales ni \emph{auto-scaling} configurados. La concurrencia interna se ajustó al escenario de prueba (2 workers de FastAPI para los Tests 01--05 y 4 workers para el Test 06), pero la capacidad total queda condicionada por los recursos de esa única instancia. La base de datos PostgreSQL es la instancia gestionada incluida en el mismo plan, sin réplicas de lectura. El frontend se sirve desde el plan estándar de Vercel.
\item
  \textbf{Ausencia de capa de caché distribuida}. No se incorporó Redis ni un CDN privado para la API; los endpoints de agregación (\texttt{/dashboard/stats}, \texttt{/sales/stats}) consultan PostgreSQL directamente en cada petición.
\item
  \textbf{Arquitectura monolítica multi-tenant compartiendo una sola base de datos}. Esta decisión es intencional y está documentada en el Capítulo 3: ofrece un coste por tenant muy bajo y simplifica la operación, pero sitúa el techo de concurrencia del producto entero en los recursos del servicio backend (2--4 workers según escenario).
\item
  \textbf{Ejecución de las pruebas desde un único cliente}. Locust se lanzó desde la estación de trabajo del equipo en Bogotá; no se utilizaron clientes distribuidos geográficamente, lo que descarta escenarios que evalúen latencia de red multi-región o saturación de la salida del cliente.
\item
  \textbf{Plan de validación dimensionado para una pyme tipo}. El régimen de carga normal (\textasciitilde8 usuarios concurrentes) corresponde al tamaño de equipo objetivo descrito en el Capítulo 1; los regímenes de 50 y 200 usuarios se incluyen para caracterizar el techo del sistema, no como criterio de aceptación del MVP.
\item
  \textbf{Alcance temporal acotado al cierre del MVP}. Las pruebas reportadas se ejecutaron en una ventana puntual de validación; no se diseñó un esquema de monitoreo continuo de SLOs ni de pruebas de carga periódicas en producción.
\end{itemize}

Estas limitaciones encuadran el alcance del capítulo: los hallazgos son representativos del comportamiento del MVP en su despliegue de producción del proyecto de grado, y delimitan claramente los aspectos que requerirían inversión adicional para evolucionar OrbitEngine hacia un servicio SaaS multi-región a gran escala. Las acciones concretas derivadas de estos hallazgos se discuten en el capítulo de Conclusiones.

    \chapter{Capítulo 6. Resultados de Usuarios}\label{capuxedtulo-6.-resultados-de-usuarios}

\begin{quote}
\textbf{Alcance de este capítulo.} Las secciones que siguen presentan exclusivamente los resultados de la \textbf{validación con usuarios reales}: eficiencia operativa, completitud de tareas, usabilidad percibida, satisfacción y hallazgos cualitativos de las entrevistas. Los \textbf{resultados de la validación técnica} (pruebas de carga con Locust y rendimiento web con Lighthouse / PageSpeed Insights / WebPageTest) se encuentran en el \textbf{Capítulo 5. Resultados Técnicos}. Las \textbf{acciones de mejora} derivadas de los hallazgos aquí descritos (tanto técnicas como funcionales) se presentan en el \textbf{Capítulo 7. Conclusiones y Trabajo Futuro}.
\end{quote}

\begin{center}\rule{0.5\linewidth}{0.5pt}\end{center}

\section{Marco Metodológico de la Validación con Usuarios}\label{marco-metodoluxf3gico-de-la-validaciuxf3n-con-usuarios}

\subsection{Objetivos}\label{objetivos}

La validación con usuarios persiguió tres objetivos concretos y verificables:

\begin{enumerate}
\def\labelenumi{\arabic{enumi}.}
\tightlist
\item
  \textbf{Medir el impacto en la eficiencia operativa} de las empresas participantes antes y después de la adopción de OrbitEngine, cuantificando la reducción de tiempos en tareas administrativas clave y la variación en la tasa de error en operaciones de inventario.
\item
  \textbf{Evaluar la usabilidad del sistema} mediante el instrumento estandarizado \emph{System Usability Scale} (SUS) y pruebas de tareas guiadas con usuarios reales, determinando si el sistema supera el umbral de usabilidad aceptable establecido en la literatura.
\item
  \textbf{Caracterizar la satisfacción específica por módulo} mediante NPS (\emph{Net Promoter Score}) y CSAT (\emph{Customer Satisfaction Score}), e identificar hallazgos cualitativos relevantes a través de entrevistas semiestructuradas de cierre.
\end{enumerate}

\subsection{Diseño de Investigación}\label{diseuxf1o-de-investigaciuxf3n}

El presente estudio adoptó un \textbf{diseño de estudio de caso múltiple} (Yin, 2018) con componente \textbf{cuasi-experimental de tipo pre/post} (Campbell \& Stanley, 1963). Se eligió este diseño porque:

\begin{itemize}
\tightlist
\item
  El número de empresas disponibles (N = 3) impide inferencia estadística a la población y hace inadecuado un diseño experimental aleatorizado.
\item
  Las empresas participantes son pymes reales con operaciones activas; no es posible asignarlas aleatoriamente a condiciones de tratamiento o control.
\item
  El interés científico radica en comprender \textbf{cómo y en qué medida} cambia la eficiencia operativa al adoptar el sistema, más que en establecer causalidad universal.
\end{itemize}

El enfoque es \textbf{mixto} (Creswell \& Plano Clark, 2018): combina datos cuantitativos (tiempos pre/post, tasas de error, scores SUS/NPS/CSAT, telemetría) con datos cualitativos (entrevistas semiestructuradas con codificación temática), siendo los cualitativos confirmatorios y explicativos de los cuantitativos.

\textbf{Nota sobre el cambio de N respecto al Capítulo 5.} Las pruebas técnicas del Capítulo 5 se ejecutaron cuando la plataforma contaba con \textbf{dos empresas reales} en producción (Frozt Bitez y Miss Peggy). En el intervalo transcurrido entre el cierre de esas pruebas y el inicio de la Fase 5 de validación con usuarios (27 de abril de 2026), se incorporó una \textbf{tercera empresa real}, Luana Handmade, que también participó en el estudio de caso. El piloto de usuarios reales comprende, por tanto, \textbf{N = 3 empresas}.

\subsection{Participantes}\label{participantes}

\textbf{Criterios de inclusión:}

\begin{itemize}
\tightlist
\item
  Empresa pyme del sector comercio o servicios con operaciones activas en Colombia.
\item
  Disposición a registrar datos reales de inventario, ventas y clientes en la plataforma durante la Fase 5.
\item
  Al menos un usuario designado con rol Administrador y al menos un usuario con rol Vendedor.
\item
  Consentimiento informado firmado por el representante legal o dueño de la empresa.
\end{itemize}

\textbf{Criterios de exclusión:}

\begin{itemize}
\tightlist
\item
  Empresas sin acceso a internet estable (requisito técnico de la plataforma).
\item
  Empresas en proceso de cierre o reestructuración durante el período de validación.
\end{itemize}

\textbf{Muestra resultante:} tres empresas (N = 3), con un total de 7 usuarios individuales distribuidos entre las tres organizaciones: 3 de Frozt Bitez (1 Administrador + 2 Vendedores), 3 de Miss Peggy (1 Administradora + 2 Vendedores) y 1 de Luana Handmade (1 Administradora).

\subsection{Instrumentos}\label{instrumentos}

{\def\LTcaptype{none} % do not increment counter
\begin{longtable}[]{@{}
  >{\raggedright\arraybackslash}p{(\linewidth - 6\tabcolsep) * \real{0.2044}}
  >{\raggedright\arraybackslash}p{(\linewidth - 6\tabcolsep) * \real{0.1602}}
  >{\raggedright\arraybackslash}p{(\linewidth - 6\tabcolsep) * \real{0.3370}}
  >{\raggedright\arraybackslash}p{(\linewidth - 6\tabcolsep) * \real{0.2983}}@{}}
\toprule\noalign{}
\begin{minipage}[b]{\linewidth}\raggedright
Instrumento
\end{minipage} & \begin{minipage}[b]{\linewidth}\raggedright
Tipo
\end{minipage} & \begin{minipage}[b]{\linewidth}\raggedright
Propósito
\end{minipage} & \begin{minipage}[b]{\linewidth}\raggedright
Momento de aplicación
\end{minipage} \\
\midrule\noalign{}
\endhead
\bottomrule\noalign{}
\endlastfoot
\textbf{SUS} (\emph{System Usability Scale}) & Cuantitativo & Evaluar la usabilidad percibida del sistema (score 0--100) & Post-uso, al cierre de la Fase 5 \\
\textbf{NPS} (\emph{Net Promoter Score}) & Cuantitativo & Medir la probabilidad de recomendación del sistema (0--10) & Post-uso, al cierre de la Fase 5 \\
\textbf{CSAT por módulo} & Cuantitativo & Satisfacción específica con cada módulo del sistema (1--5) & Post-uso, al cierre de la Fase 5 \\
\textbf{Pruebas de tareas guiadas} & Cuantitativo / observacion al & Medir completitud, tiempo y errores en tareas representativas & Durante la Fase 5 \\
\textbf{Entrevistas semiestructuradas} & Cualitativo & Explorar impacto percibido, dificultades y sugerencias & Al cierre de la Fase 5 \\
\textbf{Telemetría productiva} & Cuantitativo & Registrar uso real del sistema en producción & Continuo durante la Fase 5 \\
\textbf{Registro de tiempos pre/post} & Cuantitativo & Comparar duración de tareas administrativas antes y después & Pre-implementación (retrospectivo) y durante la Fase 5 \\
\textbf{Registro de errores de inventario} & Cuantitativo & Medir tasa de discrepancias antes y después & Pre-implementación (retrospectivo) y durante la Fase 5 \\
\end{longtable}
}

\subsection{Procedimiento y Cronograma (Fase 5)}\label{procedimiento-y-cronograma-fase-5}

La Fase 5 de validación con usuarios se ejecutó entre el \textbf{27 de abril de 2026} y el \textbf{4 de mayo de 2026} (ocho días calendario). El procedimiento siguió las etapas que se describen a continuación:

\textbf{Fase inicial (27 abril -- 30 abril 2026):}

\begin{enumerate}
\def\labelenumi{\arabic{enumi}.}
\tightlist
\item
  Firma de consentimiento informado con el representante de cada empresa.
\item
  Sesión de onboarding y carga inicial de datos (productos, clientes, stock inicial).
\item
  Registro retrospectivo de tiempos pre-implementación mediante entrevista estructurada con el usuario principal de cada empresa.
\item
  Inicio del uso productivo real con la plataforma.
\end{enumerate}

\textbf{Fase de evaluación (1 mayo -- 4 mayo 2026):}

\begin{enumerate}
\def\labelenumi{\arabic{enumi}.}
\tightlist
\item
  Sesión de pruebas de tareas guiadas con cada usuario participante (duración estimada: 45--60 minutos por sesión).
\item
  Aplicación de la encuesta SUS (10 ítems, 5--10 minutos).
\item
  Aplicación del NPS y el CSAT por módulo (5 minutos).
\item
  Entrevista semiestructurada de cierre (20--30 minutos por empresa).
\item
  Extracción de datos de telemetría de la base de datos de producción.
\end{enumerate}

\subsection{Hipótesis a Contrastar}\label{hipuxf3tesis-a-contrastar}

A partir del diseño cuasi-experimental descrito en la sección 1.2.3 del Capítulo 1, se formalizan las siguientes hipótesis de investigación:

\begin{quote}
\textbf{H1. Hipótesis de Eficiencia Operativa}
La adopción de OrbitEngine reduce en al menos un 30\% el tiempo dedicado a tareas administrativas clave (registro de ventas, actualización de inventario y generación de reportes) en las empresas piloto, medido como diferencia entre el tiempo promedio pre-implementación y el tiempo promedio post-implementación para el mismo conjunto de tareas.
\end{quote}

\begin{quote}
\textbf{H2. Hipótesis de Precisión en Inventario}
La adopción de OrbitEngine reduce en al menos un 40\% la tasa de discrepancias entre el stock físico y el registrado en las empresas piloto, medida como proporción de ítems con discrepancia sobre el total de ítems auditados, antes y después de la implementación.
\end{quote}

\begin{quote}
\textbf{H3. Hipótesis de Usabilidad}
OrbitEngine obtiene un puntaje SUS promedio superior a 68 puntos (umbral de usabilidad ``aceptable'' según Bangor et al., 2008) en la evaluación con usuarios de las empresas piloto, indicando que el sistema es percibido como usable por el segmento objetivo.
\end{quote}

\subsection{Consideraciones Éticas}\label{consideraciones-uxe9ticas}

La validación con usuarios se desarrolló bajo los siguientes principios éticos:

\begin{itemize}
\tightlist
\item
  \textbf{Consentimiento informado}: todos los participantes firmaron un formulario de consentimiento antes de iniciar cualquier actividad de recolección de datos. El formulario explicitó el propósito del estudio, el uso que se daría a los datos y el derecho a retirarse sin consecuencias.
\item
  \textbf{Anonimización}: los datos individuales de desempeño (tiempos de tarea, errores observados, respuestas a encuestas) se presentan de forma agregada. Las citas textuales se atribuyen únicamente a la empresa (no al individuo) salvo autorización explícita.
\item
  \textbf{Protección de información comercial}: los datos operativos de las empresas (inventario, ventas, clientes) son propiedad de cada organización. Solo se reportan métricas derivadas y agregadas; no se incluyen datos nominales de clientes ni precios específicos de productos.
\item
  \textbf{Conflicto de interés}: el equipo investigador es el mismo que desarrolló la plataforma. Esta limitación se reconoce explícitamente en la sección 6.10 y se mitiga mediante el uso de instrumentos estandarizados (SUS, NPS) y protocolos de entrevista predefinidos.
\item
  \textbf{Posibilidad de retiro}: cualquier empresa pudo retirarse del estudio en cualquier momento sin que ello afectara su acceso continuado a la plataforma.
\end{itemize}

\section{Caracterización de las Empresas Piloto}\label{caracterizaciuxf3n-de-las-empresas-piloto}

La siguiente tabla sintetiza los atributos más relevantes de las tres empresas que participaron en la validación con usuarios.

\textbf{Tabla 6.2.1.} Caracterización de las empresas piloto del estudio de caso.

{\def\LTcaptype{none} % do not increment counter
\begin{longtable}[]{@{}
  >{\raggedright\arraybackslash}p{(\linewidth - 6\tabcolsep) * \real{0.1684}}
  >{\raggedright\arraybackslash}p{(\linewidth - 6\tabcolsep) * \real{0.2559}}
  >{\raggedright\arraybackslash}p{(\linewidth - 6\tabcolsep) * \real{0.2290}}
  >{\raggedright\arraybackslash}p{(\linewidth - 6\tabcolsep) * \real{0.3468}}@{}}
\toprule\noalign{}
\begin{minipage}[b]{\linewidth}\raggedright
Atributo
\end{minipage} & \begin{minipage}[b]{\linewidth}\raggedright
Frozt Bitez
\end{minipage} & \begin{minipage}[b]{\linewidth}\raggedright
Miss Peggy
\end{minipage} & \begin{minipage}[b]{\linewidth}\raggedright
Luana Handmade
\end{minipage} \\
\midrule\noalign{}
\endhead
\bottomrule\noalign{}
\endlastfoot
\textbf{Sector} & Alimentos / Comercio electrónico (uvas congeladas acidulces) & Comercio minorista (naturismo y belleza) & Artesanías / Confección artesanal (tejidos en trapillo y macramé) \\
\textbf{Antigüedad de la empresa} & 1-2 años & Más de 5 años & 4-5 años \\
\textbf{Número de empleados} & 3 (fundador + 2 colaboradores de ventas) & 3 (administradora + 2 colaboradores de ventas) & 1 (fundadora y única trabajadora) \\
\textbf{Número de usuarios participantes en el estudio} & 3 & 3 & 1 \\
\textbf{Roles representados} & Administrador y Vendedor & Administrador y Vendedor & Administrador \\
\textbf{Rol del informante principal} & Fundador / Dueño & Dueña / Administradora & Fundadora / Propietaria \\
\textbf{Herramientas de gestión previas} & WooCommerce (tienda pública) + WhatsApp (atención y confirmación de pedidos) & Excel (hoja de inventario propia) & Cuaderno físico + WhatsApp (sin sistema digital) \\
\textbf{Incorporación al piloto} & Fase 5 (27 abr 2026); también presente en piloto técnico del Cap. 5 & Fase 5 (27 abr 2026); también presente en piloto técnico del Cap. 5 & Fase 5 (27 abr 2026); incorporada entre el cierre del Cap. 5 y el inicio de la validación con usuarios \\
\end{longtable}
}

\section{Eficiencia Operativa (pre/post)}\label{eficiencia-operativa-prepost}

\subsection{Tiempos en Tareas Administrativas}\label{tiempos-en-tareas-administrativas}

Los tiempos pre-implementación se recogieron mediante entrevista estructurada retrospectiva al inicio de la Fase 5, solicitando al informante principal de cada empresa que estimara el tiempo promedio que dedicaba a cada tarea antes de usar OrbitEngine. Los tiempos post-implementación se midieron durante las sesiones de prueba de la Fase 5.

\textbf{Tabla 6.3.1.} Comparación de tiempos medios por tarea administrativa (minutos).

{\def\LTcaptype{none} % do not increment counter
\begin{longtable}[]{@{}
  >{\raggedright\arraybackslash}p{(\linewidth - 8\tabcolsep) * \real{0.3939}}
  >{\raggedright\arraybackslash}p{(\linewidth - 8\tabcolsep) * \real{0.1414}}
  >{\raggedright\arraybackslash}p{(\linewidth - 8\tabcolsep) * \real{0.1616}}
  >{\raggedright\arraybackslash}p{(\linewidth - 8\tabcolsep) * \real{0.1717}}
  >{\raggedright\arraybackslash}p{(\linewidth - 8\tabcolsep) * \real{0.1313}}@{}}
\toprule\noalign{}
\begin{minipage}[b]{\linewidth}\raggedright
Tarea
\end{minipage} & \begin{minipage}[b]{\linewidth}\raggedright
Empresa
\end{minipage} & \begin{minipage}[b]{\linewidth}\raggedright
Tiempo pre (min)
\end{minipage} & \begin{minipage}[b]{\linewidth}\raggedright
Tiempo post (min)
\end{minipage} & \begin{minipage}[b]{\linewidth}\raggedright
Reducción (\%)
\end{minipage} \\
\midrule\noalign{}
\endhead
\bottomrule\noalign{}
\endlastfoot
Registro de una venta & Frozt Bitez & 6.0 & 3.5 & −42 \% \\
Registro de una venta & Miss Peggy & 5.0 & 3.0 & −40 \% \\
Registro de una venta & Luana Handmade & 6.0 & 3.7 & −38 \% \\
Actualización de stock (un producto) & Frozt Bitez & 4.0 & 1.5 & −63 \% \\
Actualización de stock (un producto) & Miss Peggy & 4.0 & 1.4 & −65 \% \\
Actualización de stock (un producto) & Luana Handmade & 4.0 & 1.8 & −55 \% \\
Generación de reporte de ventas semanal & Frozt Bitez & 35.0 & 1.2 & −97 \% \\
Generación de reporte de ventas semanal & Miss Peggy & 60.0 & 1.4 & −98 \% \\
Generación de reporte de ventas semanal & Luana Handmade & 75.0 & 1.6 & −98 \% \\
Consulta de historial de un cliente & Frozt Bitez & 10.0 & 1.5 & −85 \% \\
Consulta de historial de un cliente & Miss Peggy & 8.0 & 1.3 & −84 \% \\
Consulta de historial de un cliente & Luana Handmade & 12.0 & 1.8 & −85 \% \\
\textbf{Promedio global} & \textbf{Todas} & \textbf{19.1} & \textbf{2.0} & \textbf{−71 \%} \\
\end{longtable}
}

\subsection{Tasa de Error en Operaciones de Inventario}\label{tasa-de-error-en-operaciones-de-inventario}

La tasa de error se definió como la proporción de ítems auditados que presentaban discrepancia entre el stock físico y el stock registrado, sobre el total de ítems auditados en una sesión de conteo.

La comparación pre/post para esta métrica solo es metodológicamente válida en el caso de \textbf{Miss Peggy}, empresa que contaba con un registro de inventario formalizado en Excel antes de la implementación de OrbitEngine, lo que permitió auditar la misma muestra de referencias en ambos momentos bajo condiciones comparables. Las otras dos empresas del piloto (\textbf{Frozt Bitez} y \textbf{Luana Handmade}) no disponían de un registro de inventario previo estructurado: Frozt Bitez gestionaba el stock a través de WooCommerce sin realizar cotejos sistemáticos entre el sistema y el inventario físico, mientras que Luana Handmade llevaba el control en un cuaderno sin formato estandarizado por SKU. En ambos casos, la ausencia de una línea de base pre-implementación impide calcular una reducción de tasa de error comparable, por lo que quedan excluidas del análisis cuantitativo de esta sección.

\textbf{Tabla 6.3.2.} Tasa de error en inventario antes y después de la implementación en Miss Peggy.

{\def\LTcaptype{none} % do not increment counter
\begin{longtable}[]{@{}
  >{\raggedright\arraybackslash}p{(\linewidth - 12\tabcolsep) * \real{0.1603}}
  >{\raggedright\arraybackslash}p{(\linewidth - 12\tabcolsep) * \real{0.1450}}
  >{\raggedright\arraybackslash}p{(\linewidth - 12\tabcolsep) * \real{0.1298}}
  >{\raggedright\arraybackslash}p{(\linewidth - 12\tabcolsep) * \real{0.1679}}
  >{\raggedright\arraybackslash}p{(\linewidth - 12\tabcolsep) * \real{0.1527}}
  >{\raggedright\arraybackslash}p{(\linewidth - 12\tabcolsep) * \real{0.1374}}
  >{\raggedright\arraybackslash}p{(\linewidth - 12\tabcolsep) * \real{0.1069}}@{}}
\toprule\noalign{}
\begin{minipage}[b]{\linewidth}\raggedright
Ítems auditados (pre)
\end{minipage} & \begin{minipage}[b]{\linewidth}\raggedright
Discrepancias (pre)
\end{minipage} & \begin{minipage}[b]{\linewidth}\raggedright
Tasa de error pre
\end{minipage} & \begin{minipage}[b]{\linewidth}\raggedright
Ítems auditados (post)
\end{minipage} & \begin{minipage}[b]{\linewidth}\raggedright
Discrepancias (post)
\end{minipage} & \begin{minipage}[b]{\linewidth}\raggedright
Tasa de error post
\end{minipage} & \begin{minipage}[b]{\linewidth}\raggedright
Reducción (pp)
\end{minipage} \\
\midrule\noalign{}
\endhead
\bottomrule\noalign{}
\endlastfoot
25 & 4 & \textbf{16,0 \%} & 25 & 1 & \textbf{4,0 \%} & \textbf{−12 pp} \\
\end{longtable}
}

La tasa de error pasó de 16,0 \% a 4,0 \%, lo que representa una reducción absoluta de 12 puntos porcentuales (−75 \% relativo). Los cuatro ítems con discrepancia en la medición pre correspondían a productos de alta rotación cuyas ventas no habían sido descontadas del registro en Excel por la práctica de actualización en lotes; bajo OrbitEngine, cada transacción de venta descuenta el stock de forma automática e inmediata. La única discrepancia registrada en la medición post fue de naturaleza operacional (una recepción de mercancía pendiente de ingreso), no un error de registro sistemático.

Esta reducción de 12 pp no alcanza el umbral de 40 pp establecido en H2, por lo que el veredicto sobre esa hipótesis es mixto: la dirección del cambio es la esperada y el mecanismo causal es identificable, pero la magnitud no llega al criterio de aceptación definido. El análisis detallado de H2 se presenta en la sección 6.9.2.

\begin{quote}
\textbf{Nota metodológica.} Los datos pre-implementación son de naturaleza retrospectiva y auto-reportada. Las limitaciones asociadas a este tipo de medición se discuten en la sección 6.10.
\end{quote}

\subsection{Síntesis Cuantitativa de la Mejora}\label{suxedntesis-cuantitativa-de-la-mejora}

Las tres empresas del piloto muestran reducciones de tiempo que superan ampliamente el umbral del 30 \% establecido en H1 en las cuatro tareas administrativas medidas. La tarea con menor reducción es el registro de ventas (Miss Peggy y Frozt Bitez: −40 \% y −42 \% respectivamente; Luana: −38 \%), y la mayor es la generación del reporte de ventas semanal, que en los tres casos supera el −97 \%: de 35 minutos a 1.2 minutos en Frozt Bitez, de 60 minutos a 1.4 minutos en Miss Peggy y de 75 minutos a 1.6 minutos en Luana. Este cuello de botella, que en las tres organizaciones ocupaba entre 35 y 75 minutos de trabajo manual semanal, se convierte en la ganancia de eficiencia más dramática y uniforme del piloto. Las reducciones promedio por empresa son 72 \% (Frozt Bitez), 72 \% (Miss Peggy) y 69 \% (Luana Handmade), con un promedio global de \textbf{71 \%}, más del doble del umbral de H1.

En cuanto a la precisión de inventario, Miss Peggy es la única empresa del piloto con datos pre/post comparables: el registro en Excel previo permitió auditar la misma muestra de 25 SKUs antes y después de la implementación. La tasa de error pasó de 16.0 \% (4 discrepancias en 25 SKUs) a 4.0 \% (1 discrepancia), lo que representa una reducción absoluta de 12 pp (−75 \% relativo). Esta reducción no alcanza el umbral de 40 pp de H2, pero confirma que OrbitEngine introduce disciplina real en el control de inventario: los cuatro ítems con discrepancia pre correspondían a productos de alta rotación cuyas ventas no se habían descontado del Excel por el registro en lotes; bajo OrbitEngine, cada venta descuenta automáticamente el stock. La discrepancia post restante fue operacional (una recepción de mercancía aún no ingresada), no sistemática. Frozt Bitez y Luana, que carecían de registro formal previo, registraron tasas post de 0 \% y 5.6 \% respectivamente, mostrando que OrbitEngine introduce un nivel de control antes inexistente en ambas organizaciones.

\section{Pruebas de Tareas Guiadas}\label{pruebas-de-tareas-guiadas}

Las pruebas de tareas guiadas se realizaron con cada usuario participante en una sesión individual facilitada por un miembro del equipo investigador. Se eligieron tareas representativas de los módulos principales del sistema. Se registraron: éxito/fallo en la completitud, tiempo medio de ejecución, número de errores (acciones incorrectas o recuperaciones) y severidad observada.

\textbf{Escala de severidad de errores}: 0 = sin errores; 1 = error leve (usuario se recupera solo en \textless{} 30 s); 2 = error moderado (requiere ayuda del facilitador); 3 = error crítico (tarea no completada).

\textbf{Tabla 6.4.1.} Resultados de las pruebas de tareas guiadas (agregado por tarea).

{\def\LTcaptype{none} % do not increment counter
\begin{longtable}[]{@{}
  >{\raggedright\arraybackslash}p{(\linewidth - 12\tabcolsep) * \real{0.0180}}
  >{\raggedright\arraybackslash}p{(\linewidth - 12\tabcolsep) * \real{0.3174}}
  >{\raggedright\arraybackslash}p{(\linewidth - 12\tabcolsep) * \real{0.2156}}
  >{\raggedright\arraybackslash}p{(\linewidth - 12\tabcolsep) * \real{0.1018}}
  >{\raggedright\arraybackslash}p{(\linewidth - 12\tabcolsep) * \real{0.0958}}
  >{\raggedright\arraybackslash}p{(\linewidth - 12\tabcolsep) * \real{0.1557}}
  >{\raggedright\arraybackslash}p{(\linewidth - 12\tabcolsep) * \real{0.0958}}@{}}
\toprule\noalign{}
\begin{minipage}[b]{\linewidth}\raggedright
N°
\end{minipage} & \begin{minipage}[b]{\linewidth}\raggedright
Tarea
\end{minipage} & \begin{minipage}[b]{\linewidth}\raggedright
Usuarios que completaron (n / total)
\end{minipage} & \begin{minipage}[b]{\linewidth}\raggedright
Tasa de éxito (\%)
\end{minipage} & \begin{minipage}[b]{\linewidth}\raggedright
Tiempo medio (s)
\end{minipage} & \begin{minipage}[b]{\linewidth}\raggedright
Errores observados (media)
\end{minipage} & \begin{minipage}[b]{\linewidth}\raggedright
Severidad máxima
\end{minipage} \\
\midrule\noalign{}
\endhead
\bottomrule\noalign{}
\endlastfoot
T1 & Iniciar sesión y navegar al Dashboard & 7 / 7 & \textbf{100 \%} & 76 & 0.00 & 0 \\
T2 & Crear un producto nuevo en inventario & 7 / 7 & \textbf{100 \%} & 190 & 1.00 & 1 \\
T3 & Registrar un movimiento de entrada de stock & 7 / 7 & \textbf{100 \%} & 110 & 0.43 & 1 \\
T4 & Registrar una venta nueva (múltiples productos) & 7 / 7 & \textbf{100 \%} & 211 & 0.57 & 1 \\
T5 & Consultar el historial de ventas con filtro de fechas & 7 / 7 & \textbf{100 \%} & 101 & 0.71 & 1 \\
T6 & Buscar un cliente y revisar su historial de compras & 7 / 7 & \textbf{100 \%} & 87 & 0.00 & 0 \\
T7 & Exportar el listado de inventario a Excel & 7 / 7 & \textbf{100 \%} & 84 & 0.43 & 2 \\
T8 & Crear un usuario nuevo con rol Vendedor & 3 / 3 admins ⁵ & \textbf{100 \%} & 154 & 0.33 & 2 \\
\end{longtable}
}

\begin{quote}
⁵ T8 aplica únicamente a usuarios con rol Administrador: U1 (Frozt Bitez), U4 (Miss Peggy) y U7 (Luana Handmade).
\end{quote}

\textbf{Datos disponibles. Luana Handmade (U7).} La administradora completó las 8 tareas con una tasa de éxito del 100 \% (8/8). Tiempos por tarea (segundos): T1 = 92, T2 = 195, T3 = 112, T4 = 218, T5 = 97, T6 = 103, T7 = 82, T8 = 168. Total de errores: 5 en toda la sesión (promedio 0.63/tarea). Severidad máxima alcanzada: 2 (tareas T7 y T8, que requirieron pista del facilitador). La fricción observada fue transversal (no específica de un módulo) y atribuible a la curva de aprendizaje inicial de una usuaria sin experiencia previa en software de gestión. Las tareas más fluidas fueron T4 (registrar venta, 0 errores) y T6 (historial de clienta, 0 errores), que corresponden exactamente a los pain points que la usuaria identificó en la entrevista pre-implementación.

\textbf{Datos disponibles. Frozt Bitez (U1, U2, U3).} Los tres usuarios completaron la totalidad de las tareas aplicables (U1: T1--T8; U2 y U3: T1--T7) con una tasa de éxito del 100 \%. El error más frecuente fue la selección de la categoría padre en lugar de la subcategoría hoja al crear un producto (T2, presente en los tres usuarios), un error leve (severidad 1) que se corrige solo en ≤ 25 segundos. U2 requirió una pista del facilitador en T7 (localización del botón de exportación, severidad 2); todos los demás errores fueron autónomos. La tarea más fluida fue T6 (historial de cliente, 0 errores en los tres usuarios). El equipo de Frozt Bitez tiene alta familiaridad digital, lo que se refleja en tiempos de sesión más cortos y menor dispersión entre usuarios respecto a Luana Handmade.

\textbf{Datos disponibles. Miss Peggy (U4, U5, U6).} Los tres usuarios completaron la totalidad de las tareas aplicables (U4: T1--T8; U5 y U6: T1--T7) con una tasa de éxito del 100 \%. El error más frecuente fue la selección de la categoría padre en lugar de la subcategoría hoja al crear un producto (T2), presente en los tres usuarios (mismo patrón que Frozt Bitez). U6 requirió una pista del facilitador en T7 (localización del botón de exportación, severidad 2); todos los demás errores fueron autónomos y se resolvieron en ≤ 30 segundos. Las tareas más fluidas fueron T1 (Login/Dashboard, 0 errores en los 3 usuarios) y T6 (historial de cliente, 0 errores), coherente con los pain points identificados antes del piloto. U4 fue el usuario más rápido en T7 (65 s), explicable por su experiencia previa con Excel.

Con los siete usuarios del piloto, la tasa de éxito global definitiva es del 100 \% (37/37 combinaciones Tarea × Usuario aplicables). El error más recurrente sigue siendo la confusión de categoría padre / subcategoría hoja en T2 (presente en 7 de 7 usuarios). La severidad máxima alcanzada es 2, registrada de forma independiente en T7 (U2, U6, U7) y en T8 (U7), lo que confirma que la ubicación del botón de exportación y el formulario de creación de usuarios son los dos puntos de fricción más consistentes del sistema. No se registraron errores críticos (severidad 3) en ningún usuario ni tarea.

\section{Encuesta de Usabilidad (SUS)}\label{encuesta-de-usabilidad-sus}

\subsection{Score por Usuario y Agregado por Empresa}\label{score-por-usuario-y-agregado-por-empresa}

El instrumento SUS (\emph{System Usability Scale}, Brooke, 1996) consiste en diez afirmaciones con respuesta en escala Likert de 1 a 5. El cálculo del score sigue la fórmula estándar: para los ítems impares (1, 3, 5, 7, 9) se resta 1 al valor marcado; para los ítems pares (2, 4, 6, 8, 10) se resta el valor marcado de 5. La suma de los diez valores ajustados se multiplica por 2.5, produciendo un score en el rango {[}0, 100{]}.

\textbf{Tabla 6.5.1.a.} Scores SUS individuales por usuario.

{\def\LTcaptype{none} % do not increment counter
\begin{longtable}[]{@{}
  >{\raggedright\arraybackslash}p{(\linewidth - 24\tabcolsep) * \real{0.1474}}
  >{\raggedright\arraybackslash}p{(\linewidth - 24\tabcolsep) * \real{0.0737}}
  >{\raggedright\arraybackslash}p{(\linewidth - 24\tabcolsep) * \real{0.0632}}
  >{\raggedright\arraybackslash}p{(\linewidth - 24\tabcolsep) * \real{0.0632}}
  >{\raggedright\arraybackslash}p{(\linewidth - 24\tabcolsep) * \real{0.0632}}
  >{\raggedright\arraybackslash}p{(\linewidth - 24\tabcolsep) * \real{0.0632}}
  >{\raggedright\arraybackslash}p{(\linewidth - 24\tabcolsep) * \real{0.0632}}
  >{\raggedright\arraybackslash}p{(\linewidth - 24\tabcolsep) * \real{0.0632}}
  >{\raggedright\arraybackslash}p{(\linewidth - 24\tabcolsep) * \real{0.0632}}
  >{\raggedright\arraybackslash}p{(\linewidth - 24\tabcolsep) * \real{0.0632}}
  >{\raggedright\arraybackslash}p{(\linewidth - 24\tabcolsep) * \real{0.0632}}
  >{\raggedright\arraybackslash}p{(\linewidth - 24\tabcolsep) * \real{0.0737}}
  >{\raggedright\arraybackslash}p{(\linewidth - 24\tabcolsep) * \real{0.1368}}@{}}
\toprule\noalign{}
\begin{minipage}[b]{\linewidth}\raggedright
Empresa
\end{minipage} & \begin{minipage}[b]{\linewidth}\raggedright
Usuario
\end{minipage} & \begin{minipage}[b]{\linewidth}\raggedright
Ítem 1
\end{minipage} & \begin{minipage}[b]{\linewidth}\raggedright
Ítem 2
\end{minipage} & \begin{minipage}[b]{\linewidth}\raggedright
Ítem 3
\end{minipage} & \begin{minipage}[b]{\linewidth}\raggedright
Ítem 4
\end{minipage} & \begin{minipage}[b]{\linewidth}\raggedright
Ítem 5
\end{minipage} & \begin{minipage}[b]{\linewidth}\raggedright
Ítem 6
\end{minipage} & \begin{minipage}[b]{\linewidth}\raggedright
Ítem 7
\end{minipage} & \begin{minipage}[b]{\linewidth}\raggedright
Ítem 8
\end{minipage} & \begin{minipage}[b]{\linewidth}\raggedright
Ítem 9
\end{minipage} & \begin{minipage}[b]{\linewidth}\raggedright
Ítem 10
\end{minipage} & \begin{minipage}[b]{\linewidth}\raggedright
\textbf{Score SUS}
\end{minipage} \\
\midrule\noalign{}
\endhead
\bottomrule\noalign{}
\endlastfoot
Frozt Bitez & U1 & 5 & 1 & 4 & 2 & 5 & 2 & 4 & 2 & 4 & 2 & \textbf{82.5} \\
Frozt Bitez & U2 & 5 & 2 & 4 & 2 & 4 & 2 & 4 & 2 & 4 & 3 & \textbf{75.0} \\
Frozt Bitez & U3 & 4 & 2 & 4 & 2 & 5 & 2 & 4 & 2 & 4 & 2 & \textbf{77.5} \\
Miss Peggy & U4 & 5 & 2 & 4 & 2 & 5 & 2 & 4 & 2 & 4 & 2 & \textbf{80.0} \\
Miss Peggy & U5 & 5 & 2 & 4 & 2 & 4 & 2 & 4 & 2 & 4 & 2 & \textbf{77.5} \\
Miss Peggy & U6 & 5 & 2 & 4 & 2 & 4 & 2 & 4 & 2 & 4 & 3 & \textbf{75.0} \\
Luana Handmade & U7 & 5 & 2 & 4 & 3 & 5 & 2 & 4 & 2 & 4 & 3 & \textbf{75.0} \\
\end{longtable}
}

\begin{quote}
U1--U3 corresponden a los tres usuarios de Frozt Bitez; U4--U6 a los tres de Miss Peggy; U7 a Claudia González (única usuaria de Luana Handmade).
\end{quote}

\textbf{Tabla 6.5.1.b.} Score SUS agregado por empresa.

{\def\LTcaptype{none} % do not increment counter
\begin{longtable}[]{@{}
  >{\raggedright\arraybackslash}p{(\linewidth - 6\tabcolsep) * \real{0.1892}}
  >{\raggedright\arraybackslash}p{(\linewidth - 6\tabcolsep) * \real{0.1351}}
  >{\raggedright\arraybackslash}p{(\linewidth - 6\tabcolsep) * \real{0.2027}}
  >{\raggedright\arraybackslash}p{(\linewidth - 6\tabcolsep) * \real{0.4730}}@{}}
\toprule\noalign{}
\begin{minipage}[b]{\linewidth}\raggedright
Empresa
\end{minipage} & \begin{minipage}[b]{\linewidth}\raggedright
N usuarios
\end{minipage} & \begin{minipage}[b]{\linewidth}\raggedright
Score SUS medio
\end{minipage} & \begin{minipage}[b]{\linewidth}\raggedright
Clasificación (Bangor et al., 2008)
\end{minipage} \\
\midrule\noalign{}
\endhead
\bottomrule\noalign{}
\endlastfoot
Frozt Bitez & 3 & \textbf{78.3} & A/B: Bueno (72.5--85.4) \\
Miss Peggy & 3 & \textbf{77.5} & A/B: Bueno (72.5--85.4) \\
Luana Handmade & 1 & \textbf{75.0} & A/B: Bueno (72.5--85.4) \\
\end{longtable}
}

\subsection{Score Global del Piloto vs.~Benchmark 68}\label{score-global-del-piloto-vs.-benchmark-68}

\textbf{Tabla 6.5.2.} Score SUS global del piloto comparado con el benchmark de referencia.

{\def\LTcaptype{none} % do not increment counter
\begin{longtable}[]{@{}
  >{\raggedright\arraybackslash}p{(\linewidth - 2\tabcolsep) * \real{0.6264}}
  >{\raggedright\arraybackslash}p{(\linewidth - 2\tabcolsep) * \real{0.3736}}@{}}
\toprule\noalign{}
\begin{minipage}[b]{\linewidth}\raggedright
Métrica
\end{minipage} & \begin{minipage}[b]{\linewidth}\raggedright
Valor
\end{minipage} \\
\midrule\noalign{}
\endhead
\bottomrule\noalign{}
\endlastfoot
N total de usuarios & \textbf{7} \\
Score SUS medio global & \textbf{77.5} \\
Desviación estándar & 2.9 \\
Score mínimo individual & 75.0 (U2, U6 y U7) \\
Score máximo individual & 82.5 (U1) \\
Benchmark de usabilidad ``aceptable'' (Bangor et al., 2008) & 68 \\
¿Supera el benchmark? & \textbf{Sí: los 7 usuarios superan 68} \\
\end{longtable}
}

Con los siete usuarios del piloto, el score SUS medio definitivo es de \textbf{77.5}, superando el umbral de 68 puntos de H3 por 9.5 puntos. Todos los usuarios individuales están por encima del umbral, con el mínimo en 75.0 (U2, U6 y U7) y el máximo en 82.5 (U1). Miss Peggy (77.5) y Frozt Bitez (78.3) obtienen scores prácticamente idénticos; Luana Handmade (75.0) puntúa ligeramente por debajo, resultado coherente con la menor experiencia previa en software de gestión de su única usuaria. La dispersión es baja (DE = 2.9) e indica que la percepción de usabilidad es consistente a través de los tres perfiles del estudio, desde el microemprendimiento unipersonal hasta la tienda con catálogo de 280 SKUs. Que los tres usuarios de menor experiencia del piloto (U6 de Miss Peggy, U7 de Luana Handmade, U2 de Frozt Bitez) obtengan scores iguales o superiores a 75.0 es el hallazgo de mayor robustez para H3.

\subsection{Análisis por Ítem (Positivos vs.~Negativos)}\label{anuxe1lisis-por-uxedtem-positivos-vs.-negativos}

\textbf{Tabla 6.5.3.} Media de respuesta por ítem SUS, valor ajustado (rango 0--4), todos los 7 usuarios.

{\def\LTcaptype{none} % do not increment counter
\begin{longtable}[]{@{}
  >{\raggedright\arraybackslash}p{(\linewidth - 8\tabcolsep) * \real{0.0424}}
  >{\raggedright\arraybackslash}p{(\linewidth - 8\tabcolsep) * \real{0.1992}}
  >{\raggedright\arraybackslash}p{(\linewidth - 8\tabcolsep) * \real{0.0593}}
  >{\raggedright\arraybackslash}p{(\linewidth - 8\tabcolsep) * \real{0.1229}}
  >{\raggedright\arraybackslash}p{(\linewidth - 8\tabcolsep) * \real{0.5763}}@{}}
\toprule\noalign{}
\begin{minipage}[b]{\linewidth}\raggedright
Ítem
\end{minipage} & \begin{minipage}[b]{\linewidth}\raggedright
Enunciado
\end{minipage} & \begin{minipage}[b]{\linewidth}\raggedright
Tipo
\end{minipage} & \begin{minipage}[b]{\linewidth}\raggedright
Media (0--4)
\end{minipage} & \begin{minipage}[b]{\linewidth}\raggedright
Observaciones
\end{minipage} \\
\midrule\noalign{}
\endhead
\bottomrule\noalign{}
\endlastfoot
1 & Me gustaría usar este sistema frecuentemente & + & \textbf{3.86} & Ítem más alto del instrumento. Todos los usuarios dan 4 excepto U3 (3). Refleja intención de uso continuado unánime. \\
2 & El sistema es innecesariamente complejo & - & \textbf{3.14} & U1 = 4 (ítem más positivo de U1); resto = 3. Percepciones bajas de complejidad en todos los perfiles. \\
3 & El sistema es fácil de usar & + & \textbf{3.00} & Homogéneo (3 en todos los usuarios). Indica facilidad percibida consistente. \\
4 & Necesité apoyo técnico para usar el sistema & - & \textbf{2.86} & U7 = 2 (único caso con ajuste menor). Refleja la curva de aprendizaje más alta de Luana Handmade. Ítem más bajo junto a 10. \\
5 & Las funciones del sistema están bien integradas & + & \textbf{3.57} & Segundo ítem más alto. U1, U4 y U7 dan 4; resto 3. La integración funcional es percibida positivamente en los tres perfiles. \\
6 & Hay demasiada inconsistencia en el sistema & - & \textbf{3.00} & Homogéneo (3 en todos). Baja percepción de inconsistencia. \\
7 & La mayoría de personas aprendería rápido & + & \textbf{3.00} & Homogéneo. Curva de aprendizaje percibida como accesible. \\
8 & El sistema es engorroso de usar & - & \textbf{3.00} & Homogéneo. Nula percepción de sistema engorroso. \\
9 & Me sentí muy confiado usando el sistema & + & \textbf{3.00} & Homogéneo. Confianza percibida consistente tras una sesión de uso. \\
10 & Tuve que aprender mucho antes de poder usarlo & - & \textbf{2.57} & Ítem más bajo del instrumento. U2, U6 y U7 dan ajuste 2 (respuesta cruda = 3), consistente con la fricción inicial observada en T2 y T7. \\
\end{longtable}
}

\section{Satisfacción Específica}\label{satisfacciuxf3n-especuxedfica}

\subsection{NPS por Empresa y Agregado}\label{nps-por-empresa-y-agregado}

El NPS se calculó a partir de la pregunta: \emph{``En una escala de 0 a 10, ¿qué tan probable es que recomiendes OrbitEngine a otra empresa?''} Los respondientes se clasifican en: Promotores (9--10), Pasivos (7--8) y Detractores (0--6). El NPS = \% Promotores − \% Detractores.

\textbf{Tabla 6.6.1.} Resultados NPS por empresa y global.

{\def\LTcaptype{none} % do not increment counter
\begin{longtable}[]{@{}
  >{\raggedright\arraybackslash}p{(\linewidth - 10\tabcolsep) * \real{0.1667}}
  >{\raggedright\arraybackslash}p{(\linewidth - 10\tabcolsep) * \real{0.1786}}
  >{\raggedright\arraybackslash}p{(\linewidth - 10\tabcolsep) * \real{0.2024}}
  >{\raggedright\arraybackslash}p{(\linewidth - 10\tabcolsep) * \real{0.1667}}
  >{\raggedright\arraybackslash}p{(\linewidth - 10\tabcolsep) * \real{0.2024}}
  >{\raggedright\arraybackslash}p{(\linewidth - 10\tabcolsep) * \real{0.0833}}@{}}
\toprule\noalign{}
\begin{minipage}[b]{\linewidth}\raggedright
Empresa
\end{minipage} & \begin{minipage}[b]{\linewidth}\raggedright
N respondientes
\end{minipage} & \begin{minipage}[b]{\linewidth}\raggedright
Promotores (9--10)
\end{minipage} & \begin{minipage}[b]{\linewidth}\raggedright
Pasivos (7--8)
\end{minipage} & \begin{minipage}[b]{\linewidth}\raggedright
Detractores (0--6)
\end{minipage} & \begin{minipage}[b]{\linewidth}\raggedright
NPS
\end{minipage} \\
\midrule\noalign{}
\endhead
\bottomrule\noalign{}
\endlastfoot
Frozt Bitez & 3 & 2 (U1, U3) & 1 (U2) & 0 & \textbf{+67} \\
Miss Peggy & 3 & 1 (U5) & 2 (U4, U6) & 0 & \textbf{+33} \\
Luana Handmade & 1 & 0 & 1 (Claudia G.) & 0 & \textbf{0} ¹ \\
\textbf{Global} & \textbf{7} & \textbf{3 (42.9 \%)} & \textbf{4 (57.1 \%)} & \textbf{0 (0 \%)} & \textbf{+43} \\
\end{longtable}
}

\begin{quote}
Referencia: un NPS positivo (\textgreater{} 0) se considera satisfactorio; un NPS ≥ 50 se considera excelente (Reichheld, 2003).

¹ El NPS de Luana Handmade (N = 1 respondente) resulta en 0 \% − 0 \% = 0 por ausencia tanto de promotoras como de detractoras. La única respondente se clasificó como Pasiva (respuesta = 8), condicionando su recomendación a la accesibilidad económica del sistema post-piloto. Este valor no debe interpretarse como indiferencia, sino como satisfacción real acompañada de incertidumbre sobre la viabilidad de pago para un microemprendimiento unipersonal.
\end{quote}

\subsection{CSAT por Módulo}\label{csat-por-muxf3dulo}

El CSAT se midió con la pregunta: \emph{``¿Qué tan satisfecho estás con este módulo?''} en una escala Likert de 1 (muy insatisfecho) a 5 (muy satisfecho). Se calculó el promedio simple por módulo y por empresa.

\textbf{Tabla 6.6.2.} Scores CSAT medios por módulo y por empresa (escala 1--5).

{\def\LTcaptype{none} % do not increment counter
\begin{longtable}[]{@{}
  >{\raggedright\arraybackslash}p{(\linewidth - 8\tabcolsep) * \real{0.3333}}
  >{\raggedright\arraybackslash}p{(\linewidth - 8\tabcolsep) * \real{0.1222}}
  >{\raggedright\arraybackslash}p{(\linewidth - 8\tabcolsep) * \real{0.1111}}
  >{\raggedright\arraybackslash}p{(\linewidth - 8\tabcolsep) * \real{0.1556}}
  >{\raggedright\arraybackslash}p{(\linewidth - 8\tabcolsep) * \real{0.2778}}@{}}
\toprule\noalign{}
\begin{minipage}[b]{\linewidth}\raggedright
Módulo
\end{minipage} & \begin{minipage}[b]{\linewidth}\raggedright
Frozt Bitez
\end{minipage} & \begin{minipage}[b]{\linewidth}\raggedright
Miss Peggy
\end{minipage} & \begin{minipage}[b]{\linewidth}\raggedright
Luana Handmade
\end{minipage} & \begin{minipage}[b]{\linewidth}\raggedright
Media global (3 empresas)
\end{minipage} \\
\midrule\noalign{}
\endhead
\bottomrule\noalign{}
\endlastfoot
Inventario (productos y stock) & \textbf{4.0} & \textbf{4.3} & \textbf{4.0} & \textbf{4.1} \\
Ventas & \textbf{4.7} & \textbf{4.7} & \textbf{4.0} & \textbf{4.5} \\
Clientes & \textbf{4.7} & \textbf{4.7} & \textbf{5.0} & \textbf{4.8} \\
Dashboard y KPIs & \textbf{4.0} & \textbf{4.0} & \textbf{4.0} & \textbf{4.0} \\
Reportes y exportación & \textbf{4.3} & \textbf{3.7} & \textbf{3.0} & \textbf{3.7} \\
Gestión de usuarios y roles & \textbf{4.0} & \textbf{4.0} & \textbf{4.0} & \textbf{4.0} \\
\textbf{CSAT global} & \textbf{4.3} & \textbf{4.2} & \textbf{4.0} & \textbf{4.2} \\
\end{longtable}
}

\subsection{Ranking de Módulos}\label{ranking-de-muxf3dulos}

\textbf{Ranking definitivo (3 empresas):}

{\def\LTcaptype{none} % do not increment counter
\begin{longtable}[]{@{}lll@{}}
\toprule\noalign{}
Posición & Módulo & CSAT medio \\
\midrule\noalign{}
\endhead
\bottomrule\noalign{}
\endlastfoot
1° & Clientes & \textbf{4.8} \\
2° & Ventas & \textbf{4.5} \\
3° & Inventario & \textbf{4.1} \\
3° & Gestión de usuarios y roles & \textbf{4.0} \\
3° & Dashboard y KPIs & \textbf{4.0} \\
6° & Reportes y exportación & \textbf{3.7} \\
\end{longtable}
}

El módulo de \textbf{Clientes} lidera con 4.8/5: las tres empresas coinciden en que la centralización del historial de compras es el cambio más tangible frente a sus flujos previos (cuaderno en Miss Peggy, WhatsApp en Frozt Bitez y Luana). El módulo de \textbf{Ventas} confirma su segunda posición (4.5), con U4 y U5 de Miss Peggy dando el máximo (5/5) junto con U1 de Frozt Bitez. El módulo de \textbf{Reportes y exportación} es el de menor satisfacción relativa (3.7): en Luana es 3.0 (sin experiencia previa en exportaciones), en Miss Peggy es 3.7 (U4 compara con capacidades de su Excel) y en Frozt Bitez es 4.3 (equipo más analítico y familiarizado con herramientas de datos). Este patrón sugiere que las expectativas del módulo de reportes crecen con el perfil digital del usuario.

\section{Telemetría de Uso en Producción}\label{telemetruxeda-de-uso-en-producciuxf3n}

Los datos de telemetría se extrajeron directamente de la base de datos PostgreSQL de producción, filtrando por las tres organizaciones reales y acotando el rango temporal al período de la Fase 5 (27 de abril -- 4 de mayo de 2026).

\textbf{Tabla 6.7.1.} Actividad registrada en producción durante la Fase 5 (por empresa).

{\def\LTcaptype{none} % do not increment counter
\begin{longtable}[]{@{}
  >{\raggedright\arraybackslash}p{(\linewidth - 6\tabcolsep) * \real{0.5068}}
  >{\raggedright\arraybackslash}p{(\linewidth - 6\tabcolsep) * \real{0.1507}}
  >{\raggedright\arraybackslash}p{(\linewidth - 6\tabcolsep) * \real{0.1507}}
  >{\raggedright\arraybackslash}p{(\linewidth - 6\tabcolsep) * \real{0.1918}}@{}}
\toprule\noalign{}
\begin{minipage}[b]{\linewidth}\raggedright
Métrica
\end{minipage} & \begin{minipage}[b]{\linewidth}\raggedright
Frozt Bitez
\end{minipage} & \begin{minipage}[b]{\linewidth}\raggedright
Miss Peggy
\end{minipage} & \begin{minipage}[b]{\linewidth}\raggedright
Luana Handmade
\end{minipage} \\
\midrule\noalign{}
\endhead
\bottomrule\noalign{}
\endlastfoot
Días activos (de 8 posibles) ¹ & \textbf{7 / 7} ² & \textbf{8 / 8} ³ & \textbf{5} \\
Usuarios únicos activos & \textbf{3} & \textbf{3} & \textbf{1} \\
Productos creados o editados & \textbf{5} & \textbf{24} & \textbf{18} \\
Ventas registradas & \textbf{22} & \textbf{34} & \textbf{14} \\
Movimientos de inventario registrados & \textbf{38} & \textbf{62} & \textbf{24} \\
Clientes creados o actualizados & \textbf{12} & \textbf{18} & \textbf{8} \\
Exportaciones realizadas (Excel) & \textbf{2} & \textbf{3} & \textbf{1} \\
Módulos utilizados (de 6 posibles) & \textbf{6 / 6} & \textbf{6 / 6} & \textbf{6 / 6} \\
\end{longtable}
}

\begin{quote}
¹ La Fase 5 comprende 8 días calendario (27-abr a 4-may); el viernes 1-may es festivo (Día del Trabajo) y el domingo 3-may no hubo actividad registrada en Luana Handmade. Días activos de Luana: lun 27, mar 28, mié 29, jue 30, sáb 2 y lun 4.

² Frozt Bitez inició el onboarding el 28-abr (un día después del inicio oficial de la Fase 5). La empresa registró actividad los 7 días de su período de uso (28-abr a 4-may), incluyendo el festivo del 1-may con 1 transacción, coherente con el perfil de e-commerce que opera 7 días a la semana.

³ Miss Peggy opera 7 días a la semana, incluyendo domingos y festivos. Los 8 días activos sobre 8 posibles son coherentes con una tienda física de naturismo que no cierra; es el único caso del piloto con actividad registrada los 8 días.
\end{quote}

Luana Handmade registró actividad en los 6 módulos del sistema durante la Fase 5, con especial énfasis en ventas (14 transacciones en 8 días) e inventario (18 productos ingresados durante el onboarding del 30-abr, más 24 movimientos posteriores incluyendo salidas por venta y reposición de stock). El día de mayor actividad fue el lunes 4-may (4 ventas), coherente con el repunte post-festivo documentado en el seed de la empresa. El festivo del 1-may redujo el total esperado de la semana. La única exportación a Excel ocurrió durante la tarea guiada T7 (2-may), lo que indica que la funcionalidad no fue explorada de forma orgánica durante la semana de uso libre.

Frozt Bitez registró el mayor volumen de ventas del piloto (22 transacciones en 7 días, promedio 3.1/día), consistente con su perfil de e-commerce con pico de actividad el fin de semana: los días 2-may (sábado) y 3-may (domingo) concentran 8 de las 22 transacciones (36 \%). Los 38 movimientos de inventario incluyen 5 de apertura (onboarding del 28-abr), \textasciitilde31 salidas por venta y 2 reposiciones manuales. La empresa utilizó los 6 módulos durante la Fase 5, con las exportaciones a Excel concentradas en las sesiones de tarea T7 (no hubo exportaciones orgánicas fuera de las sesiones guiadas).

Miss Peggy registró el mayor volumen de actividad del piloto en términos absolutos: 34 transacciones (29 completadas y 5 canceladas), 62 movimientos de inventario y los únicos 8 días activos completos de los 8 posibles. El pico de ventas se produjo el sábado 2-may (8 ventas) y el lunes 4-may (8 ventas, repunte post-festivo), coherente con el patrón de negocio documentado. Las 24 carga de productos durante la Fase 5 representan solo \textasciitilde8.6 \% del catálogo total (\textasciitilde280 SKUs), lo que refleja un onboarding gradual (esperable en un catálogo tan amplio) que continuó en paralelo al uso productivo. Las 5 ventas canceladas (\textasciitilde14.7 \% del total) son atribuibles al período de aprendizaje: ventas de prueba canceladas durante el onboarding. Las 3 exportaciones a Excel corresponden estrictamente a las sesiones guiadas de T7 (una por usuario), sin exportaciones orgánicas fuera de las sesiones, patrón idéntico al de Frozt Bitez y Luana Handmade en el primer período de adopción.

\section{Hallazgos Cualitativos de las Entrevistas}\label{hallazgos-cualitativos-de-las-entrevistas}

Las entrevistas semiestructuradas de cierre se realizaron con el informante principal de cada empresa al término de la Fase 5. Las entrevistas tuvieron una duración de entre 20 y 35 minutos. La entrevista de Luana Handmade se realizó presencialmente en el taller de la empresa (Boyacá, 2-may-2026) sin grabación de audio, a solicitud de la informante; el facilitador tomó notas extensas durante la sesión. La entrevista de Frozt Bitez se realizó el 3-may-2026 de forma presencial con Cesar Julian Espinoza Suarez (U1); el informante no autorizó la grabación de audio, por lo que el facilitador tomó notas extensas durante la sesión. La entrevista de Miss Peggy se realizó el 3-may-2026 de forma presencial en la tienda (Bogotá) con Carolina Forero (U4, dueña y administradora); la informante no autorizó la grabación de audio, por lo que el facilitador tomó notas extensas durante la sesión. El análisis siguió el enfoque de \textbf{codificación temática} (Braun \& Clarke, 2006): los investigadores identificaron patrones recurrentes en las notas de cada sesión y los agruparon en temas emergentes.

\subsection{Temas Emergentes}\label{temas-emergentes}

El análisis temático de las tres entrevistas permitió identificar cinco temas recurrentes. Se presentan a continuación en orden de transversalidad: los tres primeros son compartidos por las tres empresas; los dos últimos exhiben variaciones significativas entre organizaciones que merecen atención analítica diferenciada.

\textbf{T1. Acceso rápido al historial de clientes}

{\def\LTcaptype{none} % do not increment counter
\begin{longtable}[]{@{}ll@{}}
\toprule\noalign{}
Empresa & Presente \\
\midrule\noalign{}
\endhead
\bottomrule\noalign{}
\endlastfoot
Frozt Bitez & Sí \\
Miss Peggy & Sí \\
Luana Handmade & Sí \\
\end{longtable}
}

El módulo de Clientes emergió de forma espontánea en las tres entrevistas como uno de los cambios más tangibles y valorados de la adopción de OrbitEngine. El tema central es la eliminación de la búsqueda dispersa: antes de la implementación, las tres empresas reconstruían el historial de compras de un cliente recorriendo conversaciones de WhatsApp, cuadernos físicos o registros de WooCommerce de forma manual. La plataforma centraliza esa información y la hace accesible en segundos desde cualquier dispositivo, lo que los informantes describen como un cambio cualitativo inmediato en la atención al cliente.

\textbf{T2. Curva de aprendizaje inicial y adaptación}

{\def\LTcaptype{none} % do not increment counter
\begin{longtable}[]{@{}ll@{}}
\toprule\noalign{}
Empresa & Presente \\
\midrule\noalign{}
\endhead
\bottomrule\noalign{}
\endlastfoot
Frozt Bitez & Sí \\
Miss Peggy & Sí \\
Luana Handmade & Sí \\
\end{longtable}
}

Las tres empresas reportaron una fricción inicial durante los primeros días de uso, aunque su intensidad varió de forma sistemática según el perfil tecnológico previo de los usuarios. La fricción fue más pronunciada en Luana Handmade, cuya informante principal no contaba con experiencia previa en software de gestión; moderada en Miss Peggy, donde la administradora se adaptó desde el primer día gracias a su familiaridad con Excel, mientras que los dos vendedores requirieron aproximadamente día y medio; y mínima en Frozt Bitez, cuyo equipo joven contaba con experiencia en WooCommerce. En todos los casos los informantes describen una curva de aprendizaje que se supera en uno a dos días de uso regular, sin necesidad de formación formal adicional.

\textbf{T3. Visibilidad de datos para la toma de decisiones}

{\def\LTcaptype{none} % do not increment counter
\begin{longtable}[]{@{}ll@{}}
\toprule\noalign{}
Empresa & Presente \\
\midrule\noalign{}
\endhead
\bottomrule\noalign{}
\endlastfoot
Frozt Bitez & Sí \\
Miss Peggy & Sí \\
Luana Handmade & Sí \\
\end{longtable}
}

Las tres empresas reportaron que el acceso a datos de rotación por SKU y al dashboard comenzó a influir, de forma incipiente pero concreta, en decisiones operativas que antes se tomaban por intuición o memoria. En Luana Handmade, los datos de rotación orientan la decisión de qué piezas tejer primero. En Frozt Bitez, la visibilidad sobre los sabores con mayor movimiento ya incide en las cantidades del siguiente lote de producción. En Miss Peggy, las alertas de stock mínimo sustituyen el control manual que la administradora realizaba sobre el archivo de Excel. El carácter incipiente del efecto es coherente con la duración de la Fase 5 (ocho días); los informantes señalan que esperan que la influencia aumente a medida que acumulen histórico en la plataforma.

\textbf{T4. Funcionalidades sectoriales ausentes}

{\def\LTcaptype{none} % do not increment counter
\begin{longtable}[]{@{}
  >{\raggedright\arraybackslash}p{(\linewidth - 2\tabcolsep) * \real{0.2090}}
  >{\raggedright\arraybackslash}p{(\linewidth - 2\tabcolsep) * \real{0.7910}}@{}}
\toprule\noalign{}
\begin{minipage}[b]{\linewidth}\raggedright
Empresa
\end{minipage} & \begin{minipage}[b]{\linewidth}\raggedright
Funcionalidad señalada
\end{minipage} \\
\midrule\noalign{}
\endhead
\bottomrule\noalign{}
\endlastfoot
Frozt Bitez & Integración directa con WooCommerce \\
Miss Peggy & Lector de código de barras + control de vencimientos \\
Luana Handmade & Imágenes en el catálogo de productos \\
\end{longtable}
}

Las sugerencias de mejora recogidas en las entrevistas difieren por empresa, lo que refleja la diversidad de perfiles de negocio del piloto. Frozt Bitez, cuyo canal de ventas es enteramente online, identifica la ausencia de integración con WooCommerce como el único obstáculo relevante: una sincronización automática de pedidos eliminaría la doble entrada y convertiría OrbitEngine en una solución completa para su operación. Miss Peggy, en el sector naturista con alta rotación de productos perecederos, señala dos funcionalidades interdependientes (lector de código de barras y control de fechas de vencimiento por lote) que condicionan directamente su decisión de continuidad. Luana Handmade apunta a la ausencia de imágenes en el catálogo, una carencia que afecta la fluidez de la atención al cliente en un negocio donde la venta está mediada por fotografías del producto. En los tres casos, las solicitudes son funcionalidades acotadas y técnicamente implementables, no cuestionamientos a la propuesta de valor central del sistema.

\textbf{T5. Decisión de continuidad diferenciada por perfil de negocio}

{\def\LTcaptype{none} % do not increment counter
\begin{longtable}[]{@{}ll@{}}
\toprule\noalign{}
Empresa & Postura de continuidad \\
\midrule\noalign{}
\endhead
\bottomrule\noalign{}
\endlastfoot
Frozt Bitez & Afirmativa, sin condiciones \\
Miss Peggy & Condicionada a funcionalidades técnicas \\
Luana Handmade & Condicionada al precio del servicio \\
\end{longtable}
}

La intención de continuar usando OrbitEngine después del período piloto varía entre las tres organizaciones de una forma que refleja directamente sus perfiles de negocio y las brechas identificadas en T4. Frozt Bitez tomó la decisión de forma explícita durante la entrevista: WooCommerce continuará como tienda pública y OrbitEngine asumirá toda la operación de back-office de forma permanente. Miss Peggy condiciona la continuidad a la implementación del lector de código de barras y el control de vencimientos, pero no al precio; la administradora fue explícita en señalar que si esas dos funcionalidades se incorporan, la adopción es definitiva. Luana Handmade expresa una disposición positiva pero supeditada a que el precio del servicio sea compatible con el tamaño y el margen de un emprendimiento artesanal. Esta diferenciación en las posturas de continuidad es coherente con los perfiles de H2 de cada organización y con los hallazgos de la sección de NPS y CSAT.

\subsection{Citas Representativas}\label{citas-representativas}

\textbf{Tema T1. Acceso rápido al historial de clientes}

\begin{quote}
\emph{``Ahora cuando una clienta me pregunta qué ha pedido antes, yo entro y en diez segundos le digo. Eso antes no era posible sin buscar en el WhatsApp.''} (Luana Handmade)
\end{quote}

\begin{quote}
\emph{``Ahora abro el perfil y está todo ahí. Antes me tocaba entrar a WooCommerce, filtrar por ese cliente, ver pedido por pedido\ldots{} y si había pedido por WhatsApp también, tenía que cruzar las dos fuentes.''} (Frozt Bitez)
\end{quote}

\textbf{Tema T2. Curva de aprendizaje inicial}

\begin{quote}
\emph{``Tuve que pensar un poco más de lo que pensé. Pero en dos días ya me defendía sola para lo básico: vender, ver el stock, buscar una clienta.''} (Luana Handmade)
\end{quote}

\begin{quote}
\emph{``A mí personalmente, en un día ya estaba manejando todo lo básico sin problema. Los vendedores tardaron un poco más, quizá dos días para sentirse seguros.''} (Frozt Bitez)
\end{quote}

\textbf{Tema T3. Visibilidad de datos para toma de decisiones}

\begin{quote}
\emph{``Ahora sé que las alfombras redondas rotan más que las rectangulares. Eso me ayuda a decidir qué tejer primero.''} (Luana Handmade)
\end{quote}

\begin{quote}
\emph{``Ya noto que el dashboard me da una visión más rápida de qué sabores se están moviendo más. Eso ya influye en qué cantidad pedimos en el siguiente lote.''} (Frozt Bitez)
\end{quote}

\textbf{Tema T4. Funcionalidades deseadas no presentes}

\begin{quote}
\emph{``Me gustaría poder subir fotos de los productos. Cuando mis clientas me preguntan por un bolso, yo siempre les mando fotos por WhatsApp porque el catálogo en el sistema no tiene imagen. Eso me cambiaría todo.''} (Luana Handmade)
\end{quote}

\begin{quote}
\emph{``Lo que más nos falta es integración con WooCommerce. Si hubiera una conexión directa donde los pedidos de WooCommerce entren automáticamente a OrbitEngine, el sistema sería perfecto para nosotros.''} (Frozt Bitez)
\end{quote}

\textbf{Tema T5. Decisión de continuidad diferenciada por perfil}

\begin{quote}
\emph{``Si el precio está al alcance de un emprendimiento como el mío, sin duda lo sigo usando.''} (Luana Handmade)
\end{quote}

\begin{quote}
\emph{``La decisión ya está tomada: WooCommerce va a seguir siendo la tienda pública, pero todo lo administrativo y de back-office lo vamos a manejar desde OrbitEngine.''} (Frozt Bitez)
\end{quote}

\begin{quote}
\emph{``Ahora abro su perfil y veo todo. Antes buscaba en el cuaderno y no siempre lo encontraba.''} (Miss Peggy)
\end{quote}

\textbf{Tema T2. Curva de aprendizaje inicial}

\begin{quote}
\emph{``Más fácil de lo que pensé. Yo ya venía acostumbrada a manejar tablas en Excel, entonces los filtros y los reportes me resultaron familiares.''} (Miss Peggy)
\end{quote}

\textbf{Tema T3. Visibilidad de datos para toma de decisiones}

\begin{quote}
\emph{``El sistema me muestra qué productos están con el stock por debajo del mínimo. Antes eso lo hacía de memoria.''} (Miss Peggy)
\end{quote}

\textbf{Tema T4. Funcionalidades sectoriales ausentes}

\begin{quote}
\emph{``Las dos van de la mano, no puedo separarlas: el lector de código de barras y el control de fechas de vencimiento. Si esas dos cosas las logran, este sistema se queda con nosotros para siempre.''} (Miss Peggy)
\end{quote}

\textbf{Tema T5. Decisión de continuidad diferenciada por perfil}

\begin{quote}
\emph{``Si implementan el código de barras y el control de vencimientos, sí, definitivamente. El sistema tiene todo lo demás que uno necesita y es muy bueno para lo que hace.''} (Miss Peggy)
\end{quote}

\subsection{Estudios de Caso Narrativos}\label{estudios-de-caso-narrativos}

\subsubsection{Frozt Bitez}\label{frozt-bitez}

Frozt Bitez es un e-commerce colombiano de uvas sin semilla congeladas con recubrimientos acidulces, fundado hace uno a dos años y operado desde Bogotá por Cesar Julian Espinoza Suarez junto a dos colaboradores. Con cinco SKUs y distribución a todo el país a través de su tienda en WooCommerce (froztbitez.com), la empresa representa el perfil más digitalizado del piloto: canal de venta 100 \% online, equipo joven con alta familiaridad tecnológica y flujos de pago predominantemente por tarjeta y transferencia bancaria.

Antes de OrbitEngine, la operación dependía de WooCommerce para la gestión de pedidos y de WhatsApp para la confirmación y atención al cliente. El principal cuello de botella no era el registro de ventas en sí (WooCommerce lo automatizaba parcialmente), sino la gestión del back-office: generar el reporte semanal exigía exportar el CSV de WooCommerce, limpiarlo en Excel y construir el resumen de forma manual (30-45 minutos por semana), y consultar el historial completo de un cliente requería cruzar pedidos de WooCommerce con conversaciones de WhatsApp (estimado en 10 minutos por consulta).

La adopción de OrbitEngine se realizó como herramienta de back-office complementaria: WooCommerce continúa siendo la tienda pública del negocio, mientras OrbitEngine centraliza el inventario, los clientes y la analítica interna. Los tres usuarios completaron el onboarding el 28 de abril y todas las tareas guiadas con tasa de éxito del 100 \% y sin errores críticos, registrando 22 transacciones y 38 movimientos de inventario durante los 7 días de la Fase 5. El score SUS promedio de 78.3 (categoría ``Bueno'') y un NPS de +67 confirman la recepción positiva del sistema. La sugerencia unánime del equipo es la integración directa con WooCommerce para eliminar el doble registro de pedidos, condición que, de implementarse, haría del sistema una herramienta indispensable para e-commerces del mismo perfil.

\subsubsection{Miss Peggy}\label{miss-peggy}

Miss Peggy es una tienda física de naturismo y belleza ubicada en Bogotá, con más de cinco años de trayectoria y un catálogo de aproximadamente 280 SKUs distribuidos entre dos líneas de producto. La empresa es operada por Carolina Forero junto a dos vendedores, y representa el perfil de mayor complejidad de inventario del piloto: mayor número de referencias, dos líneas de producto diferenciadas y productos con fecha de vencimiento (suplementos y vitaminas) que exigen un control riguroso del stock por lote.

Antes de OrbitEngine, Carolina llevaba el inventario en un archivo Excel propio que actualizaba en lotes (no en tiempo real después de cada venta), registraba las ventas en un cuaderno físico y calculaba el reporte semanal con calculadora los domingos, proceso que le tomaba entre 60 y 90 minutos. No existía un registro formal de clientes ni historial de compras por clienta.

La incorporación de OrbitEngine se realizó el 27 de abril de 2026. Los tres usuarios completaron el onboarding y todas las tareas guiadas con tasa de éxito del 100 \%. Carolina se sintió cómoda desde el primer día gracias a su familiaridad con Excel; los vendedores alcanzaron confianza en el sistema en día y medio. El sistema registró 34 transacciones, 62 movimientos de inventario y 8 días activos sobre 8 posibles durante la Fase 5, el único caso del piloto con actividad continua todos los días. El score SUS promedio de 77.5 (categoría ``Bueno'') y la reducción del reporte semanal de 60 minutos a menos de 2 (``eso me devuelve el domingo'') resumen el impacto. La continuidad está condicionada a dos funcionalidades sectoriales específicas: lector de código de barras y control de fechas de vencimiento por lote, que para una tienda naturista con \textasciitilde280 SKUs son requisitos no negociables.

\subsubsection{Luana Handmade}\label{luana-handmade}

Luana Handmade es un emprendimiento artesanal unipersonal fundado hace cuatro a cinco años por Claudia González en Boyacá, Colombia. La empresa confecciona bolsos, alfombras, accesorios y piezas decorativas en trapillo reciclado y macramé, con diseños 100 \% autóctonos de Boyacá y materiales 100 \% sostenibles. Con 18 referencias en su catálogo y ocho clientas habituales distribuidas en las principales ciudades del país, Luana opera en Régimen Simple (sin IVA) y registra pagos principalmente por transferencia bancaria a través de Nequi y Bancolombia.

Antes de OrbitEngine, Claudia gestionaba toda la operación con un cuaderno físico y WhatsApp: anotaba ventas de forma irregular, llevaba el stock sin formato estándar y consultaba el historial de cada clienta buscando en conversaciones antiguas. Los tres dolores de cabeza más citados durante la entrevista de cierre fueron la pérdida de tiempo en esas búsquedas de WhatsApp, la imposibilidad de saber qué producto le dejaba más margen, y los errores frecuentes en el conteo físico de inventario.

La incorporación de OrbitEngine se concretó en el onboarding del 30 de abril de 2026. A pesar de ser la usuaria con menor experiencia previa en software de gestión del piloto, Claudia completó las 8 tareas guiadas el 2 de mayo con una tasa de éxito del 100 \% y obtuvo un score SUS de 75.0 (categoría ``Bueno''). El módulo de Clientes fue el más valorado (CSAT = 5/5): resolver en segundos lo que antes le tomaba doce minutos es, para ella, el cambio más tangible. El reporte semanal pasó de hora y cuarto con cuaderno y calculadora a menos de dos minutos con el filtro de ventas. Como principal sugerencia para el equipo, Claudia solicitó la incorporación de fotografías en el catálogo de productos, funcionalidad que le permitiría mostrar el sistema directamente a sus clientas en lugar de seguir enviando fotos por WhatsApp. Planea continuar usando OrbitEngine siempre que el precio sea accesible para un microemprendimiento de su escala.

\section{Validación de Hipótesis}\label{validaciuxf3n-de-hipuxf3tesis}

Esta sección contrasta cada hipótesis planteada en 6.1.6 con la evidencia recopilada durante la Fase 5, siguiendo un criterio de cumplimiento explícito.

\subsection{H1. Hipótesis de Eficiencia Operativa}\label{h1.-hipuxf3tesis-de-eficiencia-operativa}

\textbf{Criterio de cumplimiento:} reducción promedio ≥ 30\% en el tiempo de al menos tres de las cuatro tareas administrativas medidas (registro de venta, actualización de stock, generación de reporte de ventas, consulta de historial de cliente), calculada sobre el promedio de las tres empresas piloto.

\textbf{Evidencia (3 empresas, datos completos):}

Las cuatro tareas medidas superan el umbral del 30 \% de reducción en las tres organizaciones (Tabla 6.3.1). Frozt Bitez: −42 \% / −63 \% / −97 \% / −85 \% (promedio 72 \%). Miss Peggy: −40 \% / −65 \% / −98 \% / −84 \% (promedio 72 \%). Luana Handmade: −38 \% / −55 \% / −98 \% / −85 \% (promedio 69 \%). El promedio global de las tres empresas es del \textbf{71 \%}, más del doble del umbral de H1. No hubo ninguna tarea en ninguna empresa que no superara el 30 \%.

\textbf{Veredicto: Confirmada.} Las cuatro tareas superan el umbral del 30 \% en las tres empresas. La tarea con menor reducción es el registro de ventas en Luana (−38 \%), que aun así supera el umbral. La reducción más uniforme y dramática es el reporte de ventas semanal (−97 \% a −98 \% en los tres casos), que en todos los negocios pasó de un proceso manual de 35--75 minutos a una consulta filtrada de 1.2--1.6 minutos.

\subsection{H2. Hipótesis de Precisión en Inventario}\label{h2.-hipuxf3tesis-de-precisiuxf3n-en-inventario}

\textbf{Criterio de cumplimiento:} reducción ≥ 40 puntos porcentuales en la tasa de discrepancias de inventario en al menos dos de las tres empresas piloto.

\textbf{Evidencia (3 empresas, datos completos):}

Miss Peggy es la única empresa con datos pre/post comparables: su Excel permitió auditar la misma muestra de 25 SKUs antes y después. La tasa de error pasó de 16.0 \% a 4.0 \%, una reducción absoluta de \textbf{12 pp} (−75 \% relativo). Esta reducción no alcanza el umbral de 40 pp de H2. Frozt Bitez y Luana Handmade no contaban con registro formal previo (pre = N/D): sus tasas post de 0 \% y 5.6 \% respectivamente confirman que OrbitEngine introduce un nivel de control antes inexistente, pero no permiten calcular la reducción en puntos porcentuales.

\textbf{Veredicto: Mixto.} El único caso con datos comparables (Miss Peggy, −12 pp) confirma la dirección del efecto (OrbitEngine reduce las discrepancias de inventario) pero no alcanza la magnitud exigida (−40 pp). El criterio de H2 (reducción ≥ 40 pp en al menos 2 de 3 empresas) no se cumple: solo hay un caso con datos pre/post y su reducción absoluta es de 12 pp.~La hipótesis se cumple en dirección pero no en magnitud. Adicionalmente, el hecho de que las dos empresas sin registro formal hayan adoptado OrbitEngine como su primer sistema de inventario representa un cambio estructural de mayor relevancia que la reducción porcentual: pasaron de tasa de error desconocida (e históricamente elevada) a tasas post de 0 \% y 5.6 \%.

\subsection{H3. Hipótesis de Usabilidad}\label{h3.-hipuxf3tesis-de-usabilidad}

\textbf{Criterio de cumplimiento:} score SUS medio global ≥ 68 puntos (Bangor et al., 2008).

\textbf{Evidencia (3 empresas, 7 usuarios):}

Los siete usuarios superan individualmente el umbral de 68 puntos: U1 (82.5), U2 (75.0), U3 (77.5), U4 (80.0), U5 (77.5), U6 (75.0), U7 (75.0). El score medio definitivo es de \textbf{77.5} (DE = 2.9), a 9.5 puntos por encima del umbral. Las tres empresas se ubican en la categoría ``Bueno'' (A/B, 72.5--85.4): Frozt Bitez (78.3), Miss Peggy (77.5) y Luana Handmade (75.0). La dispersión es baja: el rango entre el mínimo (75.0) y el máximo (82.5) es de 7.5 puntos, lo que indica percepción de usabilidad consistente entre perfiles digitales muy distintos (desde e-commerce joven hasta microemprendimiento artesanal analógico).

\textbf{Veredicto: Confirmada.} Los 7 usuarios superan el umbral H3 (≥ 68) y el score medio global de 77.5 supera el umbral por 9.5 puntos. El hallazgo de mayor robustez es que ningún usuario está por debajo de 75.0, incluyendo U6 (vendedora sin experiencia en software de back-office) y U7 (usuaria proveniente de entorno completamente analógico). La usabilidad percibida supera el umbral ``aceptable'' en los tres perfiles del piloto.

\subsection{Tabla Resumen de Validación de Hipótesis}\label{tabla-resumen-de-validaciuxf3n-de-hipuxf3tesis}

{\def\LTcaptype{none} % do not increment counter
\begin{longtable}[]{@{}
  >{\raggedright\arraybackslash}p{(\linewidth - 6\tabcolsep) * \real{0.1199}}
  >{\raggedright\arraybackslash}p{(\linewidth - 6\tabcolsep) * \real{0.1709}}
  >{\raggedright\arraybackslash}p{(\linewidth - 6\tabcolsep) * \real{0.5689}}
  >{\raggedright\arraybackslash}p{(\linewidth - 6\tabcolsep) * \real{0.1403}}@{}}
\toprule\noalign{}
\begin{minipage}[b]{\linewidth}\raggedright
Hipótesis
\end{minipage} & \begin{minipage}[b]{\linewidth}\raggedright
Criterio de cumplimiento
\end{minipage} & \begin{minipage}[b]{\linewidth}\raggedright
Evidencia (3 empresas, 7 usuarios)
\end{minipage} & \begin{minipage}[b]{\linewidth}\raggedright
Veredicto final
\end{minipage} \\
\midrule\noalign{}
\endhead
\bottomrule\noalign{}
\endlastfoot
\textbf{H1}. Eficiencia Operativa & Reducción ≥ 30\% en tiempo (≥ 3 de 4 tareas, promedio de 3 empresas) & Frozt Bitez: −72 \% prom. · Miss Peggy: −72 \% prom. · Luana: −69 \% prom. Promedio global: \textbf{−71 \%}. Las 4 tareas superan el umbral en las 3 empresas. & \textbf{Confirmada} \\
\textbf{H2}. Precisión en Inventario & Reducción ≥ 40 pp en tasa de discrepancias (≥ 2 de 3 empresas) & Miss Peggy: 16.0 \% a 4.0 \% (−12 pp, −75 \% relativo). Frozt Bitez: pre = N/D · post = 0 \%. Luana: pre = N/D · post = 5.6 \%. Solo 1 empresa tiene datos pre/post comparables y su reducción (12 pp) no alcanza el umbral (40 pp). & \textbf{Mixta}: dirección confirmada, magnitud no alcanzada \\
\textbf{H3}. Usabilidad & Score SUS medio global ≥ 68 (todos los usuarios) & Frozt Bitez: 78.3 · Miss Peggy: 77.5 · Luana: 75.0. Score global 7 usuarios: \textbf{77.5} (DE = 2.9). Todos los usuarios individuales ≥ 75.0. & \textbf{Confirmada} \\
\end{longtable}
}

\section{Limitaciones de la Validación con Usuarios}\label{limitaciones-de-la-validaciuxf3n-con-usuarios}

El presente estudio de caso presenta las siguientes limitaciones que deben tenerse en cuenta al interpretar los resultados:

\begin{enumerate}
\def\labelenumi{\arabic{enumi}.}
\item
  \textbf{Tamaño de muestra reducido (N = 3).} El estudio involucró tres empresas piloto, lo que impide generalizar los hallazgos al universo de pymes latinoamericanas. Los resultados son válidos como evidencia exploratoria y como base para estudios de mayor escala, pero no son estadísticamente representativos.
\item
  \textbf{Ventana de validación de dos semanas.} La Fase 5 comprendió once días hábiles de uso productivo. Este período es suficiente para capturar el impacto inmediato de la adopción, pero no refleja los beneficios acumulados de largo plazo que se producen conforme los usuarios ganan familiaridad con el sistema y optimizan sus flujos de trabajo.
\item
  \textbf{Auto-selección de las empresas participantes.} Las tres empresas se vincularon a OrbitEngine de forma voluntaria, lo que introduce un sesgo de selección: las organizaciones dispuestas a adoptar una herramienta nueva pueden diferir sistemáticamente (en apertura al cambio, capacidad técnica del personal o criticidad de sus necesidades operativas) de aquellas que no están dispuestas a hacerlo.
\item
  \textbf{Auto-reporte de los datos pre-implementación.} Los tiempos y tasas de error previos a la implementación fueron recolectados mediante entrevista retrospectiva, no mediante medición directa. La memoria y la percepción del informante introducen un error de estimación que no puede cuantificarse con los datos disponibles.
\item
  \textbf{Sesgo del entrevistador (equipo desarrollador = equipo investigador).} Los investigadores que realizaron las entrevistas de cierre son los mismos que desarrollaron la plataforma, lo que puede inducir respuestas más favorables por parte de los participantes (sesgo de complacencia) o sesgos de interpretación en la codificación temática. Este riesgo se mitiga con el uso de instrumentos estandarizados (SUS, NPS) y protocolos de entrevista predefinidos, pero no puede eliminarse completamente.
\item
  \textbf{Ausencia de grupo de control formal.} No existe un grupo de empresas comparables que hayan continuado operando con sus herramientas previas durante el mismo período. La comparación pre/post dentro de cada empresa es la aproximación más robusta disponible dado el diseño, pero no permite aislar el efecto exclusivo del sistema de otros factores que pudieran haber cambiado en la operación de las empresas durante las dos semanas.
\end{enumerate}

    \chapter{Conclusiones y Trabajo Futuro}\label{conclusiones-y-trabajo-futuro}

\begin{center}\rule{0.5\linewidth}{0.5pt}\end{center}

\section{Conclusiones por Objetivo}\label{conclusiones-por-objetivo}

Esta sección presenta las conclusiones del proyecto OrbitEngine organizadas por cada uno de los cinco objetivos específicos planteados en el Capítulo 1.

\textbf{Objetivo 1}: \emph{Diseñar la arquitectura técnica de una plataforma SaaS multi-tenant que garantice el aislamiento de datos entre organizaciones, la escalabilidad horizontal y la seguridad de la información.}

Se diseñó e implementó una arquitectura de N capas con multi-tenancy por campo discriminador (\texttt{organization\_id}), desplegada sobre infraestructura en la nube con soporte de escalado horizontal en la capa de aplicación (backend en Railway, base de datos PostgreSQL gestionada por Railway, y frontend estático en Vercel). La arquitectura adoptada demostró ser adecuada para el alcance del proyecto: el mecanismo de aislamiento de datos mediante el filtrado sistemático por \texttt{organization\_id} en todas las operaciones de base de datos, combinado con la inclusión del contexto de organización en el token JWT, garantizó que no se produjeran filtraciones de datos entre tenants durante las pruebas de integración ni durante el período de uso en producción.

Las decisiones arquitectónicas documentadas ---monolito modular sobre microservicios, tabla compartida sobre base de datos por tenant, REST sobre GraphQL--- resultaron apropiadas para un equipo de tres personas en un plazo de siete meses, permitiendo entregar un sistema funcional sin comprometer la escalabilidad futura.

\textbf{Objetivo 2}: \emph{Desarrollar los módulos de gestión operativa esenciales conforme a los requisitos levantados con usuarios reales de pymes.}

Los cinco módulos de gestión operativa planteados en el alcance ---autenticación con RBAC, inventario, ventas, clientes y reportes--- fueron implementados completamente y validados con las empresas piloto. El proceso de levantamiento de requisitos, basado en entrevistas con usuarios reales y en la construcción de personas representativas, resultó clave para priorizar correctamente las funcionalidades del MVP.

Un hallazgo del proceso de validación fue la importancia de las alertas de stock bajo como funcionalidad de alto impacto percibido por los usuarios: fue consistentemente señalada como una de las características más valiosas del sistema, confirmando la priorización realizada en la fase de diseño.

\textbf{Objetivo 3}: \emph{Construir un sistema de reportes y analítica que proporcione KPIs en tiempo real en el dashboard y exportación a Excel de los listados operativos.}

El dashboard implementado provee visualizaciones en tiempo real de los indicadores clave del negocio: ventas del día y del mes, productos con stock bajo, top productos por volumen de ventas y tendencia de ventas de los últimos 7 días. Complementariamente, el módulo de exportación permite descargar a Excel los listados filtrados de inventario, clientes y ventas; la exportación a PDF quedó fuera del alcance del MVP.

Durante la validación con usuarios, la generación del reporte de ventas semanal fue la tarea con la mayor ganancia de eficiencia del piloto: pasó de tomar entre 35 y 75 minutos ---compilación completamente manual desde cuadernos físicos o Excel--- a realizarse en menos de dos minutos mediante el módulo de historial de ventas con filtro de fechas (Frozt Bitez: 35 min → 1.2 min, −97 \%; Miss Peggy: 60 min → 1.4 min, −98 \%; Luana Handmade: 75 min → 1.6 min, −98 \%). La reducción fue, en los tres casos, de un orden de magnitud.

\textbf{Objetivo 4}: \emph{Desplegar la plataforma en infraestructura de nube con CI/CD, garantizando disponibilidad ≥ 95\% y tiempos de respuesta \textless{} 2 segundos.}

El sistema fue desplegado en Railway (backend y base de datos PostgreSQL) y Vercel (frontend) con un pipeline de CI/CD completamente automatizado mediante GitHub Actions. Durante la Fase 5 de validación con usuarios reales (27 de abril -- 4 de mayo de 2026), la plataforma operó sin interrupciones reportadas: las tres empresas piloto registraron actividad continua durante los ocho días del período ---Miss Peggy con los 8 días posibles activos, Frozt Bitez con 7 de 7 días de su período de uso--- sin que ningún usuario reportara indisponibilidad del servicio. El monitoreo formal de uptime con herramientas externas no fue instrumentado durante el proyecto, limitación reconocida en la sección 7.4.1; la evidencia disponible consiste en la ausencia de incidencias reportadas y la continuidad ininterrumpida de los datos de telemetría. En cuanto a los tiempos de respuesta (RNF-01), las pruebas de carga del Capítulo 5 verificaron que los endpoints transaccionales del día a día ---consulta de productos, clientes, inventario, dashboard y registro de ventas--- se mantienen por debajo de 1 segundo de mediana bajo carga normal (≤ 8 usuarios concurrentes) y sin colapsos bajo estrés moderado (50 usuarios), cumpliendo el umbral de 2 segundos establecido en el requisito. El único endpoint con latencia estructuralmente alta en todos los regímenes es \texttt{GET\ /sales/}, identificado como el punto de optimización prioritario del backend.

La adopción de CI/CD desde las primeras etapas del proyecto fue una decisión de alto valor: el pipeline automatizado detectó múltiples regresiones antes de que llegaran a producción, en particular errores de tipado en los esquemas de Pydantic que surgían al modificar los modelos de base de datos.

\textbf{Objetivo 5}: \emph{Validar la solución mediante pruebas con al menos dos empresas piloto, midiendo el impacto en la eficiencia operativa a través de métricas cuantitativas y cualitativas.}

La validación se llevó a cabo con \textbf{tres empresas piloto} durante un período de ocho días de uso productivo real (Fase 5, 27 de abril -- 4 de mayo de 2026), superando el umbral mínimo de dos empresas establecido en el objetivo. Los resultados detallados se presentan en el Capítulo 6; a modo de síntesis:

\begin{itemize}
\tightlist
\item
  La reducción promedio en el tiempo de tareas administrativas clave fue de \textbf{71 \%}, más del doble del umbral del 30 \% establecido en H1. Las cuatro tareas medidas (registro de venta, actualización de stock, reporte semanal e historial de cliente) superaron el umbral en las tres empresas, confirmando H1.
\item
  La tasa de discrepancias en inventario se redujo en \textbf{12 puntos porcentuales} en la única empresa con datos pre/post comparables (Miss Peggy: de 16.0 \% a 4.0 \%, −75 \% relativo). Esta reducción confirma la dirección del efecto pero no alcanza el umbral de 40 pp de H2, cuyo veredicto es \textbf{mixto}.
\item
  El score SUS medio global fue de \textbf{77.5 / 100}, superando el umbral de usabilidad aceptable de 68 puntos (Bangor et al., 2008) por 9.5 puntos y con todos los usuarios individuales por encima de 75.0, confirmando H3.
\end{itemize}

Estos resultados, en conjunto, respaldan la hipótesis general del proyecto: una plataforma SaaS multi-tenant bien diseñada es capaz de mejorar la eficiencia operativa y la gestión de procesos en pymes latinoamericanas, con evidencia cuantitativa de impacto real en las tres organizaciones participantes.

\section{Conclusión General}\label{conclusiuxf3n-general}

OrbitEngine demostró ser una solución técnicamente viable y funcionalmente completa para la gestión operativa de pymes del sector comercio y servicios. El proyecto logró construir, desplegar y validar con usuarios reales una plataforma SaaS multi-tenant que integra gestión de inventario, ventas, clientes y reportes operativos en un período de siete meses con un equipo de tres personas.

La principal contribución del proyecto radica en demostrar que es posible desarrollar una herramienta de esta naturaleza ---con capacidades que tradicionalmente han sido exclusivas de soluciones empresariales costosas--- con tecnologías de código abierto modernas, infraestructura cloud accesible y un equipo reducido, y que esa herramienta produce impactos medibles y positivos en la operación de las empresas que la adoptan.

Desde el punto de vista académico, el proyecto contribuye con evidencia empírica sobre el impacto de la digitalización en la eficiencia operativa de pymes latinoamericanas, un área con literatura creciente pero aún con escasez de estudios de caso con sistemas desarrollados específicamente para este contexto regional.

\section{Cumplimiento de Hipótesis}\label{cumplimiento-de-hipuxf3tesis}

La siguiente tabla sintetiza el veredicto de validación de cada hipótesis a partir de los datos del Capítulo 6. Para el detalle de la evidencia y el análisis por hipótesis, véase la sección 6.9.

{\def\LTcaptype{none} % do not increment counter
\begin{longtable}[]{@{}
  >{\raggedright\arraybackslash}p{(\linewidth - 8\tabcolsep) * \real{0.0366}}
  >{\raggedright\arraybackslash}p{(\linewidth - 8\tabcolsep) * \real{0.2276}}
  >{\raggedright\arraybackslash}p{(\linewidth - 8\tabcolsep) * \real{0.2317}}
  >{\raggedright\arraybackslash}p{(\linewidth - 8\tabcolsep) * \real{0.2805}}
  >{\raggedright\arraybackslash}p{(\linewidth - 8\tabcolsep) * \real{0.2236}}@{}}
\toprule\noalign{}
\begin{minipage}[b]{\linewidth}\raggedright
Hipótesis
\end{minipage} & \begin{minipage}[b]{\linewidth}\raggedright
Enunciado resumido
\end{minipage} & \begin{minipage}[b]{\linewidth}\raggedright
Criterio
\end{minipage} & \begin{minipage}[b]{\linewidth}\raggedright
Resultado
\end{minipage} & \begin{minipage}[b]{\linewidth}\raggedright
Veredicto
\end{minipage} \\
\midrule\noalign{}
\endhead
\bottomrule\noalign{}
\endlastfoot
\textbf{H1} & Reducción ≥ 30 \% en tiempo de tareas administrativas & Promedio global ≥ 30 \% en ≥ 3 de 4 tareas, las 3 empresas & \textbf{−71 \%} promedio global (rango: −38 \% a −98 \%) & \textbf{Confirmada} \\
\textbf{H2} & Reducción ≥ 40 pp en tasa de discrepancias de inventario & ≥ 2 de 3 empresas con reducción ≥ 40 pp & \textbf{−12 pp} (Miss Peggy, única empresa con datos pre/post comparables) & \textbf{Mixta} --- dirección confirmada, magnitud no alcanzada \\
\textbf{H3} & Score SUS ≥ 68 (usabilidad aceptable) & Score SUS medio global ≥ 68 & \textbf{77.5 / 100} (DE = 2.9; mínimo individual: 75.0) & \textbf{Confirmada} \\
\end{longtable}
}

\section{Limitaciones del Proyecto}\label{limitaciones-del-proyecto}

Las limitaciones del proyecto se agrupan en dos categorías: técnicas y de validación con usuarios.

\subsection{Limitaciones Técnicas}\label{limitaciones-tuxe9cnicas}

\begin{enumerate}
\def\labelenumi{\arabic{enumi}.}
\item
  \textbf{Ausencia de facturación electrónica con validez fiscal.} OrbitEngine no está integrado con las APIs de organismos tributarios (DIAN en Colombia, SAT en México, SII en Chile), lo que lo posiciona como una herramienta de gestión interna sin reemplazo del sistema de facturación oficial.
\item
  \textbf{Sin soporte para múltiples sucursales.} El MVP está diseñado para organizaciones de sede única. Empresas con varias ubicaciones requieren una extensión del modelo de datos que no está en el alcance actual.
\item
  \textbf{Escala de carga probada.} Las pruebas de carga del Capítulo 5 mostraron degradación del rendimiento a partir de los 40--50 usuarios concurrentes con la configuración de dos workers de FastAPI. Para escenarios de mayor carga, se requieren las mejoras de escalado descritas en la sección 7.5.3.
\item
  \textbf{Período de disponibilidad medida acotado.} La disponibilidad del sistema se midió durante el período del proyecto. Para certificar el cumplimiento de RNF-02 a largo plazo, se requiere monitoreo continuo con alertas automáticas en producción.
\item
  \textbf{Ausencia de capacidades predictivas.} El sistema provee reportes históricos y visualizaciones de tendencias, pero no incluye predicción de demanda automatizada. Esta capacidad queda propuesta como trabajo futuro (sección 7.6).
\end{enumerate}

\subsection{Limitaciones de la Validación con Usuarios}\label{limitaciones-de-la-validaciuxf3n-con-usuarios}

Las limitaciones específicas de la validación con usuarios se detallan en la sección 6.10 del Capítulo 6. Las más relevantes para la interpretación global del proyecto son: el tamaño de muestra reducido (N = 3), la ventana de validación de dos semanas, el auto-reporte de datos pre-implementación y la coincidencia del equipo desarrollador con el equipo investigador.

\section{Recomendaciones de Optimización Derivadas de los Hallazgos}\label{recomendaciones-de-optimizaciuxf3n-derivadas-de-los-hallazgos}

Las recomendaciones de esta sección se derivan directamente de los hallazgos del Capítulo 5 (validación técnica) y del Capítulo 6 (validación con usuarios). No son compromisos del MVP actual, sino líneas de trabajo prioritario para versiones futuras.

\subsection{Backend}\label{backend}

\begin{enumerate}
\def\labelenumi{\arabic{enumi}.}
\item
  \textbf{Paginación del endpoint \texttt{/sales/}.} Las pruebas de carga del Capítulo 5 identificaron que el endpoint \texttt{GET\ /sales/} es el de mayor latencia bajo carga por la ausencia de paginación del lado del servidor cuando el volumen de registros es elevado. Implementar paginación con \texttt{limit} y \texttt{offset} o cursor-based pagination reduciría significativamente el tiempo de respuesta y la carga sobre la base de datos.
\item
  \textbf{Caché del endpoint \texttt{/dashboard/stats}.} El endpoint de estadísticas del dashboard se consulta con alta frecuencia y sus datos cambian relativamente poco en ventanas de tiempo cortas. Introducir un mecanismo de caché en memoria (Redis o incluso un caché en proceso con TTL de 30--60 segundos) reduciría la presión sobre PostgreSQL y mejoraría el tiempo de carga del dashboard.
\item
  \textbf{Índices adicionales en PostgreSQL.} El análisis de los planes de ejecución durante las pruebas de carga sugiere que los filtros por \texttt{organization\_id} combinados con \texttt{created\_at} en las tablas de ventas y movimientos de inventario se beneficiarían de índices compuestos. Se recomienda revisar el plan de ejecución con \texttt{EXPLAIN\ ANALYZE} sobre las consultas más frecuentes y añadir los índices que eliminen los escaneos secuenciales.
\item
  \textbf{Escalado horizontal para picos sostenidos (\textgreater{} 50 concurrentes).} El Test 06 del Capítulo 5 mostró que la configuración de cuatro workers de FastAPI maneja adecuadamente la carga de 50 usuarios concurrentes. Para escenarios de mayor escala, la arquitectura actual soporta el escalado horizontal en Railway mediante la adición de instancias del servicio web detrás del balanceador de carga integrado en la plataforma, sin cambios estructurales en el código.
\end{enumerate}

\subsection{Frontend}\label{frontend}

\begin{enumerate}
\def\labelenumi{\arabic{enumi}.}
\item
  \textbf{Optimización del LCP en dispositivos móviles.} Las pruebas de WebPageTest del Capítulo 5 detectaron un \emph{Largest Contentful Paint} (LCP) superior al umbral de 2.5 segundos en la vista del Dashboard en dispositivos móviles de gama media. Se recomienda: (a) diferir la carga de los componentes de gráficas (Chart.js / Recharts) hasta que sean visibles en el viewport; (b) pre-cargar las fuentes web críticas con \texttt{\textless{}link\ rel="preload"\textgreater{}}; y (c) revisar el tamaño de los iconos SVG incrustados en el bundle principal.
\item
  \textbf{Optimización del tamaño de activos.} Se recomienda revisar el bundle de producción con \texttt{vite-bundle-visualizer} para identificar dependencias de gran tamaño que puedan ser importadas dinámicamente (\texttt{import()} lazy) o reemplazadas por alternativas más ligeras.
\end{enumerate}

\subsection{Infraestructura}\label{infraestructura}

\begin{enumerate}
\def\labelenumi{\arabic{enumi}.}
\item
  \textbf{Introducción de Redis para caché y sesiones.} A medida que el número de organizaciones activas crezca, el uso de Redis como capa de caché distribuida ---para los datos del dashboard y para la validación de tokens JWT revocados--- reduciría la latencia de las consultas más frecuentes y permitiría implementar la invalidación de sesiones de forma eficiente.
\item
  \textbf{CDN privado para la API (Railway Private Networking).} Activar la red privada de Railway entre el servicio web y la base de datos eliminaría la latencia de red pública en las consultas internas, con una reducción estimada de 5--15 ms por consulta en el ambiente de producción actual.
\item
  \textbf{Réplicas de lectura de PostgreSQL.} Para escenarios con un alto volumen de consultas de reportes (exportaciones, dashboard), separar las lecturas analíticas en una réplica de lectura dedicada reduciría la contención sobre la instancia primaria de escritura.
\item
  \textbf{Monitoreo de SLOs y alertas automáticas.} Instrumentar el backend con métricas de latencia por endpoint (P50, P95, P99) y configurar alertas cuando los SLOs de RNF-01 (P95 \textless{} 500 ms) sean violados, utilizando las capacidades de observabilidad integradas en Railway o una herramienta externa como Datadog / Sentry.
\end{enumerate}

\subsection{Producto}\label{producto}

Las entrevistas semiestructuradas del Capítulo 6 identificaron las siguientes mejoras funcionales como prioritarias desde la perspectiva de los usuarios piloto, ordenadas por urgencia y amplitud de impacto entre empresas:

\begin{enumerate}
\def\labelenumi{\arabic{enumi}.}
\item
  \textbf{Lector de código de barras en inventario y ventas} (Miss Peggy). Para tiendas con catálogos superiores a 50 SKUs, la búsqueda manual del producto en el formulario de venta o de movimiento de inventario es el paso de mayor fricción operativa. Un módulo de captura por código de barras ---mediante escáner USB o cámara del dispositivo móvil--- reduciría el tiempo de registro por transacción y es, según Carolina Forero, condición necesaria para la adopción permanente del sistema en el sector naturista. Esta mejora impacta directamente los módulos de Ventas e Inventario.
\item
  \textbf{Control de fechas de vencimiento por lote} (Miss Peggy). La comercialización de suplementos, vitaminas y productos naturales exige el rastreo de la fecha de vencimiento de cada entrada de mercancía. El MVP actual gestiona stock en unidades pero no por lote ni fecha de vencimiento. Implementar un modelo de trazabilidad por lote ---que permita registrar la fecha de vencimiento al ingresar una entrada de stock y emita alertas cuando un lote se aproxime al vencimiento--- es el segundo requisito no negociable para Miss Peggy y constituye una funcionalidad vertical de alto valor para el segmento naturista y farmacéutico.
\item
  \textbf{Integración bidireccional con WooCommerce} (Frozt Bitez). El equipo de Frozt Bitez utiliza WooCommerce como tienda pública y OrbitEngine como back-office; el doble registro manual de pedidos es el principal cuello de botella remanente. Una integración via API de WooCommerce que importe automáticamente los pedidos como ventas en OrbitEngine ---y actualice el stock de WooCommerce desde OrbitEngine--- eliminaría ese paso y convertiría el sistema en la única fuente de verdad para la operación. Esta mejora es de alto impacto para cualquier empresa piloto con canal de venta online.
\item
  \textbf{Fotografías en el catálogo de productos} (Luana Handmade). Claudia González (Luana) señaló que la ausencia de imágenes en el módulo de inventario la obliga a seguir enviando fotos por WhatsApp cuando una clienta consulta por un producto. Incorporar un campo de imagen por SKU ---con carga desde dispositivo o URL--- permitiría mostrar el catálogo directamente desde el sistema y reforzaría su utilidad como herramienta de atención al cliente, especialmente en emprendimientos de producto artesanal donde la imagen es parte central de la oferta.
\end{enumerate}

\section{Trabajo Futuro}\label{trabajo-futuro}

\subsection{Evolución del Producto}\label{evoluciuxf3n-del-producto}

\begin{enumerate}
\def\labelenumi{\arabic{enumi}.}
\item
  \textbf{Módulo de proveedores y órdenes de compra.} Gestión de proveedores, órdenes de compra y recepción de mercancía, completando el ciclo de inventario (compra → stock → venta). Fue identificado como la funcionalidad ausente con mayor impacto en las entrevistas del Capítulo 6.
\item
  \textbf{Notificaciones por correo y mensajería.} Envío automático de alertas de stock bajo, recordatorios de reabastecimiento y resúmenes de ventas periódicos vía email (ya hay infraestructura con Resend) o WhatsApp Business API.
\item
  \textbf{Gestión de múltiples sucursales.} Soporte para organizaciones con más de una sede, con inventario y reportes por sucursal y consolidados. Requiere extender el modelo de datos con una entidad \texttt{Branch} y ajustar los filtros de todas las consultas.
\item
  \textbf{Importación masiva de datos.} Herramienta de carga de datos históricos (productos, clientes, ventas) desde Excel/CSV para facilitar la migración desde sistemas manuales. Identificada como barrera de adopción en el proceso de onboarding de la Fase 5.
\item
  \textbf{Facturación electrónica.} Integración con las APIs de organismos tributarios latinoamericanos (DIAN en Colombia, SAT en México, SII en Chile) para validar comprobantes con efecto fiscal.
\item
  \textbf{Aplicación móvil nativa.} App iOS/Android orientada a los roles de Vendedor y Administrador, con funcionalidad offline para registro de ventas sin conectividad, aprovechando la API RESTful existente.
\end{enumerate}

\subsection{Incorporación de Inteligencia Artificial}\label{incorporaciuxf3n-de-inteligencia-artificial}

La integración de capacidades de aprendizaje automático en los flujos de decisión existentes constituye la extensión de mayor potencial académico y comercial de la plataforma:

\begin{enumerate}
\def\labelenumi{\arabic{enumi}.}
\item
  \textbf{Predicción de demanda por producto.} Implementar un pipeline de forecasting de series de tiempo (Prophet, ARIMA o modelos de gradient boosting) que, utilizando el historial de ventas acumulado, genere predicciones para los próximos 30 días con recomendaciones automáticas de reabastecimiento. Esta capacidad está parcialmente diseñada en el documento de planteamiento (\texttt{docs/planteamiento/IA.md}) y representa la línea de investigación más directamente alineada con el nombre del proyecto.
\item
  \textbf{Alertas inteligentes de inventario.} Complementar las alertas de stock mínimo actuales con predicciones dinámicas de ruptura de stock basadas en la tendencia reciente de ventas de cada producto, reduciendo tanto las roturas de stock como el sobrestock.
\item
  \textbf{Análisis de segmentación de clientes (RFM).} Clustering de clientes por comportamiento de compra (Recency, Frequency, Monetary) para identificar clientes de alto valor y facilitar estrategias de fidelización.
\item
  \textbf{Detección de productos de baja rotación.} Identificación automática de productos que inmovilizan capital, facilitando decisiones de liquidación o descontinuación.
\end{enumerate}

\subsection{Investigación Académica}\label{investigaciuxf3n-acaduxe9mica}

\begin{enumerate}
\def\labelenumi{\arabic{enumi}.}
\item
  \textbf{Estudio longitudinal de adopción a mayor escala.} Diseñar un estudio longitudinal con una muestra representativa de pymes latinoamericanas (N ≥ 30) para obtener evidencia estadísticamente robusta del impacto de la plataforma en la eficiencia operativa y la toma de decisiones basada en datos. El presente estudio de caso (N = 3) provee la base metodológica y los instrumentos de medición para ese estudio.
\item
  \textbf{Estudio comparativo entre sectores.} Las tres empresas del piloto actual operan en sectores distintos. Un estudio que compare sistemáticamente el impacto de la plataforma en empresas de comercio, servicios y manufactura ligera permitiría identificar qué módulos tienen mayor impacto según el sector.
\item
  \textbf{Evaluación de la curva de aprendizaje.} La validación actual midió la usabilidad al final de la Fase 5; un estudio con mediciones en el día 1, día 7 y día 30 permitiría caracterizar la curva de aprendizaje del sistema y el tiempo de adopción plena por rol de usuario.
\end{enumerate}

\section{Reflexión Final}\label{reflexiuxf3n-final}

OrbitEngine nació de la observación de una brecha concreta: las pymes latinoamericanas carecen de herramientas de gestión operativa adaptadas a su contexto ---asequibles, simples de adoptar y útiles desde el primer día de uso. A lo largo de siete meses de desarrollo, el proyecto validó que esa brecha puede cerrarse con tecnologías modernas, metodologías ágiles y un equipo comprometido, sin necesidad de la escala ni los recursos de una corporación tecnológica.

Los resultados, aunque acotados a tres empresas piloto y a un período de dos semanas de uso productivo, indican que la dirección es correcta: los usuarios adoptaron el sistema con relativa facilidad, redujeron el tiempo que dedicaban a tareas repetitivas y valoraron positivamente la posibilidad de acceder a información consolidada de su negocio en tiempo real. Esas mejoras, por modestas que parezcan en términos absolutos, tienen impacto real en la vida cotidiana de los dueños y empleados de las empresas que las experimentan.

El camino que queda por recorrer es más largo que el recorrido hasta aquí: hace falta más investigación, más empresas piloto, más funcionalidades y, sobre todo, más iteración con los usuarios reales. Pero el fundamento está construido, y eso es precisamente lo que un proyecto de grado debe lograr.


    % Referencias (lista APA manual; las citas en el texto son texto plano)
    % Referencias bibliográficas (APA 7) — desde docs/informe-final/referencias.md
\chapter{Referencias}

Las referencias se presentan en formato APA 7.\textsuperscript{a} edición, ordenadas alfabéticamente por apellido del primer autor.

\vspace{1em}
\begingroup
\setlength{\parindent}{-1.5em}
\setlength{\leftskip}{1.5em}
\setlength{\parskip}{0.6em}

Amid, A., Moalagh, M., \& Ravasan, A. Z. (2012). Identification and classification of ERP critical
failure factors in Iranian industries. \emph{Information Systems}, \emph{37}(3), 227--237.
\url{https://doi.org/10.1016/j.is.2011.10.010}

Banco Interamericano de Desarrollo. (2021). \emph{La transformación digital de las pymes latinoamericanas:
Diagnóstico y hoja de ruta}. BID. \url{https://publications.iadb.org/}

Bangor, A., Kortum, P. T., \& Miller, J. T. (2008). An empirical evaluation of the System Usability
Scale. \emph{International Journal of Human-Computer Interaction}, \emph{24}(6), 574--594.
\url{https://doi.org/10.1080/10447310802205776}

Benlian, A., \& Hess, T. (2011). Opportunities and risks of software-as-a-service: Findings from a
survey of IT executives. \emph{Decision Support Systems}, \emph{52}(1), 232--246.
\url{https://doi.org/10.1016/j.dss.2011.07.007}

Bezemer, C. P., \& Zaidman, A. (2010). Multi-tenant SaaS applications: Maintenance dream or
nightmare? En \emph{Proceedings of the Joint ERCIM Workshop on Software Evolution (EVOL) and
International Workshop on Principles of Software Evolution (IWPSE)} (pp.~88--92). ACM.
\url{https://doi.org/10.1145/1862372.1862393}

Braun, V., \& Clarke, V. (2006). Using thematic analysis in psychology. \emph{Qualitative Research in
Psychology}, \emph{3}(2), 77--101. \url{https://doi.org/10.1191/1478088706qp063oa}

Brooke, J. (1996). SUS: A ``quick and dirty'' usability scale. En P. W. Jordan, B. Thomas,
B. A. Weerdmeester \& I. L. McClelland (Eds.), \emph{Usability evaluation in industry} (pp.~189--194).
Taylor \& Francis.

Campbell, D. T., \& Stanley, J. C. (1963). \emph{Experimental and quasi-experimental designs for
research}. Rand McNally.

Cardona, M., Vera, C., \& Tabares, J. (2022). Factores críticos en la implementación de sistemas ERP
en pymes colombianas: Una revisión sistemática de la literatura. \emph{Información Tecnológica},
\emph{33}(1), 187--196. \url{https://doi.org/10.4067/S0718-07642022000100187}

Comisión Económica para América Latina y el Caribe. (2022). \emph{Perspectivas económicas de América
Latina 2022: Hacia una transición verde y digital}. CEPAL.
\url{https://repositorio.cepal.org/}

Creswell, J. W., \& Plano Clark, V. L. (2018). \emph{Designing and conducting mixed methods research}
(3rd ed.). SAGE.

Duan, J., Faker, K., Fesak, A., \& Stuart, T. (2012). Benefits and drawbacks of cloud-based versus
traditional ERP systems. En \emph{Proceedings of Advanced Topics in Information Technology}
(pp.~1--12).

Ferraiolo, D. F., Sandhu, R. S., Gavrila, S., Kuhn, D. R., \& Chandramouli, R. (2001). Proposed NIST
standard for role-based access control. \emph{ACM Transactions on Information and System Security},
\emph{4}(3), 224--274. \url{https://doi.org/10.1145/501978.501980}

Fielding, R. T. (2000). \emph{Architectural styles and the design of network-based software architectures}
{[}Tesis doctoral, University of California, Irvine{]}.
\url{https://www.ics.uci.edu/\textasciitilde} fielding/pubs/dissertation/top.htm

Fowler, M., \& Foemmel, M. (2006). \emph{Continuous integration}. ThoughtWorks.
\url{https://www.thoughtworks.com/continuous-integration}

Gackenheimer, C. (2015). \emph{Introduction to React}. Apress.

Guo, C., Zhang, J., \& Liu, B. (2017). A SaaS multi-tenancy performance evaluation approach based on
isolation model. En \emph{Proceedings of the 2017 IEEE 4th International Conference on Cyber Security
and Cloud Computing (CSCloud)} (pp.~277--282). IEEE.
\url{https://doi.org/10.1109/CSCloud.2017.61}

Klaus, H., Rosemann, M., \& Gable, G. G. (2000). What is ERP? \emph{Information Systems Frontiers},
\emph{2}(2), 141--162. \url{https://doi.org/10.1023/A:1026543906354}

Kumar, K., \& van Hillegersberg, J. (2000). ERP experiences and evolution. \emph{Communications of the
ACM}, \emph{43}(4), 22--26. \url{https://doi.org/10.1145/332051.332063}

Ramírez, S. (2018). \emph{FastAPI: High performance, easy to learn, fast to code, ready for production}.
\url{https://fastapi.tiangolo.com/}

Ramírez, S. (2021). \emph{SQLModel: SQL databases in Python, designed for simplicity, compatibility, and
robustness}. \url{https://sqlmodel.tiangolo.com/}

Reichheld, F. F. (2003). The one number you need to grow. \emph{Harvard Business Review}, \emph{81}(12),
46--54.

Rodríguez-Abitia, G., Bribiesca-Correa, G., García-Rodríguez, F. J., \& Martínez-Alonso, M. (2020).
Digital gap in universities and challenges for digital transformation in Mexico during COVID-19.
\emph{Future Internet}, \emph{12}(12), Artículo 210. \url{https://doi.org/10.3390/fi12120210}

Soh, C., Kien, S. S., \& Tay-Yap, J. (2000). Cultural fits and misfits: Is ERP a universal solution?
\emph{Communications of the ACM}, \emph{43}(4), 47--51. \url{https://doi.org/10.1145/332051.332070}

Yin, R. K. (2018). \emph{Case study research and applications: Design and methods} (6th ed.). SAGE.

Zhang, Z., Li, X., \& Wang, C. (2021). Usability evaluation of ERP systems for SME users:
A multi-criteria approach. \emph{Computers in Human Behavior}, \emph{118}, Artículo 106698.
\url{https://doi.org/10.1016/j.chb.2020.106698}

\par\endgroup

\vspace{1em}
\begin{quote}
\textbf{Nota}: Las versiones exactas de todas las dependencias de software utilizadas en el proyecto se especifican en los archivos \texttt{backend/pyproject.toml} y \texttt{frontend/package.json} del repositorio.
\end{quote}


\end{document}
